\documentclass[11pt]{article}
\usepackage[margin=1in]{geometry}
\usepackage{amsmath,amssymb,amsthm,mathtools}
\usepackage{booktabs}
\usepackage{enumitem}
\usepackage{graphicx}
\usepackage{float}
\usepackage[colorlinks=true,linkcolor=blue,citecolor=blue]{hyperref}

% ---------------------------------------------------------------
% Theorem environments
% ---------------------------------------------------------------
\newtheorem{theorem}{Theorem}
\newtheorem{lemma}{Lemma}
\newtheorem{proposition}{Proposition}
\newtheorem{corollary}{Corollary}
\theoremstyle{definition}
\newtheorem{definition}{Definition}
\newtheorem{assumption}{Assumption}
\theoremstyle{remark}
\newtheorem*{remark}{Remark}

% ---------------------------------------------------------------
% Notation macros (one symbol, one meaning)
% ---------------------------------------------------------------
\newcommand{\R}{\mathbb{R}}
\newcommand{\Aset}{\mathcal{A}}          % data set
\newcommand{\Net}{\mathcal{N}_\theta}    % network
\newcommand{\Lip}{\operatorname{Lip}_{\Aset}}
\newcommand{\supA}[1]{\|#1\|_{\infty}}
\newcommand{\Rgh}{\mathcal{R}}           % roughness function (seminorm)
\newcommand{\Lipn}{\Rgh_{\mathrm{Lip}}}  % Lipschitz seminorm on the training inputs
\newcommand{\sR}{\sigma_{\Rgh}}
\newcommand{\CR}{C_{\Rgh}}
\newcommand{\id}{\mathrm{id}}
\newcommand{\ip}[2]{\langle #1,#2\rangle}
\newcommand{\cost}{\mathcal{B}}          % cost of a direction along a path
\newcommand{\Hess}{\mathcal{H}}          % loss Hessian
% NOTE: main.tex used \mathcal{H} (\Hard, \HR) for the local complexity of a
% stepper. In this document \mathcal{H} is the loss Hessian. The local
% complexity enters only in the applications (Part III) and needs a new symbol
% there.

\title{Roughness and the error of a trained network\\[4pt]
\large Lower bounds at one time during training, and the error along the training path}
\author{Draft}
\date{}

\begin{document}
\maketitle

\begin{abstract}
This paper gives a theory of how the roughness of a dataset limits the error of a network trained on it, first at one time during training and then along the whole training path.
Roughness is measured by a seminorm on the values at the training inputs, the roughness function.
The labels have an absolute roughness, in the units of the labels; for a seminorm fixed by the training inputs it is a property of the data alone, computed before training.
The first part links this roughness to the eigenvalues of the training kernel, the kernel of the network's gradient features.
At every time during training, the part of the error on the slow eigendirections of the kernel is at least the absolute roughness of the labels, minus the roughness of the current network output, minus a leak term, divided by a constant of the seminorm.
Consequently a rough dataset pushes the error onto the slow eigendirections, and one training step of a given length can reduce the error only by an amount that decreases with the roughness of the labels.
These statements hold for every update rule, and the network enters only through the eigenvalues of its kernel, the roughness of its current output, and its peak feature roughness, the most roughness that a unit parameter step can add to the output.
For a single training step, the supremum of the bound over all seminorms equals the smallest error that the step reaches to first order in its length, which is the smallest error actually reached up to a remainder of second order; it is attained by a unique linear functional, which is the limit of a Krylov series built from the error by repeated application of the kernel.
The second part integrates the first along training.
Training is a closed system in one state variable, the parameters, driven by a tower of derivative tensors of the network.
Along every gradient-flow path the loss obeys exact identities, the energy identity and the speed identity, in which the whole path enters through one scalar function of time, the curvature $\kappa$ of the loss along the gradient; $\kappa$ splits into a part fixed by where the error lies in the spectrum of the kernel and a part through which feature learning enters the loss curve.
The length of the parameter path is the coordinate in which the kernel moves.
With it, the roughness bound of the first part becomes a bound on the error after a stretch of training, in terms of the absolute roughness of the dataset, the initial output, the path length, and the derivatives of the network at the start.
Read the other way, a rough dataset needs a minimum path length, and under gradient flow a minimum training time, before the error can fall to a given value; this holds at every time during training.
The error is expanded exactly to all orders, by a Dyson series in the motion of the kernel and by a Picard scheme from the initial state that converges to the true path on every finite window, at a cost that grows with the window; the paper states why none of the expansions gives a bound uniform over all of training, and which assumptions on the path would close the system.
The central result is a general prediction framework: for any design trained by any update rule that is a smooth flow or map on the parameters and the optimizer state, including gradient descent with schedules, minibatches, momentum and Adam, the error at every time in a finite window is predicted from the initial state by an explicit convergent scheme, with a two-sided error bound computed along with the prediction and a check that certifies it.
Combined with the lower bounds, it gives a general principle for comparing designs at a given training time.
As a first application, a design that starts closer to its target provably keeps the smaller error throughout a training budget against a design that cannot close its own gap in time; for learned time steppers this proves that, for every training budget, the direct step is worse than the Euler step below an explicit time step.
Further applications are deferred to the last part.
\end{abstract}

\tableofcontents

% ===============================================================
\section{Introduction}\label{sec:intro}
% ===============================================================

A network is trained on a dataset of inputs and labels.
Some datasets are learned quickly and to a small error, others slowly or not at all.
This paper asks how the difficulty of a dataset can be measured before training, and proves that a harder dataset leaves a larger error, at every time during training and along the whole training path.

The difficulty of a dataset is the roughness of its labels.
The Lipschitz constant on the set from which the inputs are drawn is one measure of roughness.
It takes the worst case over all pairs of inputs, and it gives a weak lower bound, because its ratio to the root-mean-square norm grows with the number and density of the training inputs (Section~\ref{sec:examples}).
This paper measures roughness by a seminorm on the values at the training inputs.
Any seminorm on sampled functions can be used.
Each defines two constants, a norm constant and a peak feature roughness, and the bounds are stated in terms of these.
Section~\ref{sec:examples} gives three examples: the variation within clusters of nearby inputs, the manifold high-pass built from a graph Laplacian on the inputs, and the Lipschitz seminorm.
The sharpest seminorm is not fixed by the inputs.
It is the Krylov functional of Section~\ref{sec:step}, which is built from the current error and the training kernel.

The lower bound links the roughness of the labels to the eigenvalues of the training kernel, the kernel of the network's gradient features.
The kernel is not held fixed.
The paper has two theoretical parts, which differ in what they say about the motion of the kernel.

\paragraph{Part~\ref{part:onetime}: the error at one time.}
Every statement of Part~\ref{part:onetime} is made at one time during training, with the kernel at that time.
These results concern the rate of change of the error and what one training step can do.
The contributions are the following.
\begin{enumerate}[label=(\arabic*)]
\item The best roughness function (Section~\ref{sec:step}).
Every roughness function gives a lower bound on the error after one training step of a given length.
The supremum of these bounds over all roughness functions equals the smallest error that the step reaches in the first-order Taylor expansion about the current parameters (Theorem~\ref{thm:stepsup}), and the error actually reached is within a remainder of second order in the step length of it (Proposition~\ref{prop:taylorstep}).
It is attained by a unique linear functional, which has a kernel-ridge closed form and is the limit of a Krylov series built from the error by repeated application of the kernel (Theorem~\ref{thm:stepkrylov}).
\item A lower bound on the part of the error that lies on the slow eigendirections of the kernel, in terms of the roughness of the labels and the roughness of the current output (Theorem~\ref{thm:floor}).
Each eigendirection is weighted by how slowly gradient flow removes it, with no threshold between slow and fast directions.
\item The main statement of Part~\ref{part:onetime} (Theorem~\ref{thm:slow}): at every time during training and for every update rule, the error is at least the absolute roughness of the dataset, minus the roughness of the current output, minus a leak term, divided by a constant of the seminorm; and a training step of a given length can reduce the error only by an amount that decreases with the absolute roughness of the dataset (Theorem~\ref{thm:dataset}).
The network enters only through the eigenvalues of its kernel, the roughness of its current output, and its peak feature roughness, and the best choice of seminorm makes the roughness-of-the-error form of both bounds exact.
\end{enumerate}

\paragraph{Part~\ref{part:path}: the error along training.}
The error after a stretch of training requires integrating the rate along a parameter path whose kernel changes, and the path is not known in advance.
Part~\ref{part:path} sets up the framework in which that gap is closed as far as it can be closed without approximation.
Every remainder there is an exact integral or an infinite series with all terms, and every bound on a series is a bound on its whole tail.
\begin{enumerate}[label=(\arabic*),resume]
\item The closed system (Section~\ref{sec:state}): training has one state variable, the parameters, and a tower of derivative tensors, the Taylor series of the network at the start of the path, kept whole.
\item Exact identities along every gradient-flow path (Section~\ref{sec:identities}): the energy identity, the speed identity, which routes the whole path through one scalar function of time, the curvature $\kappa$ of the loss along the gradient, and the time derivatives of the error to all orders.
\item The path calculus (Section~\ref{sec:path}): the length of the parameter path is the coordinate in which the kernel moves.
With it, the floor of Part~\ref{part:onetime} becomes a bound on the error after a stretch of training (Theorems~\ref{thm:path} and~\ref{thm:time0}): the absolute roughness of the dataset, minus the roughness of the initial output, minus the path length times the peak feature roughness, minus a series in the path length set by the derivatives of the network at the start.
Read the other way (Corollary~\ref{cor:mintime}), this gives a minimum path length, and under gradient flow a minimum training time, before the error can reach a given value; it grows with the absolute roughness of the dataset and holds at every time, with no horizon.
A companion two-sided bound expands the error exactly about the kernel at the start, with the whole motion of the kernel kept in the remainder (Theorem~\ref{thm:duhamel}).
\item The error to all orders (Section~\ref{sec:allorders}): a Dyson series in the motion of the kernel, exact for every time, and a Picard scheme from the initial state, whose every iterate is computable from the network near the start, from the tower at the start within its radius of convergence, and whose limit is the true path on every finite window; and the reasons why none of them gives a bound uniform over all of training.
\item The general prediction framework (Section~\ref{sec:predict}), the central result: training of any design by any update rule that is a continuously differentiable flow or map on the parameters and the optimizer state is predicted from the initial state, at every time in a finite window, by a convergent scheme with a two-sided error bound and a check that certifies it (Theorem~\ref{thm:predict}); gradient flow, gradient descent with schedules, minibatches, momentum and Adam are instances.
With the lower bounds of the earlier items, it gives a general principle for comparing two designs at a given training time (Corollary~\ref{cor:compare}).
\item The loss curve read through the theory (Section~\ref{sec:loss}): the kernel term of $\kappa$, set by the kernel at each time, is bounded by the slow part of the error of Part~\ref{part:onetime}, and the rest is one scalar function of time through which feature learning enters the loss curve.
Section~\ref{sec:closing} states the assumptions on the path that would close the system, each with its consequence proved.
\end{enumerate}

Together, the two parts give one account of training.
The labels demand roughness $\Rgh(Y)$.
The network can supply roughness only at rate $\sR$ per unit of parameter motion, and only along eigendirections whose eigenvalue is at most $\sR^2$ divided by their squared roughness; so the rough part of the error sits on the slow eigenvalues, and one step can remove only a small part of it (Part~\ref{part:onetime}).
Along training, the parameter motion is the path length, the error falls at the squared speed of the parameters, and the speed decays at the curvature $\kappa$, whose kernel term is small once the kernel energy of the error lies on the slow eigenvalues; the kernel itself moves by an amount set by the path length and the higher derivatives of the network (Part~\ref{part:path}).
The error at any time is then predicted from the initial state, two-sided and with a proved error, at a cost that grows with the length of the window (Section~\ref{sec:predict}); bounds that hold at every time without computing the prediction are lower bounds, or upper bounds under stated assumptions on the path.
None of it needs a model of the network's approximation or generalization.

\paragraph{Applications.}
Part~\ref{part:apps} collects the applications of the framework.
The first, Section~\ref{sec:designcompare}, is the closed-form case of the comparison principle of Corollary~\ref{cor:compare}: under gradient flow the error of a design never exceeds its initial error, and the theory bounds from below how fast another design can close its own gap, so a design that starts closer to its target provably wins within a training budget (Theorem~\ref{thm:designcompare}).
For learned time steppers, where the labels are the part of the flow map that a known part leaves to the network, this proves that for every training budget the direct step is worse than the Euler step below an explicit time step (Corollary~\ref{cor:eulerdirect}).
Numerical checks on sine targets, the Lorenz-63 comparison of the two steppers, and gradient descent at the edge of stability are to be added.
Table~\ref{tab:symbols} in Appendix~\ref{app:symbols} lists the symbols.

% ===============================================================
\part{The error at one time during training}\label{part:onetime}
% ===============================================================

% ===============================================================
\section{Setting}\label{sec:setting}
% ===============================================================

The data are $n$ training inputs $x_1,\dots,x_n$ in a compact set $\Aset$, the data set, and labels $Y_1,\dots,Y_n \in \R^{n_{\mathrm y}}$.
The pair $(X, Y)$ is the dataset and does not change during training.
Nothing below assumes anything about how the labels are produced.
For a learned time stepper they are the part of the flow map of a dynamical system that a known part of the step leaves to the network, and $\Aset$ is an invariant set of the dynamics (Part~\ref{part:apps}).

A sampled function is a collection of one point per training input, $u = (u_1,\dots,u_n)$ with $u_i \in \R^{n_{\mathrm y}}$.
A map $u$ on $\Aset$ defines the sampled function $(u(x_1),\dots,u(x_n))$, written with the same letter.
Sampled functions carry the inner product and norm
\begin{equation}\label{eq:inner}
\ip{a}{b} = \frac1n\sum_{i=1}^n a_i\cdot b_i, \qquad \|a\| = \ip{a}{a}^{1/2}.
\end{equation}
For a map $u$ on $\Aset$,
\begin{equation}\label{eq:rmssup}
\|u\| \le \supA{u} = \sup_{x\in\Aset}|u(x)|,
\end{equation}
because every training input lies in $\Aset$.

\paragraph{Network, output scale and error.}
The network $\Net$ maps inputs to $\R^{n_{\mathrm y}}$.
Its output enters the prediction at a scale $c_p \neq 0$: the prediction at the training inputs is the sampled function
\begin{equation}\label{eq:YYhat}
\hat Y = c_p\,\Net, \qquad E = \hat Y - Y,
\end{equation}
and the network is trained on the labels $Y/c_p$.
The usual case is $c_p = 1$.
Other values arise when the targets are normalized to unit size, $c_p = \|Y\|$, and in learned time steppers, where the network may supply a difference quotient with $c_p$ equal to the time step.
$Y$, $\hat Y$ and $E$ are in the units of the labels, and every bound below is in these units.
$\hat Y = \hat Y(t)$ is the state of the network at time $t$; since $Y$ is constant, $\dot E = \dot{\hat Y}$: the rate of change of the error is the rate of change of the outputs.

% ===============================================================
\section{The training kernel}\label{sec:kernel}
% ===============================================================

This section defines the objects that describe how training changes the error.
Everything is evaluated at one time during training.
Nothing below depends on how the kernel changes over time; that is the subject of Part~\ref{part:path}.

\paragraph{Gradient features and kernel.}
The network $\Net$ has parameters $\theta\in\R^{d_\theta}$.
Let $\Pi$ be a diagonal matrix of positive per-parameter step sizes, such as the preconditioner of Adam.
The gradient feature of the training input $x_i$ is
\begin{equation}\label{eq:feature}
\Psi_i = \nabla_\theta\Net(x_i)\,\Pi^{1/2} \in \R^{n_{\mathrm y}\times d_\theta}.
\end{equation}
Let $\Psi$ be the matrix whose $i$th block of $n_{\mathrm y}$ rows is $\Psi_i$.
For a parameter direction $g\in\R^{d_\theta}$, $\Psi g$ is a sampled function: the change of the network output at the training inputs produced by the parameter step $\Pi^{1/2}g$.
The training kernel is
\begin{equation}\label{eq:kernel}
K = \frac1n\,\Psi\Psi^\top .
\end{equation}
For a sampled function $w$, set $g_w = \tfrac1n\Psi^\top w = \tfrac1n\sum_i\Psi_i^\top w_i$.
Then $Kw = \Psi g_w$ and $\ip{w}{Kw} = |g_w|^2$.
The kernel is symmetric with respect to Eq.~\eqref{eq:inner} and positive semidefinite.
Its eigenvalues are $\lambda_1\ge\lambda_2\ge\dots\ge0$, with eigenvectors $v_i$ orthonormal in Eq.~\eqref{eq:inner}.
For a sampled function $u$, write $\hat u_i = \ip{v_i}{u}$ for its components.
The kernel size of a sampled function $\varphi$ is $\|\varphi\|_K = \ip{\varphi}{K\varphi}^{1/2} = |g_\varphi|$.
Since $\ip{\varphi}{\Psi g} = g\cdot g_\varphi$, the kernel size is the largest change of $\ip{\varphi}{\Psi g}$ per unit length of $g$.

\paragraph{The error and its rate of change.}
Let the network be continuously differentiable in $\theta$.
By the chain rule,
\begin{equation}\label{eq:Edot}
\dot E = c_p\,\Psi\,\hat g,
\end{equation}
where $\hat g = \Pi^{-1/2}\dot\theta$ is the rescaled parameter velocity.
This holds for every update rule.
Under preconditioned gradient flow on the loss $\tfrac12\|\Net - Y/c_p\|^2$, $\hat g = -g_E/c_p$, and
\begin{equation}\label{eq:Edotgf}
\dot E = -K E .
\end{equation}
The component of the error along $v_i$ then decays at the instantaneous rate $\lambda_i$, with $v_i$ and $\lambda_i$ taken at the current time: $\frac{d}{dt}\ip{v_i}{E} = -\lambda_i\ip{v_i}{E}$ there.

\paragraph{The slow part of a sampled function.}
For $\gamma > 0$ and a sampled function $u$, define
\begin{equation}\label{eq:ugamma}
u_\gamma = \gamma(K+\gamma I)^{-1}u = \sum_i \frac{\gamma}{\lambda_i+\gamma}\,\hat u_i\,v_i, \qquad \Upsilon_u(\gamma) = \sum_i\frac{\lambda_i\,\hat u_i^2}{(\lambda_i+\gamma)^2}.
\end{equation}
This is a theoretical object: a family of linear maps $u \mapsto u_\gamma$, indexed by $\gamma$, built from the kernel at the current time.
Nothing is integrated over training.
We call $u_\gamma$ the slow part of $u$ at $\gamma$, for two reasons.

In the eigenbasis, $u_\gamma$ keeps the component of $u$ along $v_i$ with the weight $\gamma/(\lambda_i+\gamma) \in (0,1]$.
By Eq.~\eqref{eq:Edotgf}, $\lambda_i$ is the rate at which gradient flow removes the component along $v_i$ at the current time.
Components on directions with $\lambda_i \ll \gamma$, which gradient flow removes slowly, are kept almost fully; those with $\lambda_i \gg \gamma$ are reduced by the factor $\gamma/\lambda_i$; at $\lambda_i = \gamma$ the weight is one half.
The weight is a smooth function of $\lambda_i/\gamma$, so no eigenvector is labelled slow or fast; $\gamma$ only sets the rate against which the weights are measured.

In terms of the gradient features, $u_\gamma$ is the residual of the ridge fit of $u$ by the features, $\min_g \|u - \Psi g\|^2 + \gamma|g|^2$, whose fitted part is $K(K+\gamma I)^{-1}u$: the part of $u$ that the features cannot supply when parameter motion is charged at the price $\gamma$.
Proposition~\ref{prop:floorstep} (with $u$ in place of $E$) turns the price into a length: $\|u_\gamma\|$ is the smallest value of $\|u + c_p\Psi g\|$ over $|g| \le s_\gamma = \Upsilon_u(\gamma)^{1/2}/|c_p|$, and $u_\gamma$ is the minimizer, the smallest error that a parameter step of that length reaches in the first-order Taylor expansion of the network about the current parameters; by Proposition~\ref{prop:taylorstep}, the error that such steps actually reach differs from it by at most a term of second order in the step length.
So each $\gamma$ corresponds to a step length, a small $\gamma$ to a long step, and as $\gamma \to 0$, $u_\gamma$ tends to the part of $u$ outside the range of $K$, which the first-order term of no step can remove.
This is the reading the bounds use.
The quantity $\Upsilon_u(\gamma) = \ip{w}{Kw}$ with $w = (K+\gamma I)^{-1}u$ is the squared rescaled length of the shortest parameter direction $g$ whose network-output change $\Psi g$ is the removed part $u - u_\gamma = Kw = \Psi g_w$, so that the corresponding change of the error, $c_p\Psi g$, needs a step of length $s_\gamma$; it is at most $\|u\|^2/(4\gamma)$.
For the error we write $E_\gamma$ and $\Upsilon(\gamma) = \Upsilon_E(\gamma)$.
$E_\gamma$ is the slow part of the error: the part that a step of length $s_\gamma$ cannot remove.

% ===============================================================
\section{Roughness functions}\label{sec:seminorm}
% ===============================================================

\begin{definition}[Roughness function, local complexity function]\label{def:lcf}
A roughness function is a nonzero seminorm $\Rgh$ on sampled functions.
Its norm constant is
\begin{equation}\label{eq:CR}
\CR = \max_{u\neq 0}\frac{\Rgh(u)}{\|u\|}.
\end{equation}
A roughness function that depends only on the training inputs, and not on the labels or the network, is a local complexity function.
\end{definition}

A seminorm can vanish on nonzero sampled functions.
Those are the sampled functions that it counts as perfectly smooth.
Two consequences of the definition are used throughout.
The reverse triangle inequality gives $|\Rgh(u) - \Rgh(v)| \le \Rgh(u - v)$.
For a map $u$ on $\Aset$, Eqs.~\eqref{eq:CR} and~\eqref{eq:rmssup} give
\begin{equation}\label{eq:Rsup}
\Rgh(u) \le \CR\,\supA{u}.
\end{equation}
Every lower bound below holds for every roughness function, including seminorms built from the network and from the current error.
In particular, for any sampled function $\varphi$, the rank-one seminorm $u \mapsto |\ip{\varphi}{u}|$ is a roughness function; it measures the size of $u$ along $\varphi$, so size along a direction is a roughness in the sense used here.
The best roughness function, the Krylov functional of Section~\ref{sec:step}, is of this kind and depends on both the error and the kernel.
The restriction to local complexity functions serves one purpose: it makes the roughness of the labels, defined next, a property of the data alone, which can be computed before training.

\begin{definition}[Absolute and relative roughness of a dataset]\label{def:absolute}
For a roughness function $\Rgh$ and a dataset $(X, Y)$, the absolute roughness of the dataset under $\Rgh$ is $\Rgh(Y)$, the roughness, in the units of the labels, of what the network must supply.
More generally, for a sampled function $u \neq 0$, the absolute roughness of $u$ is $\Rgh(u)$ and the relative roughness of $u$ is $\Rgh(u/\|u\|)$, the roughness of its shape, a number in $[0, \CR]$.
Since a seminorm is absolutely homogeneous,
\begin{equation}\label{eq:sizeshape}
\Rgh(u) = \|u\|\;\Rgh\big(u/\|u\|\big):
\end{equation}
absolute roughness is size times relative roughness.
\end{definition}

When $\Rgh$ is a local complexity function, $\Rgh(Y)$ depends only on the training inputs and labels, not on the network.
The network itself never sees $Y$: it is trained on $Y/c_p$, or on targets normalized to unit size, and the roughness of its own labels is $\Rgh(Y/c_p) = \Rgh(Y)/|c_p|$, which for normalized targets is the relative roughness.
Every bound below is in the units of the labels and is stated in terms of the absolute roughness.
The distinction matters when two problems with labels of different sizes are compared, as for the learned time steppers of Part~\ref{part:apps}: measuring roughness on the network's own labels keeps the relative roughness and drops the size.

\begin{definition}[Peak feature roughness]\label{def:sigma}
The peak feature roughness under $\Rgh$ is
\begin{equation}\label{eq:sigma}
\sR = \max_{g\neq 0}\frac{\Rgh(\Psi g)}{|g|},
\end{equation}
the largest roughness of the output change produced by a unit step in any parameter direction.
\end{definition}

The peak feature roughness is the only property of the network, apart from the eigenvalues of $K$ and the roughness of its current output, that enters the lower bounds of Part~\ref{part:onetime}.

\subsection{Examples}\label{sec:examples}

A roughness function is any seminorm: $\Rgh(u) \ge 0$, $\Rgh(\alpha u) = |\alpha|\,\Rgh(u)$, and $\Rgh(u+v) \le \Rgh(u) + \Rgh(v)$.
The three examples below count as smooth, respectively, the functions constant on each cluster, the span of the first $m$ graph eigenvectors, and the constants.
None of them is needed for the bounds.
Section~\ref{sec:step} shows that, for one training step, the best bound over all of them is attained by a rank-one seminorm $u \mapsto |\ip{\varphi}{u}|$ built from the error and the kernel and computed by the Krylov series, and the bounds of Section~\ref{sec:lower} hold for every roughness function at once.
The examples serve to show what roughness means and to guide the argument.

\paragraph{Clusters.}
Partition the training indices into clusters $C$ of nearby inputs, and let $\bar u_C$ be the mean of the values $u_i$ over $i \in C$.
The cluster seminorm is $\Rgh_{\mathrm{cl}}(u) = \big(\tfrac1n\sum_C\sum_{i\in C}|u_i - \bar u_C|^2\big)^{1/2}$, the variation of $u$ within clusters.
It has the form $\Rgh(u) = \|Pu\|$ for a nonzero orthogonal projector $P$ in Eq.~\eqref{eq:inner}, a projector seminorm.
Every projector seminorm has $\CR = 1$, and $\sR^2 = \lambda_{\max}(PKP)$, the largest value of $\ip{u}{Ku}/\|u\|^2$ over the range of $P$: the fastest rate at which gradient flow can remove a sampled function that the seminorm counts as rough.
(With $B = n^{-1/2}P\Psi$, $\Rgh(\Psi g)^2 = |Bg|^2$, and $B^\top B$ and $BB^\top = PKP$ have the same nonzero eigenvalues.)
Subtracting a least-squares polynomial of degree $p$ in the input on each cluster, in place of the mean, gives a seminorm for every $p$.

\paragraph{Manifold high-pass.}
Connect each training input to its $k_{\mathrm g}$ nearest neighbours with weights $\omega_{ij} = \exp(-(d_{ij}/r_{\mathrm g})^2)$, where $d_{ij} = |x_i - x_j|$ and $r_{\mathrm g}$ is the median distance to the $k_{\mathrm g}$th neighbour, symmetrized by $\omega_{ij} = \max(\omega_{ij},\omega_{ji})$.
The graph Laplacian $L_{\mathrm g}$ satisfies $a^\top L_{\mathrm g}a = \tfrac12\sum_{i,j}\omega_{ij}(a_i - a_j)^2$, and its eigenvectors $v^{\mathrm g}_1,\dots,v^{\mathrm g}_n$, orthonormal in Eq.~\eqref{eq:inner}, have eigenvalues $0 = \mu^{\mathrm g}_1 \le \dots \le \mu^{\mathrm g}_n$; eigenvectors with small eigenvalues change slowly between neighbouring inputs.
For an integer $1 \le m < n$, let $P_m$ remove from each output component of a sampled function its projection onto $v^{\mathrm g}_1,\dots,v^{\mathrm g}_m$.
The manifold high-pass is $\Rgh_m(u) = \|P_m u\|$, the part of $u$ above the $m$th graph-Laplacian eigenvector.
It is a projector seminorm, so $C_{\Rgh_m} = 1$ and $\sigma_{\Rgh_m}^2 = \lambda_{\max}(P_mKP_m)$.
Like the cluster seminorm, it measures variation between neighbouring inputs averaged over the data, but with a smooth cutoff set by the graph spectrum rather than by a partition, and as $n$ grows the graph eigenvectors approach the eigenfunctions of the Laplace--Beltrami operator of a manifold from which the inputs are drawn; nothing below needs this limit.

\paragraph{Lipschitz seminorm.}
The Lipschitz seminorm on the training inputs is $\Lipn(u) = \max_{i\neq j}|u_i - u_j|/|x_i - x_j|$, and $\Lipn(u) \le \Lip(u)$ for a map $u$ on $\Aset$, where $\Lip(u)$ is its Lipschitz constant on $\Aset$.
Its norm constant is at least $\sqrt n/d_{\min}$, where $d_{\min} = \min_{i\neq j}|x_i - x_j|$; the sampled function that equals $\sqrt n$ at one point of a closest pair and zero elsewhere shows this.
The lower bounds of Section~\ref{sec:lower} divide by the norm constant, so the Lipschitz seminorm gives a weak bound when the inputs are many or dense.

\subsection{The best roughness function: one step and the Krylov series}\label{sec:step}\label{sec:krylov}

This subsection asks what one training step can do and which roughness function gives the best lower bound on the error after it.
It uses only the kernel of Section~\ref{sec:kernel} and the definitions above.

A discrete training step changes the parameters by $\Delta\theta$, with $\Pi$ held fixed during the step.
Its rescaled step is $g = \Pi^{-1/2}\Delta\theta$.
By the fundamental theorem of calculus, the error changes by
\begin{equation}\label{eq:avgstep}
\Delta E = c_p\,\bar\Psi\,g, \qquad \bar\Psi = \int_0^1 \Psi(\theta + \tau\Delta\theta)\,d\tau,
\end{equation}
where $\bar\Psi$ is the matrix of gradient features averaged along the step.
The statements below hold for any one feature matrix, written $\Psi$, with kernel $K = \tfrac1n\Psi\Psi^\top$.
Two choices are used, and neither assumes that the features stay fixed during training.
With $\Psi = \bar\Psi$, the averaged features of a step actually taken, Eq.~\eqref{eq:avgstep} is exact and every lower bound below holds for that step with no remainder; $\bar\Psi$ is known only after the step.
With $\Psi = \Psi(t)$, the features at the current parameters, $c_p\Psi g$ is the first-order term of the Taylor expansion of the output change about $\theta(t)$, and Proposition~\ref{prop:taylorstep} below bounds the remainder.
The first-order output changes of steps of length at most $s$ form the set $\mathcal{D}_s = \{c_p\Psi g : |g| \le s\}$, and the smallest error among them is
\begin{equation}\label{eq:Fs}
F(s) = \min_{|g|\le s}\|E + c_p\Psi g\| .
\end{equation}
For a roughness function $\Rgh$, the step bound is
\begin{equation}\label{eq:bR}
b_\Rgh(s) = \frac{1}{\CR}\Big(\Rgh(E) - |c_p|\,s\,\sR\Big),
\end{equation}
the roughness of the error minus the most roughness that a step of length $s$ can remove, divided by the norm constant.

\begin{theorem}[Supremum over roughness functions]\label{thm:stepsup}
Let $s \ge 0$, let $\hat e_i = \ip{v_i}{E}$, and let $\Upsilon(0^+) = \lim_{\gamma\to0}\Upsilon(\gamma) = \sum_{\lambda_i>0}\hat e_i^2/\lambda_i$.
\begin{enumerate}[label=(\alph*)]
\item For every roughness function $\Rgh$, $b_\Rgh(s) \le F(s)$.
\item If $F(s) > 0$, then $\sup_\Rgh b_\Rgh(s) = F(s)$, where the supremum is over all roughness functions. It is attained by the rank-one seminorm $\Rgh^\star(u) = |\ip{\varphi^\star}{u}|$ with $\varphi^\star = E^\star/\|E^\star\|$, where $E^\star = E + c_p\Psi g^\star$ is the minimizer in Eq.~\eqref{eq:Fs}. The minimizer $E^\star$ is unique.
\item If $0 < (c_ps)^2 < \Upsilon(0^+)$, let $\gamma > 0$ solve $\Upsilon(\gamma) = (c_ps)^2$. Then $E^\star = E_\gamma = \gamma(K+\gamma I)^{-1}E$ and
\begin{equation}\label{eq:ridgestep}
F(s)^2 = \sum_i\frac{\gamma^2\hat e_i^2}{(\lambda_i+\gamma)^2}.
\end{equation}
If $(c_ps)^2 \ge \Upsilon(0^+)$, then $F(s)^2 = \sum_{\lambda_i=0}\hat e_i^2$. For $s = 0$, $F(0) = \|E\|$.
\item If $E \neq 0$, then $F(s) \ge \|E\| - |c_p|\,s\,\|\hat E\|_K$ for every $s$, and $F(s) = \|E\| - |c_p|\,s\,\|\hat E\|_K + O(s^2)$ as $s \to 0$, where $\hat E = E/\|E\|$.
\end{enumerate}
\end{theorem}

Part (a) says that every roughness function gives a lower bound on the error after the step.
Part (b) says that the supremum of these bounds over all roughness functions is exactly $F(s)$, and that a single linear functional attains it.
Part (c) says that the best functional points along the slow part $E_\gamma$ of the error, at a rate $\gamma$ set by the step length.
Part (d) says that for short steps the bound is first order in the step length; Proposition~\ref{prop:rates}(c) in Section~\ref{sec:lower} states the same rate for the time derivative of the error.
Restricted to local complexity functions, which are fixed by the training inputs, the supremum is at most $F(s)$.
Parts (a) and (b) are the duality between the distance from a point to a compact convex set and the support function of that set \cite{hiriart2001}, and part (c) is the resulting ridge problem.
What the theorem adds is the identification of the dual variables with roughness functions.

With $\Psi = \Psi(t)$, $F(s)$ is the smallest error of the first-order Taylor expansion of the step about $\theta(t)$.
The error that steps actually reach differs from it by the Taylor remainder, which the next proposition bounds with no assumption on how the features change along the step.
Let $B_s = \{\theta(t) + \Pi^{1/2}g : |g|\le s\}$ and
\begin{equation}\label{eq:H2bar1}
\bar H_2(s) = \max_{\theta'\in B_s}\Big(\tfrac1n\sum_i\big\|\nabla^2_\theta\Net(x_i;\theta')\big[\Pi^{1/2}\cdot,\Pi^{1/2}\cdot\big]\big\|^2\Big)^{1/2},
\end{equation}
the largest root-mean-square size of the second derivative of the network over $B_s$, in the rescaled parameters; here $\|\cdot\|$ of a bilinear form is $\sup|B[u_1,u_2]|$ over unit $u_1,u_2$, and for vector outputs the largest such norm of $\phi\cdot\nabla^2_\theta\Net(x_i)$ over unit output vectors $\phi$.
With $\Pi = I$, as in Part~\ref{part:path}, $B_s$ is the Euclidean ball of radius $s$ and $\bar H_2(s)$ is the quantity of Lemma~\ref{lem:drift}.

\begin{proposition}[Steps actually taken]\label{prop:taylorstep}
Let $\Net$ be twice continuously differentiable in $\theta$, let $\Psi = \Psi(t)$, and let $s \ge 0$.
\begin{enumerate}[label=(\alph*)]
\item Every step of rescaled length at most $s$ leaves an error $E_{\mathrm{new}}$ with
\[
\|E_{\mathrm{new}}\| \ \ge\ F(s) - \tfrac12|c_p|\,\bar H_2(s)\,s^2 .
\]
\item Any minimizer $g^\star$ in Eq.~\eqref{eq:Fs}, which exists by compactness, gives a step that leaves an error with $\|E_{\mathrm{new}}\| \le F(s) + \tfrac12|c_p|\,\bar H_2(s)\,s^2$.
\end{enumerate}
So the smallest error that steps of length at most $s$ actually reach lies within $\tfrac12|c_p|\bar H_2(s)s^2$ of $F(s)$.
With $\Psi = \bar\Psi$ of a step actually taken, of length $s$, its error is at least $F(s)$ computed from $\bar\Psi$, with no remainder.
\end{proposition}

\begin{proof}
For a step $\Delta\theta = \Pi^{1/2}g$, Taylor's theorem with integral remainder at each training input gives $E_{\mathrm{new}} = E + c_p\Psi g + r$ with $r = c_p\int_0^1(1-\tau)\,v(\tau)\,d\tau$, where $v(\tau)$ is the sampled function with values $\nabla^2_\theta\Net(x_i;\theta + \tau\Delta\theta)[\Delta\theta,\Delta\theta]$.
By Minkowski's inequality for the integral, $\|r\|\le|c_p|\int_0^1(1-\tau)\|v(\tau)\|\,d\tau$, and $\|v(\tau)\|\le\bar H_2(s)|g|^2$ because $\theta + \tau\Delta\theta\in B_s$; so $\|r\|\le\tfrac12|c_p|\bar H_2(s)s^2$.
(a) is $\|E_{\mathrm{new}}\|\ge\|E + c_p\Psi g\| - \|r\|\ge F(s) - \|r\|$, and (b) is $\|E_{\mathrm{new}}\|\le\|E + c_p\Psi g^\star\| + \|r\| = F(s) + \|r\|$.
The last statement is Eq.~\eqref{eq:avgstep} and $\|E + c_p\bar\Psi g\|\ge F(s)$.
\end{proof}

The expansion can be carried to any order, with the remainder bounded by the next derivative over $B_s$ in the same way; the first order is the one at which the best roughness function has the closed form of Theorem~\ref{thm:stepsup}.

\paragraph{The Krylov series.}
The best functional $\varphi^\star$ of Theorem~\ref{thm:stepsup} is the roughness function tuned to the dataset and to the network at time $t$.
This paragraph writes it, and its bound, in terms of the three objects that the theory connects: the labels $Y$, which are the dataset and do not change; the output $\hat Y(t)$, which is the state of the network; and the kernel $K(t)$, which is what the network can do next.
The supremum in Theorem~\ref{thm:stepsup}(b) is attained by a rank-one seminorm, $E = \hat Y - Y$, and the sign of $\varphi$ is free, so
\begin{equation}\label{eq:Fdataset}
F(s) \ =\ \max_{\|\varphi\|\le 1}\Big(\ip{\varphi}{Y} - \ip{\varphi}{\hat Y} - |c_p|\,s\,\|\varphi\|_K\Big).
\end{equation}
The three terms are the roughness of the labels along $\varphi$, the part of it that the current output already carries, and the most that the first-order term of a step of length $s$ can supply along $\varphi$.
This is the step bound $b_\Rgh(s)$ for the rank-one roughness function $u \mapsto |\ip{\varphi}{u}|$, whose norm constant is $\|\varphi\|$ and whose peak feature roughness is $\|\varphi\|_K$, with the signed functional $\ip{\varphi}{\cdot}$ in place of the seminorm, which keeps the sign of the output's component along $\varphi$; Section~\ref{sec:slow} uses this signed form to make the separation of the roughness of the labels from that of the output exact, Eq.~\eqref{eq:mainexact}.
The supremum over all roughness functions is reached among the rank-one ones, so no seminorm has to be chosen.

What remains is to find $\varphi^\star$ without solving the minimization in Eq.~\eqref{eq:Fs}.
By Theorem~\ref{thm:stepsup}(c), when the step is too short to remove the whole reachable part of the error, $(c_ps)^2 < \Upsilon(0^+)$, $\varphi^\star$ is a positive multiple of $(K+\gamma I)^{-1}(\hat Y - Y)$, the direction of the slow part of the error; the maximizer of Eq.~\eqref{eq:Fdataset} is $-\varphi^\star$, since that equation is written with $Y - \hat Y$.
On the finitely many eigenvalues of $K$ on which $\hat Y - Y$ has a component, the function $\lambda \mapsto 1/(\lambda+\gamma)$ agrees with a polynomial.
So $\varphi^\star$ is a polynomial in $K$ applied to $\hat Y - Y$: it lies in the span of $\hat Y - Y$, $K(\hat Y - Y)$, $K^2(\hat Y - Y)$, and so on, all formed with the kernel at time $t$.
When the step is long enough, $(c_ps)^2 \ge \Upsilon(0^+)$, $\varphi^\star$ is along the component of $\hat Y - Y$ outside the range of $K$, which is also a polynomial in $K$ applied to $\hat Y - Y$ (Theorem~\ref{thm:stepkrylov}).
The first of these is the error, the gap between the output and the labels.
By Eq.~\eqref{eq:Edotgf}, the second is minus the rate at which gradient flow closes that gap at time $t$.
Each further one is one more application of the kernel at time $t$.
Nothing about the kernel at later times enters.
This suggests searching for the best functional in the Krylov subspaces
\begin{equation}\label{eq:krylovstep}
V_M(Y, \hat Y, K) = \operatorname{span}\big\{\hat Y - Y,\ K(\hat Y - Y),\ K^2(\hat Y - Y),\ \dots,\ K^{M-1}(\hat Y - Y)\big\},
\end{equation}
each of which costs one more application of $K$, that is, one product with $\Psi^\top$ and one with $\Psi$.
The best bound from a functional in $V_M$ is Eq.~\eqref{eq:Fdataset} restricted to $V_M$,
\begin{equation}\label{eq:FMstep}
F_M(s) \ =\ \max_{\varphi\in V_M(Y,\hat Y,K),\ \|\varphi\|\le1}\Big(\ip{\varphi}{Y} - \ip{\varphi}{\hat Y} - |c_p|\,s\,\|\varphi\|_K\Big).
\end{equation}
$F_M(s)$ is a function of the labels, the current output, the current kernel and the step length, and of nothing else.
The dataset enters through $Y$ in two places: as the roughness $\ip{\varphi}{Y}$ that the labels demand along $\varphi$, and in the subspace $V_M$, which is built from the gap $\hat Y - Y$.

\begin{theorem}[Krylov series]\label{thm:stepkrylov}
Let $F(s) > 0$ and $E = \hat Y - Y$.
\begin{enumerate}[label=(\alph*)]
\item $F_M(s)$ is the larger of $0$ and the supremum of $b_\Rgh(s)$ over the rank-one seminorms $u \mapsto |\ip{\varphi}{u}|$ with nonzero $\varphi \in V_M$, and $F_1(s) \le F_2(s) \le \dots \le F(s)$. In particular $F_1(s) = \max\big(0,\ \|\hat Y - Y\| - |c_p|\,s\,\|\hat E\|_K\big)$ with $\hat E = E/\|E\|$.
\item Let $q_E$ be the number of distinct eigenvalues of $K$, including zero, on which $\hat Y - Y$ has a nonzero component. Then $\varphi^\star \in V_{q_E}$, so $F_M(s) = F(s)$ for all $M \ge q_E$.
\end{enumerate}
\end{theorem}

Theorem~\ref{thm:stepkrylov} says that the Krylov values increase to $F(s)$ and reach it after at most $q_E$ terms.
Written as a series,
\begin{equation}\label{eq:series}
F(s) = F_1(s) + \sum_{M\ge1}\big(F_{M+1}(s) - F_M(s)\big).
\end{equation}
The first term is the first-order bound of Theorem~\ref{thm:stepsup}(d): the distance from the output to the labels, less what a step of length $s$ can close along that direction.
Each further term adds the direction of one more application of the kernel to the gap $\hat Y - Y$.
Proposition~\ref{prop:krylovrate} in Appendix~\ref{app:aux} bounds the gap $F(s) - F_M(s)$ above by a multiple of the distance $d_M$ from $\varphi^\star$ to $V_M$ and below by a multiple of $d_M^2$, and, when $0 < (c_ps)^2 < \Upsilon(0^+)$, shows that this distance decays geometrically at the Chebyshev rate of the conjugate gradient method.
A long step gives a slow rate, so the series converges slowly for steps that nearly fit the error.

\paragraph{Output scale.}
The scale $c_p$ at which the network output enters the prediction is a choice: the same prediction is obtained with the scale $c\,c_p$ and the network $\Net/c$, for example by normalizing the targets.
Every bound of this subsection and of Section~\ref{sec:lower} is unchanged by this choice, with the rate $\gamma$ rescaled to $\gamma/c^2$ where it appears (Proposition~\ref{prop:scale} in Appendix~\ref{app:aux}), so the bounds compare problems by quantities in the units of the labels only, whatever the scale at which each network is trained.
What does change is the speed of plain gradient flow: the kernel is divided by $c^2$, which acts as a change of the learning rate.
Adam is approximately invariant to the scale.

% ===============================================================
\section{Lower bounds from roughness}\label{sec:lower}
% ===============================================================

This section proves that a rough dataset forces a large error that decreases slowly; the roughness function that gives the best bound is the one of Section~\ref{sec:step}.
As in Section~\ref{sec:setting}, $Y$ collects the labels, $\hat Y = c_p\Net$ the network outputs at the training inputs, and $E = \hat Y - Y$ the error, all in the units of the labels.
The pair $(X, Y)$ is the dataset and does not change during training.
$\hat Y = \hat Y(t)$ is the state of the network at time $t$.

\paragraph{The objects in the argument.}
Five objects from Sections~\ref{sec:kernel} and~\ref{sec:seminorm} carry the argument.
\begin{itemize}
\item \emph{Roughness function.} $\Rgh$ is a nonzero seminorm on sampled functions (Definition~\ref{def:lcf}), with norm constant $\CR = \max_u\Rgh(u)/\|u\|$, Eq.~\eqref{eq:CR}.
The norm constant converts roughness into size: $\|u\| \ge \Rgh(u)/\CR$.
It equals 1 for a projector seminorm.
$\Rgh(Y)$ is the absolute roughness of the dataset (Definition~\ref{def:absolute}).
\item \emph{Gradient features.} $\Psi$ is the matrix of gradient features, Eq.~\eqref{eq:feature}.
For a parameter direction $g$, $\Psi g$ is the change of the network output at the training inputs produced by the parameter step $\Pi^{1/2}g$.
The rescaled parameter velocity is $\hat g = \Pi^{-1/2}\dot\theta$, and a parameter step $\Delta\theta$ has rescaled length $|\Pi^{-1/2}\Delta\theta|$.
\item \emph{Kernel.} $K = \tfrac1n\Psi\Psi^\top$ is the training kernel, Eq.~\eqref{eq:kernel}, with eigenvalues $\lambda_1 \ge \lambda_2 \ge \dots \ge 0$ and orthonormal eigenvectors $v_i$.
Under gradient flow, the component of the error along $v_i$ decays at the instantaneous rate $\lambda_i$, with $v_i$ and $\lambda_i$ taken at the current time.
A small $\lambda_i$ is a slow direction.
\item \emph{Peak feature roughness.} $\sR = \max_g\Rgh(\Psi g)/|g|$ is the most roughness that a unit parameter step can add to the output, Eq.~\eqref{eq:sigma}.
It depends on the network at the current time.
For a step of rescaled length at most $s$, $\bar\sigma_\Rgh(s)$ is the largest value of $\sR$ at any parameters that the step can reach.
\item \emph{Slow part of the error and leak term.} $E_\gamma = \gamma(K+\gamma I)^{-1}E$, Eq.~\eqref{eq:ugamma}, is the smallest error that a parameter step of rescaled length $s_\gamma = \Upsilon(\gamma)^{1/2}/|c_p|$ reaches in the first-order Taylor expansion about the parameters at time $t$ (Proposition~\ref{prop:floorstep}), and the error actually reached is within $\tfrac12|c_p|\bar H_2(s_\gamma)s_\gamma^2$ of it (Proposition~\ref{prop:taylorstep}); a long step is a small $\gamma$.
In the eigenbasis, $E_\gamma$ keeps the component of $E$ along $v_i$ with the smooth weight $\gamma/(\lambda_i+\gamma)$, so no eigenvector is labelled slow or fast.
$\Upsilon(\gamma) = \sum_i\lambda_i\hat e_i^2/(\lambda_i+\gamma)^2$, with $\hat e_i = \ip{v_i}{E}$, is the squared rescaled length of the shortest $g$ with $\Psi g = E - E_\gamma$, so that the error step that removes this part has length $s_\gamma = \Upsilon(\gamma)^{1/2}/|c_p|$; it is at most $\|E\|^2/(4\gamma)$.
The product $\sR\,\Upsilon(\gamma)^{1/2}$ is the most roughness that the removed part can carry, since the removed part is the network-output change $\Psi g$ of a parameter direction of length $\Upsilon(\gamma)^{1/2}$.
We call it the leak term.
\end{itemize}

\paragraph{The argument in five inequalities.}
Everything else in this section sharpens the following five statements.
\begin{enumerate}[label=(\arabic*)]
\item \emph{The dataset sets the roughness of the error.}
For every roughness function, $\Rgh(E) = \Rgh(\hat Y - Y) \ge \Rgh(Y) - \Rgh(\hat Y)$.
The error is at least as rough as the labels, minus the roughness of the current output.
The labels are fixed, so the network can lower the roughness of the error only by making its own output as rough as the labels.
\item \emph{The output moves only along the gradient features.}
By Eq.~\eqref{eq:Edot}, $\dot{\hat Y} = \dot E = c_p\Psi\hat g$ for every update rule.
Every change of the output, and hence of the error, is a combination of the columns of $\Psi$.
Under gradient flow, $\dot E = -KE$, and the component of $E$ along $v_i$ decays at the instantaneous rate $\lambda_i$.
\item \emph{Each feature is at most $\sR$ rough.}
By the definition of $\sR$, $\Rgh(\Psi g) \le \sR\,|g|$.
Hence the roughness of the output changes at rate at most $|c_p|\,\sR\,|\hat g|$ (Proposition~\ref{prop:rates}).
The network can supply the roughness $\Rgh(Y)$ that the labels demand only at a bounded rate per unit of parameter motion, and until it has done so, the error is rough by (1).
\item \emph{Fast eigenvectors are smooth.}
Applied to $Kv_i = \lambda_i v_i$, (3) gives
\begin{equation}\label{eq:budget-intro}
\lambda_i\,\Rgh(v_i)^2 \ \le\ \sR^2
\end{equation}
(Corollary~\ref{cor:budget}).
An eigenvector of roughness $\Rgh(v_i)$ has eigenvalue at most $\sR^2/\Rgh(v_i)^2$.
This is how the slow eigenvalues depend on roughness: the eigenvalues are set by the network, not by the labels, but every eigenvector that is rough under $\Rgh$ is slow.
\item \emph{A rough error lies on the slow eigenvectors, so a rough dataset is learned slowly.}
The removed part $E - E_\gamma$ of the error is the part that gradient flow takes away quickly, and it is smooth: $\Rgh(E - E_\gamma) \le \sR\,\Upsilon(\gamma)^{1/2} \le \sR\|E\|/(2\sqrt\gamma)$.
The roughness of $E$ must therefore be carried by the slow part (Theorem~\ref{thm:fast}):
\begin{equation}\label{eq:roughslow-intro}
\|E_\gamma\| \ \ge\ \frac{1}{\CR}\Big(\Rgh(E) - \sR\,\Upsilon(\gamma)^{1/2}\Big).
\end{equation}
The last term is the leak term.
It is small when $\sR^2 \ll \gamma$.
Combined with (1), the slow part of the error is at least $(\Rgh(Y) - \Rgh(\hat Y))/\CR$ minus the leak term (Theorem~\ref{thm:floor}).
Combined with (3), after a training step of rescaled length $s$ the error is still at least $(\Rgh(Y) - \Rgh(\hat Y(t)) - |c_p|\,s\,\bar\sigma_\Rgh(s))/\CR$ (Theorem~\ref{thm:dataset}).
Theorem~\ref{thm:slow}, the main statement of Part~\ref{part:onetime}, collects these two bounds and identifies the roughness function that makes them exact, the one of Section~\ref{sec:step}.
\end{enumerate}
In one sentence: the labels demand roughness $\Rgh(Y)$; the network can supply roughness only at rate $\sR$ per unit of parameter motion, and only along eigendirections whose eigenvalue is at most $\sR^2$ divided by their squared roughness; so the rough part of the error sits on the slow eigenvalues, and one step can remove only a small part of it.

Section~\ref{sec:rates} proves (2) to (5).
Section~\ref{sec:slow} states the main result, Theorem~\ref{thm:slow}.
Section~\ref{sec:step} has shown that, for one training step, the best of these bounds over all roughness functions is exactly the smallest error that the step reaches to first order, within a second-order remainder of the error actually reached (Theorem~\ref{thm:stepsup} and Proposition~\ref{prop:taylorstep}), and that a Krylov series computes it (Theorem~\ref{thm:stepkrylov}); Theorem~\ref{thm:slow}(c) uses this.
Proofs are in Appendix~\ref{app:proofs}.
Every statement holds for a given network at a given time during training, with the kernel at that time, and uses no model of how well networks approximate or generalize.

\subsection{From the dataset to the slow part of the error}\label{sec:rates}\label{sec:smooth}\label{sec:floor}

By Eq.~\eqref{eq:Edot}, $\dot E = c_p\Psi\hat g$ for every update rule.
Two rate bounds follow.
The component of the error along a sampled function $\varphi$ changes no faster than the kernel size of $\varphi$ allows, and the roughness of the error changes no faster than the peak feature roughness allows.

\begin{proposition}[Rates at one training time]\label{prop:rates}
Let $\theta$ be differentiable at $t$.
For every sampled function $\varphi$ and every roughness function $\Rgh$:
\begin{enumerate}[label=(\alph*)]
\item $\big|\tfrac{d}{dt}\ip{\varphi}{E}\big| \le |c_p|\,\|\varphi\|_K\,|\hat g|$.
\item $\limsup_{\tau\to0}\big|\Rgh(E(t+\tau)) - \Rgh(E(t))\big|/|\tau| \le |c_p|\,\sR\,|\hat g|$.
\item If $E \neq 0$, then $-\tfrac{d}{dt}\|E\| \le |c_p|\,|\hat g|\,\|\hat E\|_K$ with $\hat E = E/\|E\|$. Gradient flow attains equality, $-\tfrac{d}{dt}\|E\| = \|E\|_K^2/\|E\|$.
\end{enumerate}
\end{proposition}

Proposition~\ref{prop:rates} uses the update rule only through $\hat g$, so it holds for Adam with momentum and for minibatch steps (Appendix~\ref{app:momentum}).
Part (c) says that gradient flow is the update of a given length that reduces the error fastest.
Part (b) is statement (3).
Integrated along a training path, part (a) is the starting point of the path calculus of Part~\ref{part:path} (Theorem~\ref{thm:path}).

Statement (4) starts from the same fact about the features.
Every change of the output is a combination of the columns of $\Psi$.
If every parameter step changes the output smoothly, then no combination of steps can change it roughly.

\begin{lemma}[Rough output changes need rough features]\label{lem:mech}
For every roughness function $\Rgh$ and every sampled function $w$,
\begin{equation}\label{eq:mech}
\Rgh(Kw) \le \sR\,\ip{w}{Kw}^{1/2}.
\end{equation}
\end{lemma}

The simplest way to see that rough eigenvectors are slow is a budget that the eigenvectors share.
Take the rank-one seminorm $\Rgh(u) = |\ip{\varphi}{u}|$ of Section~\ref{sec:step}, for a sampled function $\varphi$ with $\|\varphi\|_K > 0$.
Then $\Rgh(v_i) = |\hat\varphi_i|$, and expanding $\varphi = \sum_i\hat\varphi_i v_i$ in the eigenvectors of $K$ gives
\begin{equation}\label{eq:budgetsum}
\sum_i\lambda_i\,\Rgh(v_i)^2 \ =\ \sum_i\lambda_i\,\hat\varphi_i^2 \ =\ \ip{\varphi}{K\varphi} \ =\ \|\varphi\|_K^2 .
\end{equation}
The budget $\|\varphi\|_K^2$ depends only on $\varphi$ and the kernel, not on the labels.
The shares $\lambda_i\Rgh(v_i)^2/\|\varphi\|_K^2$ are nonnegative and sum to one, so none exceeds one: $\lambda_i\,\Rgh(v_i)^2 \le \|\varphi\|_K^2$, and an eigenvector of roughness $\Rgh(v_i) > 0$ has eigenvalue at most $\|\varphi\|_K^2/\Rgh(v_i)^2$.
Since the budget is fixed, the rougher an eigenvector is, the smaller its eigenvalue, and the more slowly gradient flow removes the error along it.
The budget is the peak feature roughness of this seminorm: by Eq.~\eqref{eq:sigma} and $\ip{\varphi}{\Psi g} = g\cdot g_\varphi$, $\sR = \max_{g\neq 0}|g\cdot g_\varphi|/|g| = |g_\varphi| = \|\varphi\|_K$.
The cap is therefore $\lambda_i\,\Rgh(v_i)^2 \le \sR^2$.

Applied to an eigenvector of $K$, Lemma~\ref{lem:mech} gives the same cap for every roughness function.

\begin{corollary}[Fast eigenvectors are smooth]\label{cor:budget}
For every roughness function $\Rgh$ and every eigenpair $(\lambda_i, v_i)$ of $K$,
\begin{equation}\label{eq:budget}
\lambda_i\,\Rgh(v_i)^2 \ \le\ \sR^2 ,
\end{equation}
so $\lambda_i \le \sR^2/\Rgh(v_i)^2$ whenever $\Rgh(v_i) > 0$.
\end{corollary}

Corollary~\ref{cor:budget} is the link between roughness and the spectrum.
The eigenvalues are set by the gradient features, but a direction that is rough under $\Rgh$ cannot be a fast eigendirection, because the features would then have to change the output roughly.
The converse fails: a smooth eigenvector can also be slow, if the features have no component along it.
The next theorem extends the cap from eigenvectors to arbitrary sampled functions.
The slow part $u_\gamma$ replaces a split into slow and fast eigenvectors.

\begin{theorem}[Rough sampled functions are removed slowly]\label{thm:fast}
For every roughness function $\Rgh$, every $\gamma > 0$, and every sampled function $u$:
\begin{enumerate}[label=(\alph*)]
\item The removed part satisfies
\begin{equation}\label{eq:fastsmooth}
\Rgh(u - u_\gamma) \le \sR\,\Upsilon_u(\gamma)^{1/2} \le \frac{\sR}{2\sqrt{\gamma}}\,\|u\|.
\end{equation}
\item The slow part satisfies
\begin{equation}\label{eq:slowlower}
\|u_\gamma\| \ge \frac{1}{\CR}\Big(\Rgh(u) - \sR\,\Upsilon_u(\gamma)^{1/2}\Big).
\end{equation}
\end{enumerate}
\end{theorem}

Part (a) says that the part of a sampled function that the kernel removes quickly is smooth when $\sR^2 \ll \gamma$.
Part (b) says that the roughness which the removed part cannot carry must stay in the slow part.
Part (b) is tight, so the leak term cannot be dropped: for $u = v_1$ and $\Rgh = \|\cdot\|$, for which $\CR = 1$ and $\sR^2 = \lambda_1$, the two sides are equal.
Eq.~\eqref{eq:budget} is part (b) for $u = v_i$ in the limit $\gamma \to 0$.

Applied to the error, with statement (1), Theorem~\ref{thm:fast} gives the floor.

\begin{theorem}[Rough errors are slow: the floor]\label{thm:floor}
Let $(X, Y)$ be any dataset, let $\hat Y = c_p\Net$ be the network output on its inputs, let $E = \hat Y - Y$, and let $\Rgh$ be a roughness function.
At every time during training, for every update rule and every $\gamma > 0$,
\begin{equation}\label{eq:floor}
\|E_\gamma\| \ \ge\ \frac{1}{\CR}\Big(\Rgh(Y) - \Rgh(\hat Y) - \sR\,\Upsilon(\gamma)^{1/2}\Big).
\end{equation}
\end{theorem}

We call the right side of Eq.~\eqref{eq:floor} the floor.
All of its terms are in the units of the labels.
$\Rgh(Y)$ depends only on the data.
$\Rgh(\hat Y)$ is the roughness of what the network currently supplies.
The slow part of the error can fall below $\Rgh(Y)/\CR$ only if the output is itself rough, or through the leak term, and the floor is zero when the output is as rough as its target.
A larger $\gamma$ keeps more of the error in $E_\gamma$ and gives a smaller leak term.
The leak term does not change under a rescaling of the output if $\gamma$ is kept a fixed fraction of $\lambda_1$.

The theorem is a statement about $E$, $K$ and $\Rgh$ at one time; the update rule does not enter.
The weight $\gamma/(\lambda_i+\gamma)$ is called slow with respect to the preconditioned gradient flow with the kernel at time $t$.
Under another update rule, $\|E_\gamma\|$ is still a lower bound on $\|E\|$, but the word slow is then only an analogy.
Nothing in the theorem depends on how the labels were produced; for a learned time stepper, $\Rgh(Y)$ is the local complexity of the stepper (Part~\ref{part:apps}).
Since the weights $\gamma/(\lambda_i+\gamma)$ lie in $(0,1]$, $\|E\| \ge \|E_\gamma\|$, so the floor also bounds the error itself.
The left side of Eq.~\eqref{eq:floor} does not depend on $\Rgh$, so the bound can be taken over any family of roughness functions, even one chosen after looking at the data.
The bound also holds on validation inputs, with the kernel built from the validation inputs.

\subsection{Rough datasets are learned slowly}\label{sec:slow}

The floor depends on the network only through $\sR$, the eigenvalues, and the roughness $\Rgh(\hat Y(t))$ of the current output.
The last of these is not fixed by the data.
By statement (3), however, it can grow only at rate $\sR$ per unit of parameter motion.
This closes the chain.
The roughness of the labels can be matched only by parameter motion, and until it is matched, the error sits on the slow eigenvalues.
Theorem~\ref{thm:dataset} makes this a statement about one training step of a given length.
Theorem~\ref{thm:slow}, the main statement, then collects the bounds and identifies the roughness function of Section~\ref{sec:step} as the one that makes them exact.

\begin{theorem}[A rough dataset costs parameter motion]\label{thm:dataset}
Let $\Rgh$ be a roughness function, let $t$ be any time during training, and let $s \ge 0$.
Let $B_s = \{\theta(t) + \Pi^{1/2}g : |g| \le s\}$ be the set of parameters that a step of rescaled length at most $s$ can reach, with $\Pi$ held fixed during the step, and let $\bar\sigma_\Rgh(s) = \max_{\theta' \in B_s}\sR(\theta')$.
Every step of rescaled length at most $s$ from $\theta(t)$, by any update rule, leaves outputs $\hat Y_{\mathrm{new}}$ with
\begin{equation}\label{eq:dataset}
\|\hat Y_{\mathrm{new}} - Y\| \ \ge\ \frac{\Rgh(Y) - \Rgh(\hat Y(t)) - |c_p|\,s\,\bar\sigma_\Rgh(s)}{\CR}.
\end{equation}
\end{theorem}

In Eq.~\eqref{eq:dataset}, $\Rgh(Y)$ is how rough the outputs must become to fit the labels, $\Rgh(\hat Y(t))$ is how rough they are now, and $\bar\sigma_\Rgh(s)$ is the most roughness that a unit parameter step can add to the outputs at any parameters the step can reach.
The reduction in error that a step can achieve therefore decreases with the roughness of the dataset, and the step length needed to reach a given error grows with it.
The theorem makes no assumption on how the gradient features change along the step, because $\bar\sigma_\Rgh(s)$ is taken over the whole set $B_s$.
It is a statement about one step; Theorem~\ref{thm:path} of Part~\ref{part:path} is its form for a whole stretch of training, with the step length replaced by the length of the parameter path.
Its rate form is Proposition~\ref{prop:rates}(b), which is a statement about $\hat Y$ because $Y$ is constant: the outputs acquire roughness at rate at most $\sR$ per unit of parameter motion, whatever the labels are, while the labels demand roughness $\Rgh(Y)$.

The statement is one-sided.
It does not show that a smooth dataset is learned quickly, since a network can also lack gradient features along smooth directions.
A two-sided statement would need a lower bound on the kernel on smooth sampled functions, which is a property of the architecture rather than of the data.

\begin{theorem}[Main statement of Part~\ref{part:onetime}: the error is bounded below by the absolute roughness of the dataset]\label{thm:slow}
Let $(X, Y)$ be any dataset, let $\Rgh$ be any roughness function, and let $t$ be any time during training under any update rule, with output $\hat Y(t)$ and error $E(t) = \hat Y(t) - Y$.
\begin{enumerate}[label=(\alph*)]
\item \emph{Slow part.} For every $\gamma > 0$,
\begin{equation}\label{eq:main}
\|E(t)\| \ \ge\ \|E_\gamma(t)\| \ \ge\ \frac{\Rgh(E(t)) - \sR\,\Upsilon(\gamma)^{1/2}}{\CR} \ \ge\ \frac{\Rgh(Y) - \Rgh(\hat Y(t)) - \sR\,\Upsilon(\gamma)^{1/2}}{\CR}.
\end{equation}
\item \emph{One step.} For every $s \ge 0$, every step of rescaled length at most $s$ from $\theta(t)$ leaves an error $E_{\mathrm{new}}$ with
\begin{equation}\label{eq:mainstep}
\|E_{\mathrm{new}}\| \ \ge\ \frac{\Rgh(E(t)) - |c_p|\,s\,\bar\sigma_\Rgh(s)}{\CR} \ \ge\ \frac{\Rgh(Y) - \Rgh(\hat Y(t)) - |c_p|\,s\,\bar\sigma_\Rgh(s)}{\CR}.
\end{equation}
\item \emph{Best roughness function.} Let $E(t) \neq 0$.
The supremum over all roughness functions of the middle term of Eq.~\eqref{eq:main} is $\|E_\gamma(t)\|$, attained by $\Rgh(u) = |\ip{\varphi_\gamma}{u}|$ with $\varphi_\gamma = E_\gamma(t)/\|E_\gamma(t)\|$.
With $\Psi = \Psi(t)$ and $F(s) > 0$, the supremum over all roughness functions of the step bound $b_\Rgh(s) = \big(\Rgh(E(t)) - |c_p|\,s\,\sR\big)/\CR$, which is the middle term of Eq.~\eqref{eq:mainstep} with the peak feature roughness at time $t$ in place of $\bar\sigma_\Rgh(s)$, is $F(s)$, Eq.~\eqref{eq:Fs} in Section~\ref{sec:step}; and every step of length at most $s$ leaves $\|E_{\mathrm{new}}\|\ge F(s) - \tfrac12|c_p|\bar H_2(s)s^2$, with a step that leaves at most $F(s) + \tfrac12|c_p|\bar H_2(s)s^2$ (Proposition~\ref{prop:taylorstep}).
The supremum is attained by the Krylov functional $\varphi^\star$ of Theorem~\ref{thm:stepsup} in Section~\ref{sec:step}, which Theorem~\ref{thm:stepkrylov} computes from $Y$, $\hat Y(t)$ and finitely many applications of $K$.
For these two functionals, the last term of each chain is also exact in its signed form, in which $\Rgh(Y) - \Rgh(\hat Y(t))$ is replaced by $\ip{\varphi}{\hat Y(t)} - \ip{\varphi}{Y} = \ip{\varphi}{E(t)}$, that is, by the chain written for $-E(t) = Y - \hat Y(t)$ with the functional $\ip{-\varphi}{\cdot}$, using $\Rgh(-E) = \Rgh(E)$:
\begin{equation}\label{eq:mainexact}
\begin{aligned}
\|E_\gamma(t)\| &= \ip{\varphi_\gamma}{\hat Y(t)} - \ip{\varphi_\gamma}{Y} - \|\varphi_\gamma\|_K\,\Upsilon(\gamma)^{1/2},\\
F(s) &= \ip{\varphi^\star}{\hat Y(t)} - \ip{\varphi^\star}{Y} - |c_p|\,s\,\|\varphi^\star\|_K .
\end{aligned}
\end{equation}
\end{enumerate}
\end{theorem}

In words: at every time, the error is at least the absolute roughness of the dataset, minus the roughness of what the network currently supplies, minus what the kernel removes quickly, in (a), or what one step can remove, in (b), all in the units of the labels and divided by the norm constant.
The dataset enters through $\Rgh(Y)$ alone, and for a local complexity function $\Rgh(Y)$ is fixed before training.
The network enters through $\sR$, the eigenvalues of $K$, and the roughness of its current output.
Part (a) is Theorem~\ref{thm:floor}: a rougher dataset raises the floor, the part of the error on the slow eigendirections.
Part (b) is Theorem~\ref{thm:dataset}: the reduction in error that one step can achieve shrinks as the labels get rougher.
Part (c) says that the middle terms are not loose in $\Rgh$: the best roughness function makes (a) an equality, and in (b) it makes the first-order step bound an equality, which is the error actually reached up to the second-order remainder of Proposition~\ref{prop:taylorstep}.
For a general roughness function, the last term of each chain, which separates the roughness of the dataset from the roughness of the output, is an equality when the output is a multiple $\alpha Y$ of the labels with $0 \le \alpha \le 1$, and can lose otherwise, for instance when the output is rougher than the labels.
For the best roughness functions in their signed form, Eq.~\eqref{eq:mainexact}, this term is exact as well, so the chain from $\|E_\gamma(t)\|$, or from $F(s)$, on collapses to an equality, while $\|E_{\mathrm{new}}\|$ for the best step is within $\tfrac12|c_p|\bar H_2(s)s^2$ of $F(s)$ (Proposition~\ref{prop:taylorstep}); the leak term of (a) is then exactly $\ip{\varphi_\gamma}{E(t) - E_\gamma(t)}$, the component of the removed part of the error along $\varphi_\gamma$.
What remains is the first inequality of (a), $\|E(t)\| \ge \|E_\gamma(t)\|$, which does not involve $\Rgh$ and is an equality only when the error lies entirely outside the range of $K$.
In practice the bounds are computed by the Krylov series of Section~\ref{sec:step}, which uses the signed form, so the only loss in computing $F(s)$ is that of truncating the series, $F_M(s) \le F(s)$; relative to the error actually reached after a step, the remainder of Proposition~\ref{prop:taylorstep} remains.

\subsection{The dataset in the coordinates of the kernel}\label{sec:kernelcoords}

The bound of Theorem~\ref{thm:fast}(b) for the error, $\big(\Rgh(E) - \sR\Upsilon(\gamma)^{1/2}\big)/\CR$, is the step bound $b_\Rgh(s_\gamma)$ of Eq.~\eqref{eq:bR} in Section~\ref{sec:step} for the step of length $s_\gamma = \Upsilon(\gamma)^{1/2}/|c_p|$, and $F(s_\gamma) = \|E_\gamma\|$.
Its supremum over all roughness functions is therefore $\|E_\gamma\|$ itself, attained by the functional along $E_\gamma$ (Proposition~\ref{prop:floorstep} in Appendix~\ref{app:aux}); this is Theorem~\ref{thm:slow}(c).
The Krylov functional of Section~\ref{sec:step} measures the roughness of the dataset in the coordinates of the kernel.
Since $u \mapsto u_\gamma$ is linear, $E_\gamma = \hat Y_\gamma - Y_\gamma$, and so
\begin{equation}\label{eq:Yslow}
\|E_\gamma\| \ \ge\ \|Y_\gamma\| - \|\hat Y_\gamma\| .
\end{equation}
The slow part of the error is at least the slow part of the labels minus the slow part of the current output.
This is statement (1) in the coordinates of the kernel, and by Proposition~\ref{prop:floorstep} the left side is $F(s_\gamma)$, the smallest error of the first-order expansion of a step of length $s_\gamma$.
The slow part of the labels, $\|Y_\gamma\|$, is the roughness of the dataset as the kernel sees it.
It dominates every other measure up to the leak term: by Theorem~\ref{thm:fast}(b), $\|Y_\gamma\| \ge \big(\Rgh(Y) - \sR\,\Upsilon_Y(\gamma)^{1/2}\big)/\CR$ for every roughness function $\Rgh$.
The Krylov values $F_M(s)$ of Eq.~\eqref{eq:FMstep} in Section~\ref{sec:step} compute all of this from the three objects alone, $Y$, $\hat Y$ and finitely many applications of $K$, without an eigendecomposition and without choosing a seminorm.
This is how the one-time theory closes: the dataset $Y$ sets the roughness, the output $\hat Y(t)$ says how much of it has been supplied, the kernel $K(t)$ says how much the first-order term of a step can supply, and the Krylov series turns the three into the exact smallest error of that first-order expansion, which Proposition~\ref{prop:taylorstep} ties to the error actually reached.

\paragraph{Summary of Part~\ref{part:onetime}.}
For every roughness function, at every time during training and under every update rule,
\begin{equation}\label{eq:summary}
\|E\| \ \ge\ \frac{\Rgh(Y) - \Rgh(\hat Y) - \sR\,\Upsilon(\gamma)^{1/2}}{\CR}
\qquad\text{for every } \gamma > 0,
\end{equation}
where $\Rgh(Y)$ is the absolute roughness of the dataset in the units of the labels.
The best roughness function turns $\big(\Rgh(E) - \sR\Upsilon(\gamma)^{1/2}\big)/\CR$ into $\|E_\gamma\| = F(s_\gamma)$, and in its signed form, Eq.~\eqref{eq:mainexact}, the right side as well; $F(s_\gamma)$ is the smallest error of the first-order expansion of a step of length $s_\gamma$, and the Krylov series of Section~\ref{sec:step} computes it from $Y$, $\hat Y$ and the kernel.

\paragraph{What Part~\ref{part:onetime} does not say.}
Every statement above uses the gradient features and the kernel at one time, or, in Theorems~\ref{thm:dataset} and~\ref{thm:slow}(b), the peak feature roughness over the set that one step can reach.
The results describe where the error lies in the eigenbasis of the kernel, the rate of change of the error, and what one step of a given length can do.
They do not describe the error after a stretch of training, because the kernel moves along the path and the path is not known in advance.
Equivalently, the floor holds at time $t$ with the output $\hat Y(t)$ and the kernel $K(t)$ of that time, and neither is known at time 0.
Part~\ref{part:path} integrates the theory along the path, and states how far that integration can go without approximation.

% ===============================================================
\part{The error along training}\label{part:path}
% ===============================================================

The results of Part~\ref{part:onetime} hold at one time during training, with the kernel at that time, or for one training step.
This part sets up the framework in which the gap to a stretch of training is closed as far as it can be closed without approximation.
Section~\ref{sec:state} states the closed system that training is: one state variable, the parameters, and a tower of derivative tensors that is the Taylor series of the network at the start of the path, kept whole.
Section~\ref{sec:identities} proves the identities that hold exactly along every training path: the energy identity, the speed identity, which rewrites the energy identity in terms of the curvature of the loss along the motion, a change of variables and not an integration, and the time derivatives of the error to all orders.
Section~\ref{sec:onetime} places the one-time theory of Part~\ref{part:onetime} along the path.
Section~\ref{sec:path} is the path calculus: the length of the parameter path as the coordinate for the motion of the kernel, the series that bounds every kernel quantity along the path by the derivatives of the network at the start, and the lower and two-sided bounds on the error that follow; it is here that the roughness bounds of Part~\ref{part:onetime} become bounds on the error after a stretch of training.
Section~\ref{sec:allorders} carries the expansions to all orders: the Dyson series in the motion of the kernel, which is exact for every time, and the Picard scheme from the initial state, whose every iterate is computable from the tower at the start and whose limit is the true path; it also states why no single series from the initial state covers a training run.
Section~\ref{sec:predict} is the central result of this part: it states the prediction of the error from the initial state for every update rule that is a smooth flow or map, with a certified two-sided bound, and the principle for comparing designs that it gives.
Section~\ref{sec:loss} reads the loss curve through the theory: the kernel term of the curvature along the gradient is set by where the error lies in the spectrum, and so by the slow part of Part~\ref{part:onetime}, and the rest is one scalar-valued function of time through which feature learning enters the loss curve at each instant.
Section~\ref{sec:closing} lists the assumptions that would close the system, each with its consequence proved, and Section~\ref{sec:measure} what to measure.

Nothing in this part is approximated.
Every remainder is an exact integral or an infinite series with all terms, and every bound on a series is a bound on its whole tail.

\paragraph{Conventions.}
Throughout Part~\ref{part:path}, $\Pi = I$ and $c_p = 1$, so that $\hat Y = \Net$ at the training inputs and $E = \hat Y - Y$; the loss is $L = \frac12\|E\|^2$, with gradient $g_E = \nabla_\theta L = \frac1n\Psi^\top E$.
A fixed step-size matrix $\Pi$ is absorbed by the change of variables $\theta\mapsto\Pi^{-1/2}\theta$, which contracts each derivative with $\Pi^{1/2}$ in every slot; an output scale $c_p\neq1$ is absorbed as in Proposition~\ref{prop:scale}, with the loss of Part~\ref{part:onetime}, $\frac12\|\Net - Y/c_p\|^2 = L/c_p^2$, so that gradient-flow time is rescaled by $c_p^2$.
Both are omitted to keep the formulas short.
All other objects are those of Part~\ref{part:onetime}: the kernel $K$ with eigenvalues $\lambda_1\ge\lambda_2\ge\dots\ge0$ and orthonormal eigenvectors $v_k$, the components $\hat e_k = \ip{v_k}{E}$, the gradient $g_\varphi = \frac1n\Psi^\top\varphi$ and kernel size $\|\varphi\|_K = |g_\varphi|$ of a sampled function $\varphi$, the slow part $E_\gamma$ and $\Upsilon(\gamma)$, and roughness functions $\Rgh$ with norm constant $\CR$ and peak feature roughness $\sR$.

% ===============================================================
\section{State, closure, and the Taylor tower}\label{sec:state}
% ===============================================================

\paragraph{The state.}
The only state variable of training is $\theta$.
Under gradient flow,
\begin{equation}\label{eq:gf}
\dot\theta = -g_E = -\nabla_\theta L(\theta),
\end{equation}
and under gradient descent $\theta_{j+1} = \theta_j - \eta\,g_{E_j}$, whose iterates are joined by the piecewise-linear path.
Every other object of the theory, the error, the kernel, its eigenvalues, the Hessians $S(\varphi)$ defined below, is a function of $\theta$ alone, evaluated along the path.

\paragraph{The tower.}
For a sampled function $\varphi$ let $f_\varphi(\theta) = \ip{\varphi}{\Net}$, a scalar function of the parameters, and for $m\ge0$ let
\begin{equation}\label{eq:Tm}
T_m(\varphi;\theta) = \nabla^m_\theta f_\varphi(\theta),
\end{equation}
a symmetric $m$-linear form on parameter space; $T_0 = f_\varphi$, $T_1 = g_\varphi$, and $T_2 = S(\varphi) = \frac1n\sum_i\varphi_i\cdot\nabla^2_\theta\Net(x_i)$, the Hessian of $\ip{\varphi}{\Net}$ in $\theta$.
Their sizes at $\theta$ are
\begin{equation}\label{eq:Dm}
D_m(\varphi;\theta) = \|T_m(\varphi;\theta)\|,\qquad
H_m(\theta) = \Big(\tfrac1n\sum_i\|\nabla^m_\theta\Net(x_i;\theta)\|^2\Big)^{1/2},
\end{equation}
with $\|T\| = \sup|T[u_1,\dots,u_m]|$ over unit vectors $u_1,\dots,u_m$ the operator norm of an $m$-linear form, written also $\|\cdot\|_{op}$, and for vector outputs $\|\nabla^m_\theta\Net(x_i)\|$ the largest operator norm of $\phi\cdot\nabla^m_\theta\Net(x_i)$ over unit output vectors $\phi$; by the Cauchy--Schwarz inequality $D_m(\varphi;\theta)\le\|\varphi\|H_m(\theta)$.
$D_1(\varphi;\theta) = \|\varphi\|_K$, $H_1^2$ is at most the trace of $K$, with equality for scalar outputs, and $H_2$ is the second-derivative size of the network.

\begin{proposition}[The loss Hessian and the kernel]\label{prop:hessian}
If $\Net$ is twice differentiable in $\theta$ at the training inputs, the Hessian of the loss is
\begin{equation}\label{eq:hessian}
\Hess = \nabla^2_\theta L = G + S(E),\qquad G = \tfrac1n\Psi^\top\Psi,
\end{equation}
and $G$ has the same nonzero eigenvalues as $K$. Hence
\begin{equation}\label{eq:sharpness}
\big|\lambda_{\max}(\Hess) - \lambda_1\big|\ \le\ \|S(E)\|_{op}\ \le\ H_2\,\|E\| .
\end{equation}
\end{proposition}

\begin{proof}
$L$ is quadratic in the outputs, so differentiating it twice gives Eq.~\eqref{eq:hessian} with no remainder; $\frac1n\Psi^\top\Psi$ and $\frac1n\Psi\Psi^\top$ share their nonzero eigenvalues; the first bound is Weyl's inequality, and the second is $\|S(E)\|_{op}\le\frac1n\sum_i|E_i|\,\|\nabla^2_\theta\Net(x_i)\|$ with the Cauchy--Schwarz inequality.
\end{proof}

The largest eigenvalue of the loss Hessian, the sharpness, is the top eigenvalue of the training kernel up to a term proportional to the error.
The kernel term $G$ is the kernel read in parameter space; $S(E)$ is the part through which the second derivative of the network enters.

\begin{proposition}[The tower moves by contraction with the gradient]\label{prop:tower}
Let $\Net$ be $m + 1$ times continuously differentiable in $\theta$.
Along gradient flow, for every sampled function $\varphi$,
\begin{equation}\label{eq:tower}
\frac{d}{dt}\,T_m(\varphi;\theta(t)) = -\,T_{m+1}(\varphi;\theta(t))\big[g_E,\,\cdot\,,\dots,\,\cdot\,\big],
\end{equation}
the next tensor with $g_E$ in one slot.
In particular:
\begin{enumerate}[label=(\alph*)]
\item ($m = 0$) $\frac{d}{dt}\ip{\varphi}{\Net} = -g_\varphi\cdot g_E = -\ip{\varphi}{KE}$, that is, $\dot E = -KE$, Eq.~\eqref{eq:gf} read in the output;
\item ($m = 1$) $\dot g_\varphi = -S(\varphi)g_E$, and for the kernel, $\frac{d}{dt}\ip{\varphi}{K\psi} = -g_\psi\cdot S(\varphi)g_E - g_\varphi\cdot S(\psi)g_E$;
\item ($m = 2$) $\dot S(\varphi) = -T_3(\varphi)[g_E,\cdot,\cdot]$.
\end{enumerate}
Along the segment of one step of gradient descent, Eq.~\eqref{eq:tower} holds with $\frac{d}{d\tau}$ in place of $\frac{d}{dt}$ and the constant direction $\eta g_{E_j}$ in place of $g_E$; the forms in (a) and (b) in terms of $K$ and $E$ do not, since $g_{E_j}$ is not $g_E$ at the points of the segment. Integrating over $\tau\in[0,1]$ gives the one-step identities of Section~\ref{sec:step} with their remainders.
\end{proposition}

\begin{proof}
The chain rule, $\frac{d}{dt}T_m(\varphi;\theta(t)) = T_{m+1}(\varphi;\theta(t))[\dot\theta,\cdot,\dots,\cdot]$, with Eq.~\eqref{eq:gf}.
In (a), $\ip{\varphi}{KE} = g_\varphi\cdot g_E$ by the definition of $K$; in (b), $\ip{\varphi}{K\psi} = g_\varphi\cdot g_\psi$ and the product rule.
\end{proof}

The tower is infinite, and each level is driven by the next.
No finite part of it is closed: the error needs the kernel, the kernel needs the Hessians $S$, and so on.
This is the whole difficulty, stated once.
The literature closes the tower by truncation: at $m = 1$, holding $K$ fixed, which is the linearized or fixed-kernel picture; at $m = 2$, holding $S$ fixed, which is the quadratic model.
This paper does neither.
It closes the tower by the Taylor series of the network, with every term kept.

\paragraph{The top eigenvalue.}
Part (b) applied to $\varphi = \psi = v_1$, with $v_1$ frozen at the current time, gives the rate of the top eigenvalue of the kernel, by the envelope (Hellmann--Feynman) argument in the proof below.

\begin{proposition}[Rate of the top eigenvalue]\label{prop:lambda1}
Let $\Net$ be twice continuously differentiable in $\theta$ and let $\lambda_1$ be simple.
Then $\nabla_\theta\lambda_1 = 2a_1$ with $a_1 = S(v_1)g_{v_1}$.
Let $q = \Psi a_1$, the output change produced by the parameter direction $a_1$.
Under gradient flow,
\begin{equation}\label{eq:lambda1dot}
\dot\lambda_1 = -2\ip{q}{E},
\end{equation}
and if $\lambda_1$ is simple on $[0,t]$, then $\lambda_1(t) - \lambda_1(0) = -2\int_0^t\ip{q}{E}\,ds$ and $\lambda_1(t) - \lambda_1(0)\le2\int_0^t\|q\|\,\|E\|\,ds$.
\end{proposition}

\begin{proof}
$\lambda_1 = \max_{\|v\|=1}|g_v|^2 = |g_{v_1}|^2$, and for a simple eigenvalue $\lambda_1$ is continuously differentiable with $\nabla_\theta\lambda_1$ equal to the gradient of $|g_v|^2$ at the fixed maximizer $v = v_1$ (the envelope theorem for a simple eigenvalue), so $\nabla_\theta\lambda_1 = 2(\nabla_\theta g_{v_1})^\top g_{v_1} = 2S(v_1)g_{v_1}$.
Under gradient flow $\dot\theta = -g_E$, and $2a_1\cdot g_E = 2\ip{\Psi a_1}{E}$.
The integrated form is the fundamental theorem of calculus, and the bound is the Cauchy--Schwarz inequality.
\end{proof}

The rate of growth of the top eigenvalue is a linear functional of the error.
When $q\neq0$, the rank-one seminorm $u\mapsto|\ip{q}{u}|$ is a roughness function, with norm constant $\|q\|$ and peak feature roughness $\|q\|_K$, so Theorem~\ref{thm:fast} applies to it and splits the drive by where the error lies: $|\ip{q}{E - E_\gamma}|\le\|q\|_K\Upsilon(\gamma)^{1/2}$, and whatever exceeds this is driven by the slow part.
To raise $\lambda_1$, training must carry error.
The theory bounds the drive and does not give its sign; the application to the edge of stability in Part~\ref{part:apps} uses this proposition.

\paragraph{The derivatives of the loss.}
The loss is quadratic in the outputs, so every one of its derivatives splits into a part linear in the error and a part that does not involve the error.

\begin{proposition}[Derivatives of the loss to all orders]\label{prop:lossderiv}
Let $\Net$ be $k$ times continuously differentiable in $\theta$.
Then
\begin{equation}\label{eq:lossderiv}
\begin{aligned}
\nabla^k_\theta L &= S_k(E) + G_k,\qquad
S_k(E) = \tfrac1n\sum_iE_i\cdot\nabla^k_\theta\Net(x_i),\\
G_k[u_1,\dots,u_k] &= \tfrac12\sum_{\substack{A\subseteq\{1,\dots,k\}\\ 1\le|A|\le k-1}}\tfrac1n\sum_i\nabla^{|A|}_\theta\Net(x_i)[u_A]\cdot\nabla^{k-|A|}_\theta\Net(x_i)[u_{A^c}],
\end{aligned}
\end{equation}
where $u_A$ denotes the arguments with indices in $A$.
In particular $\nabla L = S_1(E) = g_E$, $G_1 = 0$, $S_2 = S$, $G_2 = G$, and
\begin{equation}\label{eq:lossderivbound}
\|S_k(E)\|\le H_k\,\|E\|,\qquad \|G_k\|\le\tfrac12\sum_{j=1}^{k-1}\binom kj H_jH_{k-j} .
\end{equation}
\end{proposition}

\begin{proof}
$L = \frac1{2n}\sum_if_i\cdot f_i$ with $f_i = \Net(x_i) - Y_i$, and $\nabla^jf_i = \nabla^j\Net(x_i)$ for $j\ge1$.
The Leibniz rule for the $k$-th derivative of $\frac12f_i\cdot f_i$ is $\frac12\sum_{A}\nabla^{|A|}f_i[u_A]\cdot\nabla^{k-|A|}f_i[u_{A^c}]$ over all subsets $A$; the two terms with $A = \emptyset$ and $A = \{1,\dots,k\}$ are equal and give $f_i\cdot\nabla^kf_i$, which summed over $i$ is $S_k(E)$, and the rest is $G_k$.
The bounds are the Cauchy--Schwarz inequality on the sum over $i$, with the operator norm of a contraction at most the product of the operator norms, and $\binom kj$ subsets of size $j$.
\end{proof}

Proposition~\ref{prop:lossderiv} is Proposition~\ref{prop:hessian} carried to every order.
For a network every $S_k(E)$ is a contraction of the error with the $k$-th derivative of the network, and every $G_k$ is built from lower derivatives and carries no error.

\paragraph{The Taylor closure.}
Let $\theta_0$ be the start of a stretch of training and $\Delta(t) = \theta(t) - \theta_0 = -\int_0^tg_E\,ds$ the displacement.

\begin{proposition}[Taylor closure]\label{prop:closure}
Let $t$ be such that for every training input $\theta\mapsto\Net(x_i;\theta)$ extends holomorphically to the complex ball $|\theta - \theta_0| < R_0$ in $\mathbb{C}^{d_\theta}$, with $R_0 > \sup_{s\le t}|\Delta(s)|$.
Then for every $s\le t$, every sampled function $\varphi$ and every $m\ge0$,
\begin{equation}\label{eq:closure}
T_m(\varphi;\theta(s)) = \sum_{k\ge0}\frac{1}{k!}\,T_{m+k}(\varphi;\theta_0)\big[\Delta(s),\dots,\Delta(s),\,\cdot\,,\dots,\,\cdot\,\big],
\end{equation}
with $k$ copies of $\Delta(s)$ in the first slots.
In particular the error and the gradient are
\begin{equation}\label{eq:closureE}
E(s) = E_0 + \Psi_0\Delta(s) + \sum_{k\ge2}\frac{1}{k!}\,\nabla^k_\theta\Net(\theta_0)\big[\Delta(s)^{\otimes k}\big],\qquad
g_\varphi(\theta(s)) = \sum_{k\ge0}\frac{1}{k!}\,T_{k+1}(\varphi;\theta_0)\big[\Delta(s)^{\otimes k},\,\cdot\,\big],
\end{equation}
with $g_E(\theta(s)) = \sum_{k\ge0}\frac{1}{k!}T_{k+1}(E(s);\theta_0)[\Delta(s)^{\otimes k},\cdot]$ by linearity of $T_{k+1}$ in $\varphi$, and training is the closed system $\dot\Delta = -g_E(\theta_0 + \Delta)$ with $E$ and $g_E$ given by these series.
\end{proposition}

\begin{proof}
$f_\varphi$ is a finite linear combination of the $\Net(x_i;\cdot)$, so it is holomorphic on the complex ball, and so is every derivative $\nabla^mf_\varphi$.
A holomorphic function on a ball equals its Taylor series at the centre on the whole ball, and Eq.~\eqref{eq:closure} is that series for $\nabla^mf_\varphi$, evaluated at $\theta(s)$, which lies in the ball.
\end{proof}

Eq.~\eqref{eq:closureE} is exact and contains no kernel: the first two terms are the first-order Taylor expansion of the error about $\theta_0$, and the series is every higher order, kept whole.
The state of the whole tower at time $s$ is the state at $\theta_0$ evaluated at $\Delta(s)$, so what has to be controlled along a path is $\Delta$, and the quantity that controls it is the length of the path (Section~\ref{sec:path}).
Where analyticity or convergence is not available, every statement below has a form with an exact integral remainder instead of the series; both are stated.

% ===============================================================
\section{Exact identities along training}\label{sec:identities}
% ===============================================================

The identities of this section hold along every gradient-flow path, with the kernel moving as it does, and need no closure.
They are the statements about $E(t)$, $\dot E(t)$, $\ddot E(t)$, and the time derivatives of $E$ to every order, that the theory has.

\begin{theorem}[Energy, speed, and the loss curve]\label{thm:identities}
Let $\Net$ be twice continuously differentiable in $\theta$, and let $\theta(t)$ be the gradient-flow path with $g_E(0)\ne0$.
Write $\Hess = \nabla^2_\theta L = G + S(E)$ for the loss Hessian, $G = \frac1n\Psi^\top\Psi$ (Proposition~\ref{prop:hessian}).
\begin{enumerate}[label=(\alph*)]
\item \emph{Energy.} $\dot E = -KE$, $\frac{d}{dt}\frac12\|E\|^2 = -\|E\|_K^2 = -|\dot\theta|^2$, and
\begin{equation}\label{eq:energy}
\|E(t)\|^2 = \|E_0\|^2 - 2\int_0^t\|E(s)\|_{K(s)}^2\,ds = \|E_0\|^2 - 2\int_0^t|\dot\theta(s)|^2\,ds .
\end{equation}
\item \emph{Speed.} $\dot g_E = -\Hess\,g_E$. Hence $g_E(t)\ne0$ for all $t$, and with
\begin{equation}\label{eq:kappa}
\kappa(t) = \frac{g_E^\top\Hess\,g_E}{|g_E|^2},
\end{equation}
the curvature of the loss along the gradient,
\begin{equation}\label{eq:speed}
\|E(t)\|_{K(t)}^2 = |g_E(t)|^2 = \|E_0\|_{K_0}^2\exp\Big(-2\int_0^t\kappa(s)\,ds\Big).
\end{equation}
\item \emph{The loss curve.}
\begin{equation}\label{eq:losscurve}
\|E(t)\|^2 = \|E_0\|^2 - 2\,\|E_0\|_{K_0}^2\int_0^t\exp\Big(-2\int_0^s\kappa(\tau)\,d\tau\Big)ds ,
\qquad
\frac{d^2}{dt^2}\,\tfrac12\|E\|^2 = 2\,\kappa\,\|E\|_K^2 .
\end{equation}
\item \emph{The split of the curvature.} $\kappa = \kappa_G + \kappa_S$ with
\begin{equation}\label{eq:kappasplit}
\begin{aligned}
\kappa_G &= \frac{g_E^\top Gg_E}{|g_E|^2} = \frac{\|KE\|^2}{\|E\|_K^2} = \frac{\sum_k\lambda_k^2\hat e_k^2}{\sum_k\lambda_k\hat e_k^2},\\
\kappa_S &= \frac{g_E^\top S(E)g_E}{|g_E|^2} = \frac{\ip{E}{\nabla^2_\theta\Net[g_E,g_E]}}{\|E\|_K^2},
\qquad |\kappa_S|\le H_2(\theta)\,\|E\| .
\end{aligned}
\end{equation}
$\kappa_G$ is the mean of the eigenvalues of $K$ under the weights $\lambda_k\hat e_k^2/\|E\|_K^2$, so it lies between the smallest and the largest eigenvalue on which the error has a component, and $\kappa_G\le\lambda_1$.
\item \emph{Acceleration of the error.}
\begin{equation}\label{eq:accel}
\ddot E = K^2E + \nabla^2_\theta\Net[g_E,g_E] + \Psi\,S(E)\,g_E ,
\end{equation}
where $\nabla^2_\theta\Net[g_E,g_E]$ is the sampled function with values $\nabla^2_\theta\Net(x_i)[g_E,g_E]$.
\end{enumerate}
\end{theorem}

\begin{proof}
(a) Proposition~\ref{prop:tower}(a); $\ip{E}{KE} = |g_E|^2 = |\dot\theta|^2$; integrate.
(b) $g_E = \nabla_\theta L$ is continuously differentiable, so $\dot g_E = \Hess\dot\theta = -\Hess g_E$, and $\frac{d}{dt}|g_E|^2 = -2g_E^\top\Hess g_E = -2\kappa|g_E|^2$ wherever $g_E\ne0$.
Since $|g_E^\top\Hess g_E|\le\|\Hess\|_{op}|g_E|^2$ everywhere and $\|\Hess\|_{op}\le M$ on the compact path segment, $\frac{d}{dt}|g_E|^2\ge-2M|g_E|^2$, so by Gr\"onwall's inequality $|g_E(t)|^2\ge|g_E(0)|^2e^{-2Mt} > 0$; then $\kappa$ is defined and bounded, and integrating $\frac{d}{dt}\log|g_E|^2 = -2\kappa$ gives Eq.~\eqref{eq:speed}.
(c) Insert Eq.~\eqref{eq:speed} into Eq.~\eqref{eq:energy}; differentiate Eq.~\eqref{eq:energy} twice.
(d) $g_E^\top Gg_E = \|\Psi g_E\|^2 = \|KE\|^2$ and $|g_E|^2 = \|E\|_K^2 = \sum_k\lambda_k\hat e_k^2$, $\|KE\|^2 = \sum_k\lambda_k^2\hat e_k^2$; $g_E^\top S(E)g_E = \frac1n\sum_iE_i\cdot\nabla^2_\theta\Net(x_i)[g_E,g_E]$ by the definition of $S$; the bound is $\|S(E)\|_{op}\le\|E\|H_2$ from Proposition~\ref{prop:hessian}.
(e) Differentiate $\dot E = -KE$: $\ddot E = -\dot KE - K\dot E = -\dot KE + K^2E$, and by Proposition~\ref{prop:tower}(b), $\ip{\varphi}{\dot KE} = -g_E\cdot S(\varphi)g_E - g_\varphi\cdot S(E)g_E = -\ip{\varphi}{\nabla^2_\theta\Net[g_E,g_E]} - \ip{\varphi}{\Psi S(E)g_E}$ for every $\varphi$.
\end{proof}

\paragraph{The theorem as kinematics.}
Theorem~\ref{thm:identities} is the kinematics of the path $\theta(t)$ and nothing more.
Since $\dot\theta = -\nabla_\theta L$, differentiating once more along the path gives $\ddot\theta = -\Hess\,\dot\theta$: the parameter acceleration is the loss Hessian applied to the velocity.
Hence $\frac{d}{dt}|\dot\theta|^2 = -2\,\Hess[\dot\theta,\dot\theta]$, and
\begin{equation}\label{eq:kappavel}
\kappa = \frac{\Hess[\dot\theta,\dot\theta]}{|\dot\theta|^2}
\end{equation}
is the second directional derivative of the loss in the direction in which the path is moving at time $t$, the Hessian contracted twice with the current velocity and normalized by the speed; $g_E = -\dot\theta$ is the same vector with the opposite sign.
Part (b) says that gradient flow decelerates at the rate of the curvature it is running into, and accelerates where that curvature is negative.
Part (a) says that the loss falls at the squared speed, $\dot L = \nabla_\theta L\cdot\dot\theta = -|\dot\theta|^2$, and part (c) is (a) with the speed of (b) inserted.
None of this uses the network: the loss forms $\dot L = -|\dot\theta|^2$, part (b), and $L(t) = L(0) - |g_E(0)|^2\int_0^te^{-2\int_0^s\kappa}\,ds$ hold for gradient flow on any twice differentiable function.
Part (d) is where the network enters: for $L = \frac12\|E\|^2$ the two pieces of the Hessian give two contractions with the velocity,
\begin{equation}\label{eq:kappavelsplit}
\kappa_G = \frac{\|\Psi\dot\theta\|^2}{|\dot\theta|^2} = \frac{\|\dot E\|^2}{|\dot\theta|^2},\qquad
\kappa_S = \frac{\ip{E}{\nabla^2_\theta\Net[\dot\theta,\dot\theta]}}{|\dot\theta|^2}:
\end{equation}
$\kappa_G$ is the squared speed of the output over the squared speed of the parameters, how much output motion the current parameter motion buys, and $\kappa_S$ is the second derivative of the output map along the direction of motion, projected on the error.
Part (e) is the same statement one level down: since $\dot E = \Psi\dot\theta$, $\ddot E = \nabla^2_\theta\Net[\dot\theta,\dot\theta] + \Psi\ddot\theta$, the curvature of the output map along the motion plus the features applied to the parameter acceleration, and substituting $\ddot\theta = -\Hess\dot\theta$ gives the three terms of Eq.~\eqref{eq:accel}.

Parts (b) and (c) are not an integration of $\dot E$, and they add no information to part (a).
Any positive function can be written as $f(t) = f(0)\exp(-2\int_0^t\kappa)$ with $\kappa = -\frac12\frac{d}{dt}\log f$; applied to the speed this is part (b), and inserted into part (a) it is part (c).
The whole path is inside $\kappa(t)$ exactly as it was inside $|\dot\theta(t)|^2$, and integrating $\kappa$ through time is the same problem as integrating the speed, which is a scalar summary of $\dot E$.
What the theorem adds to the energy identity is one identification: the log-derivative of the loss rate equals $\Hess[\dot\theta,\dot\theta]/|\dot\theta|^2$, a local geometric quantity that can be computed at an instant from one Hessian-vector product, without knowledge of the curve.
That identification is a change of variables with a measurable consequence, not progress on the integration; the progress on integration in this part is in Sections~\ref{sec:path} and~\ref{sec:allorders}, and its cost grows with the length of the window.
Part (d) separates, within that identification, what the kernel at time $t$ determines, $\kappa_G$, from what feature learning adds, $\kappa_S$.
Part (e) says the same for the error itself: the first term, $K^2E$, involves only the kernel at time $t$, the second is the curvature of the output along the gradient, which is the same sampled function whose component along $E$ is the numerator of $\kappa_S$, and the third is the change of the kernel acting on the error.
The loss is convex in $t$ exactly where $\kappa\ge0$, and $\kappa_S$ is the only part of $\kappa$ that can be negative.

\paragraph{Why this function.}
$\kappa(t)$, a scalar-valued function of time, is not an abbreviation chosen for convenience; at each instant it is the only part of the Hessian that the loss curve can see, and as a function it determines the loss curve for any objective.
Gradient flow is a curve with tangent $-g_E$.
The Hessian enters its dynamics only through $\Hess g_E$, and
\begin{equation}\label{eq:Hsplit}
\Hess g_E = \kappa\,g_E + (\Hess g_E)_\perp,\qquad (\Hess g_E)_\perp\perp g_E :
\end{equation}
the first part changes the speed $|g_E|$, the second turns the path and leaves the speed alone.
Since $\dot L = -|g_E|^2$, the rate of the loss is the squared speed, and so
\begin{equation}\label{eq:Lddot}
\ddot L = -2\kappa\,\dot L .
\end{equation}
This is a first-order linear equation for the loss rate whose only coefficient is $\kappa$, and it holds for gradient flow on any twice differentiable objective, not only for the quadratic loss of a network.
This is a routing identity, not a truncation.
The loss curve, Eq.~\eqref{eq:losscurve}, depends on the path only through the one scalar function $\kappa(t)$, and $\kappa$ at an instant depends on the Hessian only through its quadratic form along the gradient; nothing is dropped, and the whole tower is inside $\kappa(t)$.
The tower enters as soon as $\kappa$ is to be predicted rather than read off a run: the rate of $\kappa$, part (d) of the proposition below, brings in the turning part $(\Hess g_E)_\perp$ and the third derivative of the loss, the rate of those brings in the fourth, and the hierarchy Eq.~\eqref{eq:moments} closes at no order.
What is special about the second derivative is only that the loss is quadratic in $E$, so its rate is the squared speed and one derivative of the speed is the whole of $\ddot L$.

\begin{proposition}[What $\kappa$ determines]\label{prop:whykappa}
Let $\Net$ be three times continuously differentiable in $\theta$ (twice suffices for (a)--(c)), and consider a gradient-flow path with $g_E(0)\neq0$, so that $E\neq0$ and $g_E\neq0$ at every time (Theorem~\ref{thm:identities}(b)).
Let $\rho =|g_E|^2/L = 2\|E\|_K^2/\|E\|^2$ be the instantaneous rate of the loss, so that $\dot L = -\rho L$, and let $w_k = \hat e_k^2/\|E\|^2$ be the distribution of the error over the eigenvectors of $K(t)$, with $\mathrm{mean}_w$ and $\mathrm{Var}_w$ the mean and the variance of $\lambda$ under $w$.
\begin{enumerate}[label=(\alph*)]
\item \emph{The active eigenvalue.} $\rho = 2\,\mathrm{mean}_w(\lambda)$ and $\kappa_G = \mathrm{mean}_w(\lambda^2)/\mathrm{mean}_w(\lambda)$, so $\kappa_G\ge\rho/2$, with equality exactly when the error lies on one eigenvalue.
\item \emph{The shape of the loss curve.} Exactly,
\begin{equation}\label{eq:rho}
\dot\rho = \rho\,(\rho - 2\kappa),\qquad
\frac{d^2}{dt^2}\log L = \rho\,(2\kappa - \rho),\qquad
2\kappa - \rho = \frac{2\,\mathrm{Var}_w(\lambda)}{\mathrm{mean}_w(\lambda)} + 2\kappa_S .
\end{equation}
Hence the kernel term alone never makes $\log L$ concave: at every time, with the kernel of that time, $2\kappa_G - \rho\ge0$, with equality exactly when the error lies on one eigenvalue of $K(t)$.
$\log L$ is concave at time $t$, so that the loss is falling faster than exponentially, if and only if
\begin{equation}\label{eq:logconcave}
\kappa_S(t) < -\,\frac{\mathrm{Var}_w(\lambda)}{\mathrm{mean}_w(\lambda)} .
\end{equation}
\item \emph{The bound by the sharpness.} $\lambda_{\min}(\Hess)\le\kappa\le\lambda_{\max}(\Hess)$, with $\kappa = \lambda_{\max}(\Hess)$ exactly when $g_E$ is a top eigenvector of $\Hess$.
\item \emph{The rate of $\kappa$.} Exactly, along the moving kernel,
\begin{equation}\label{eq:kappadot}
\dot\kappa = -\,\frac{2\,|(\Hess g_E)_\perp|^2}{|g_E|^2} - \frac{\nabla^3_\theta L[g_E,g_E,g_E]}{|g_E|^2},
\end{equation}
with $(\Hess g_E)_\perp$ as in Eq.~\eqref{eq:Hsplit}.
The first term is never positive, so $\kappa$ can rise only through the third derivative of the loss along the gradient, $\nabla^3_\theta L[g_E,g_E,g_E] = S_3(E)[g_E,g_E,g_E] + G_3[g_E,g_E,g_E]$ by Proposition~\ref{prop:lossderiv}.
\end{enumerate}
\end{proposition}

\begin{proof}
(a) $\|E\|_K^2 = \sum_k\lambda_k\hat e_k^2$ and $\|KE\|^2 = \sum_k\lambda_k^2\hat e_k^2$ in Eq.~\eqref{eq:kappasplit}; $\mathrm{mean}_w(\lambda^2)\ge\mathrm{mean}_w(\lambda)^2$, with equality if and only if $\lambda$ is constant on the support of $w$.
(b) $\dot\rho = \frac{d}{dt}|g_E|^2/L - |g_E|^2\dot L/L^2 = -2\kappa\rho + \rho^2$ by Theorem~\ref{thm:identities}(b) and $\dot L = -\rho L$; $\frac{d}{dt}\log L = -\rho$; and $2\kappa_G - \rho = 2\big(\mathrm{mean}_w(\lambda^2) - \mathrm{mean}_w(\lambda)^2\big)/\mathrm{mean}_w(\lambda)$ by (a).
(c) The Rayleigh quotient of a symmetric matrix.
(d) With $m_k = \nabla^k_\theta L[g_E,\dots,g_E]$ as in Proposition~\ref{prop:timederiv}(c), whose proof is the chain rule and does not use this proposition, $\kappa = m_2/m_1$, $\dot m_1 = -2m_2$ and $\dot m_2 = -m_3 - 2|\Hess g_E|^2$, so $\dot\kappa = (\dot m_2m_1 - m_2\dot m_1)/m_1^2 = -(m_3 + 2|\Hess g_E|^2)/|g_E|^2 + 2\kappa^2$, and $|\Hess g_E|^2 - \kappa^2|g_E|^2 = |(\Hess g_E)_\perp|^2$ by Eq.~\eqref{eq:Hsplit}.
\end{proof}

Part (a) holds at every time with the kernel of that time, which is not fixed: $\kappa_G(t)$ is the mean eigenvalue of $K(t)$ currently being removed, weighted by how much error each mode is losing; $1/\rho$ is the time scale on which the loss falls at time $t$, and $1/\kappa$ the one on which the speed decays.
By the floor of Part~\ref{part:onetime}, for a rough dataset the error lies on the slow eigenvalues, so the loss rate $\rho = 2\,\mathrm{mean}_w(\lambda)$ is small: by Proposition~\ref{prop:kappaG}, $\rho\le2\gamma + 2\lambda_1^{3/2}\Upsilon(\gamma)^{1/2}/\|E\|$ for every $\gamma > 0$.
$\kappa_G$ itself is a mean under the weights $\lambda_k\hat e_k^2$, which favour the fast modes, so it is small only once the kernel energy $\|E\|_K^2$, and not only the error, lies on the slow eigenvalues.
Part (d) is the exact rate of that time scale along the moving kernel.
Its first term is the turning of the path and is never positive.
Its second term is the only one through which $\kappa$ can rise, and it is the next entry of the tower in which $\kappa_S$ sits: a contraction of the third derivative of the loss with the gradient.
Part (b) says what feature learning has to do to show in the loss curve at all.
The kernel term, set by the kernel of each instant, can only flatten the loss on a logarithmic axis at that instant; a drop that steepens needs $\kappa_S$ negative, and by more than the spread of the current error over the spectrum.
So the open sign of $\kappa_S$ is not one assumption among several: it is the one scalar-valued function of time through which training can do what the kernel of the moment cannot, and Eq.~\eqref{eq:logconcave} is the threshold it has to beat at each instant.
Part (c) ties $\kappa$ to the sharpness: a bound on the sharpness along the path is a bound on $\kappa$, with equality when the gradient lies on the top eigenvector of the Hessian, which is the direction on which the oscillation at the edge of stability concentrates it (Part~\ref{part:apps}).
What $\kappa$ does not determine is the direction $E/\|E\|$, and with it where the error lies; that is the content of the path calculus of Section~\ref{sec:path}, and it is why $\kappa$ is only the first ratio, $m_2/m_1$, of the hierarchy of Proposition~\ref{prop:timederiv}(c).

\paragraph{The tower is not ordered by $\kappa$.}
Nothing above makes the higher levels of the tower small.
In this paragraph and in Proposition~\ref{prop:timederiv}, $n$ in $E^{(n)}$ is the order of a time derivative, not the number of training inputs.
By Proposition~\ref{prop:timederiv}(b), a term of $E^{(n)}$ with $m$ factors of $E$ is bounded by a product of derivative sizes $H_{k_1}\cdots H_{k_p}$ with $k_i\ge2$ and $\sum_i(k_i - 1) = m - 1$, a power of $\lambda_1$ fixed by dimensions once that product is given, and $\|E\|^m$: each application of $D$ that raises the number of factors of $E$ by one raises the order of one tensor by one, Eq.~\eqref{eq:Dgen}, and each factor of $\Psi$ counts as $\lambda_1^{1/2}$.
Against the kernel term $\lambda_1^n\|E\|$ the level $m$ therefore carries $m - 1$ factors of $\|E\|$ weighed against derivative sizes: $H_2\|E\|/\lambda_1$ for $m = 2$, and $H_3\|E\|^2/\lambda_1^{3/2}$ or $H_2^2\|E\|^2/\lambda_1^2$ for $m = 3$.
These are bounds on terms, not sizes.
When the ratios are small the bounds order the levels, with the term $(-K)^nE$ leading and the two terms of level two in Eq.~\eqref{eq:accel}, $\nabla^2_\theta\Net[g_E,g_E]$ and $\Psi S(E)g_E$, both bounded by $H_2\lambda_1\|E\|^2$, as the first corrections; but the number of terms at each level grows with $n$, and no remainder for cutting $E^{(n)}$ at a level is proved here.
The control of a tail that is proved is that of Section~\ref{sec:allorders}, in the variables of that section, with a number of terms that grows with the length of the time window.
Outside that control the levels are comparable and no level can be dropped.
The only exact statements are then the identities of this section, which route the whole tower through $\kappa$ without ordering it, and the full series of Section~\ref{sec:allorders}.

\paragraph{Every order.}
Parts (a) and (e) are the first two time derivatives of the error.
Every further one is obtained by the same rule, and the rule is a single derivation.

\begin{proposition}[Time derivatives to all orders]\label{prop:timederiv}
Let $V(\theta) = -\nabla_\theta L(\theta)$, and for any $\theta$-dependent tensor $Q$ let $DQ = \nabla_\theta Q[V]$, the derivative of $Q$ along the vector field of gradient flow.
$D$ is a derivation: it obeys the Leibniz rule on contractions.
Let $\Net$ be $n + 1$ times continuously differentiable.
\begin{enumerate}[label=(\alph*)]
\item Along gradient flow, $\frac{d}{dt}Q(\theta(t)) = (DQ)(\theta(t))$ for every $Q$, so
\begin{equation}\label{eq:Dn}
\begin{aligned}
E^{(n)}(t) &= (D^n\Net)(\theta(t)),\qquad g_E^{(n)}(t) = (D^n\nabla_\theta L)(\theta(t)),\\
\frac{d^n}{dt^n}\,\tfrac12\|E\|^2 &= \tfrac12\sum_{j=0}^n\binom nj\ip{E^{(j)}}{E^{(n-j)}} .
\end{aligned}
\end{equation}
\item On the generators, $D$ acts by
\begin{equation}\label{eq:Dgen}
\begin{aligned}
D\,\nabla^k_\theta\Net &= \nabla^{k+1}_\theta\Net[V,\,\cdot\,],\qquad DE = -KE,\qquad DV = -\Hess\,V,\\
V &= -S_1(E),\qquad \nabla^k_\theta L = S_k(E) + G_k ,
\end{aligned}
\end{equation}
so $E^{(n)}$ is a finite sum of contractions of $\nabla_\theta\Net,\dots,\nabla^n_\theta\Net$ with between one and $n$ factors of $E$, obtained from $\Net$ by $n$ applications of the Leibniz rule, and every term of it is bounded by a product of the $H_m$, $\lambda_1^{1/2}$ and $\|E\|$.
In a term with $m$ factors of $E$, the derivatives of order $k_i\ge2$ that occur satisfy $\sum_i(k_i - 1) = m - 1$.
The first three are $E^{(1)} = -KE$, $E^{(2)}$ as in Eq.~\eqref{eq:accel}, and, with $g = g_E$ and $\dot K$ the operator with $\ip{\varphi}{\dot K\psi} = -g_\psi\cdot S(\varphi)g - g_\varphi\cdot S(\psi)g$,
\begin{equation}\label{eq:E3}
\begin{aligned}
E^{(3)} = {}&-K^3E + \dot KKE + K\dot KE - \nabla^3_\theta\Net[g,g,g] - 2\,\nabla^2_\theta\Net[\Hess g,\,g] - \nabla^2_\theta\Net[g,\,S(E)g]\\
&- \Psi\,S(KE)\,g - \Psi\,S_3(E)[g,\,\cdot\,,\,\cdot\,]\,g - \Psi\,S(E)\,\Hess g ,
\end{aligned}
\end{equation}
where $S_3(E)[g,\cdot,\cdot]$ is the matrix $\frac1n\sum_iE_i\cdot\nabla^3_\theta\Net(x_i)[g,\cdot,\cdot]$.
\item \emph{The loss along the gradient, at every order.} Let $\Net$ be $k + 1$ times continuously differentiable, and let $m_k =\nabla^k_\theta L[g_E,\dots,g_E]$ with $k$ copies of $g_E$, so that $m_1 = |g_E|^2$ and $m_2 = \kappa|g_E|^2$. Then
\begin{equation}\label{eq:moments}
\dot m_k = -\,m_{k+1} - k\,\nabla^k_\theta L[\Hess g_E,\,g_E,\dots,g_E] .
\end{equation}
\end{enumerate}
\end{proposition}

\begin{proof}
(a) The chain rule along $\dot\theta = V$; the last identity is the Leibniz rule for $\ip{E}{E}$.
(b) $D$ of a tensor-valued function of $\theta$ is its derivative contracted with $V$, which gives the first rule; the second is Theorem~\ref{thm:identities}(a); the third is $\dot g_E = -\Hess g_E$ of Theorem~\ref{thm:identities}(b) with $V = -g_E$; the last two are Proposition~\ref{prop:lossderiv}.
Each application of $D$ to a contraction raises the order of one tensor by one and inserts $V = -\frac1n\Psi^\top E$, or replaces one $E$ by $-KE$; both keep the term a contraction of derivatives of $\Net$ with factors of $E$, and the count of factors of $E$ rises by at most one per application.
An application that adds a factor of $E$ does so either through $V = -g_E$ inserted into a tensor of order at least one, or through $S(E)$ from $\Hess g$, and in both cases it raises the order of one tensor by one or turns a factor $\Psi$ into $\nabla^2\Net$; the replacement $E\mapsto-KE$ adds only factors of $\Psi$.
Starting from $E$ itself, with one factor of $E$ and no derivative of order two or more, this gives $\sum_i(k_i - 1) = m - 1$.
For Eq.~\eqref{eq:E3}, apply $D$ to the three terms of Eq.~\eqref{eq:accel}: $D(K^2E) = \dot KKE + K\dot KE - K^3E$; $D(\nabla^2\Net[g,g]) = -\nabla^3\Net[g,g,g] + 2\nabla^2\Net[Dg,g]$ with $Dg = -\Hess g$; and $D(\Psi S(E)g) = (D\Psi)S(E)g + \Psi(DS(E))g + \Psi S(E)Dg$ with $D\Psi = -\nabla^2\Net[g,\cdot]$ and $DS(E) = S(DE) + \frac1n\sum_iE_i\cdot\nabla^3\Net(x_i)[V,\cdot,\cdot] = -S(KE) - S_3(E)[g,\cdot,\cdot]$.
The bound on each term is Eq.~\eqref{eq:lossderivbound} together with $\|\Psi w\|\le\lambda_1^{1/2}|w|$ and $|g_E|\le\lambda_1^{1/2}\|E\|$.
(c) $\frac{d}{dt}\nabla^kL(\theta(t))[g,\dots,g] = \nabla^{k+1}L[\dot\theta,g,\dots,g] + k\,\nabla^kL[\dot g,g,\dots,g]$ by symmetry of $\nabla^kL$, with $\dot\theta = -g$ and $\dot g = -\Hess g$.
\end{proof}

Eq.~\eqref{eq:Dn} is the rule that generates the whole series in time of Section~\ref{sec:allorders}; nothing beyond the Leibniz rule is needed to write any order.
Eq.~\eqref{eq:moments} is the tower for the loss along the gradient: the rate of each moment is the next moment plus a term in which the Hessian has acted once.
Its $k = 1$ case is the speed identity, $\dot m_1 = -2m_2$, and its $k = 2$ case gives the rate of $\kappa|g_E|^2$, which is $-m_3 - 2|\Hess g_E|^2$: the curvature along the gradient falls through the square of the Hessian and moves through the third derivative of the loss, which by Proposition~\ref{prop:lossderiv} is $S_3(E) + G_3$.
For a network the hierarchy does not close at any finite $k$.

\paragraph{The same identities for one step.}
For gradient descent the identities hold with exact remainders.
Let $\bar\Hess_j = \int_0^1\Hess(\theta_j - \tau\eta g_j)\,d\tau$ and $\bar\Psi_j = \int_0^1\Psi(\theta_j - \tau\eta g_j)\,d\tau$ be the loss Hessian and the features averaged along the step, with $g_j = g_{E_j}$.

\begin{proposition}[One step]\label{prop:onestep}
Let $\Net$ be twice continuously differentiable in $\theta$. For one step $\theta_{j+1} = \theta_j - \eta g_j$:
\begin{enumerate}[label=(\alph*)]
\item $E_{j+1} = E_j - \eta\bar\Psi_jg_j = (I - \eta K_j)E_j + r_j$ with $r_j = \eta^2\int_0^1(1 - \tau)\nabla^2_\theta\Net(\theta_j - \tau\eta g_j)[g_j,g_j]\,d\tau$; if each $\Net(x_i;\cdot)$ extends holomorphically to a complex ball of radius larger than $\eta|g_j|$ about $\theta_j$, so that its Taylor series converges to it there,
\begin{equation}\label{eq:rseries}
r_j = \sum_{m\ge2}\frac{(-\eta)^m}{m!}\,\nabla^m_\theta\Net(\theta_j)\big[g_j^{\otimes m}\big].
\end{equation}
\item $g_{j+1} = g_j - \eta\bar\Hess_jg_j$, so
\begin{equation}\label{eq:speedstep}
|g_{j+1}|^2 = |g_j|^2 - 2\eta\,g_j^\top\bar\Hess_jg_j + \eta^2|\bar\Hess_jg_j|^2 ,
\end{equation}
and under the same holomorphy hypothesis,
\begin{equation}\label{eq:Hbarseries}
\bar\Hess_j = \sum_{m\ge0}\frac{(-\eta)^m}{(m+1)!}\,\nabla^{m+2}_\theta L(\theta_j)\big[g_j^{\otimes m},\cdot,\cdot\big] .
\end{equation}
\item $\|E_{j+1}\|^2 - \|E_j\|^2 = -2\eta\,g_j\cdot\tfrac1n\bar\Psi_j^\top E_j + \eta^2\|\bar\Psi_jg_j\|^2$.
\end{enumerate}
\end{proposition}

\begin{proof}
(a) Taylor's theorem with integral remainder for $\theta\mapsto\Net(x_i;\theta)$ along the segment, at each training input; the series is the same expansion carried to all orders.
(b) The fundamental theorem of calculus for $\nabla_\theta L$ along the segment, then expand the square; the series is the Taylor expansion of $\Hess$ along the segment, integrated term by term in $\tau$.
(c) Expand $\|E_j - \eta\bar\Psi_jg_j\|^2$ and use $\ip{E_j}{\bar\Psi_jg_j} = g_j\cdot\frac1n\bar\Psi_j^\top E_j$.
\end{proof}

Eq.~\eqref{eq:speedstep} is the discrete form of the speed identity.
Its middle term is $2\eta\,\bar\kappa_j|g_j|^2$, with $\bar\kappa_j = g_j^\top\bar\Hess_jg_j/|g_j|^2$ the step-averaged curvature along the gradient, and the last term, absent in Eq.~\eqref{eq:speed}, allows the gradient to grow in one step even when $\bar\kappa_j > 0$: at the edge of stability, where $\bar\Hess_j$ has an eigenvalue near $2/\eta$ along the top direction, the component of $g_j$ along it is multiplied by a factor near $-1$, which is the oscillation at the edge of stability seen in the gradient (Part~\ref{part:apps}).
Part (a) is Eq.~\eqref{eq:avgstep} of Section~\ref{sec:step} for a gradient step, with the averaged features expanded about the start of the step.

% ===============================================================
\section{The one-time theory along the path}\label{sec:onetime}
% ===============================================================

Every statement of Part~\ref{part:onetime} holds at each point $\theta(t)$ of the path, with the kernel $K(t)$, the error $E(t)$ and the features $\Psi(t)$ of that point: the rates (Proposition~\ref{prop:rates}), the floor (Theorem~\ref{thm:floor}), the one-step bounds and their supremum over roughness functions (Theorems~\ref{thm:stepsup}, \ref{thm:dataset} and~\ref{thm:slow}), and the Krylov series (Theorem~\ref{thm:stepkrylov}).
They are the local theory.
What they do not do is relate the point $\theta(t)$ to the point $\theta_0$, and that is what the next two sections do.
The link is Proposition~\ref{prop:rates}(a): the component of the error along any fixed sampled function $\varphi$ changes at rate at most $\|\varphi\|_{K(t)}$ times the speed of the parameters.
Integrated along the path, this turns the one-step bounds of Section~\ref{sec:step}, in which a step of length $s$ is charged $s\|\varphi\|_K$, into bounds for a stretch of training, in which the path is charged the integral of the kernel size of $\varphi$ against the speed.

% ===============================================================
\section{The path calculus}\label{sec:path}
% ===============================================================

\subsection{Path length and the cost of a direction}

The length of the path up to time $s$ is
\begin{equation}\label{eq:ell}
\ell(s) = \int_0^s|\dot\theta(\tau)|\,d\tau = \int_0^s\|E(\tau)\|_{K(\tau)}\,d\tau ,
\end{equation}
the second form by Theorem~\ref{thm:identities}(a).
For gradient descent, $\ell_N = \eta\sum_{j<N}|g_j|$.
Under the hypotheses of Proposition~\ref{prop:closure} every tensor of the tower at $\theta(s)$ is the tower at $\theta_0$ evaluated at $\Delta(s)$, and in every case $|\Delta(s)|\le\ell(s)$, so the path length is the coordinate in which the motion of the kernel is measured.
The kernel moves with the parameters, not with the clock.
The path length is the counterpart, for a stretch of training, of the rescaled step length $s$ of Section~\ref{sec:step}.

\begin{proposition}[Path length of gradient flow]\label{prop:ell}
Under gradient flow,
\begin{equation}\label{eq:ellbound}
\ell(t)\ \le\ \sqrt{t}\,\Big(\tfrac12\big(\|E_0\|^2 - \|E(t)\|^2\big)\Big)^{1/2}\ \le\ \sqrt{t/2}\,\|E_0\| ,
\qquad
\ell(t) = \|E_0\|_{K_0}\int_0^te^{-\int_0^s\kappa}\,ds .
\end{equation}
\end{proposition}

\begin{proof}
The Cauchy--Schwarz inequality in $s$ applied to Eq.~\eqref{eq:ell} with Eq.~\eqref{eq:energy}; the equality is Eq.~\eqref{eq:speed}.
\end{proof}

The first bound is pessimistic when the error decays, since the speed falls with it; the exact form shows that $\ell$ grows linearly in $t$, at the speed $\|E_0\|_{K_0}$, while $\int\kappa$ is small, and stays bounded as $t\to\infty$ if and only if $\int_0^\infty e^{-\int_0^s\kappa}\,ds < \infty$, for instance when $\kappa\ge\kappa_{\min} > 0$ eventually.

For a sampled function $\varphi$, the cost of the direction $\varphi$ along the path up to time $t$ is
\begin{equation}\label{eq:cost}
\cost_t(\varphi) = \int_0^t\|\varphi\|_{K(s)}\,|\dot\theta(s)|\,ds ,
\end{equation}
the most that the path can have moved $\ip{\varphi}{\Net}$, since $\frac{d}{ds}\ip{\varphi}{\Net} = g_\varphi\cdot\dot\theta$ and $|g_\varphi| = \|\varphi\|_K$ (Proposition~\ref{prop:rates}(a)).
Its first-order term in the path length, by Lemma~\ref{lem:drift} below, is $\ell(t)\|\varphi\|_{K_0}$, the support function of the set $\mathcal{D}_{\ell(t)}$ of Section~\ref{sec:step} with the features at $\theta_0$.

\subsection{Kernel drift along a path: the full series}

\begin{lemma}[Kernel drift along a path]\label{lem:drift}
Let $\Net$ be twice continuously differentiable in $\theta$ for (a) and the general forms of (d), and $M + 1$ times for (c); let $\theta(\cdot)$ be any continuous, piecewise continuously differentiable parameter path of length $\ell = \ell(t)$, and write $D_m(\varphi) = D_m(\varphi;\theta_0)$, $H_m = H_m(\theta_0)$, and $\bar D_m(\varphi;r)$, $\bar H_m(r)$ for the largest values of $D_m(\varphi;\theta)$, $H_m(\theta)$ over the ball $|\theta - \theta_0|\le r$.
For every sampled function $\varphi$ and every $s\le t$:
\begin{enumerate}[label=(\alph*)]
\item \emph{Exact identity.} $g_\varphi(\theta(s)) = g_\varphi(\theta_0) + \int_0^sS(\varphi;\theta(\tau))\,\dot\theta(\tau)\,d\tau$.
\item \emph{Full series.} Under the hypotheses of Proposition~\ref{prop:closure}, with the right sides read as $+\infty$ where the series diverge,
\begin{equation}\label{eq:costseries}
\begin{aligned}
\|\varphi\|_{K(s)}\ &\le\ \sum_{m\ge1}\frac{D_m(\varphi)}{(m-1)!}\,\ell(s)^{m-1},\\
\cost_t(\varphi)\ &\le\ \sum_{m\ge1}\frac{D_m(\varphi)}{m!}\,\ell^m\ \le\ \ell\,\|\varphi\|_{K_0} + \|\varphi\|\sum_{m\ge2}\frac{H_m}{m!}\,\ell^m .
\end{aligned}
\end{equation}
Convergence of the Taylor series of $\Net$ does not by itself make these series of operator norms converge; they do, for instance, when $\Net$ extends to a bounded holomorphic map on a complex ball of radius $R > \ell$ about $\theta_0$, since then Cauchy's estimate along complex lines bounds the diagonal values of $T_m(\varphi;\theta_0)$, and for a symmetric multilinear form on a Hilbert space the operator norm equals the supremum over the diagonal, so that $D_m(\varphi)\le\|\varphi\|\,m!\,\sup\|\Net\|/R^m$.
\item \emph{Resummed tail.} Without analyticity, for every $M\ge1$,
\begin{equation}\label{eq:costresum}
\cost_t(\varphi)\ \le\ \sum_{m=1}^{M}\frac{D_m(\varphi)}{m!}\,\ell^m + \frac{\bar D_{M+1}(\varphi;\ell)}{(M+1)!}\,\ell^{M+1}
\ \le\ \ell\,\|\varphi\|_{K_0} + \tfrac12\,\|\varphi\|\,\bar H_2(\ell)\,\ell^2 ,
\end{equation}
the last form being $M = 1$.
\item \emph{Drift of the gradient and of the kernel.} Let $\sigma(r) = \sum_{m\ge2}\frac{H_m}{(m-1)!}r^{m-1}$ under the hypotheses of (b), or $\sigma(r) = \bar H_2(r)\,r$ in general. Then $|g_\varphi(\theta(s)) - g_\varphi(\theta_0)|\le\|\varphi\|\,\sigma(\ell(s))$, and
\begin{equation}\label{eq:Kdrift}
\begin{aligned}
\|K(s) - K_0\|_{op}\ &\le\ 2\lambda_1(0)^{1/2}\,\sigma(\ell(s)) + \sigma(\ell(s))^2 ,\\
\|(K(s) - K_0)v_k\|\ &\le\ \sigma(\ell(s))\big(\lambda_k^{1/2} + \sigma_k(\ell(s))\big) + \lambda_1(0)^{1/2}\,\sigma_k(\ell(s)),
\end{aligned}
\end{equation}
for every eigenpair $(\lambda_k, v_k)$ of $K_0$, where $\sigma_k(r) = \sum_{m\ge2}\frac{D_m(v_k)}{(m-1)!}r^{m-1}$, or $\bar D_2(v_k;r)\,r$ in general, is the drift of the gradient of the direction $v_k$ itself.
\end{enumerate}
\end{lemma}

\begin{proof}
(a) The chain rule, $\frac{d}{d\tau}g_\varphi(\theta(\tau)) = S(\varphi;\theta(\tau))\,\dot\theta(\tau)$, integrated along the path; no update rule is used.
(b) Eq.~\eqref{eq:closure} with $m = 1$, each term bounded by its operator norm and $|\Delta(s)|\le\ell(s)$; insert into Eq.~\eqref{eq:cost} with $|\dot\theta(s)| = \ell'(s)$ and $\int_0^t\ell^{m-1}\ell'\,ds = \ell^m/m$, term by term, which is allowed because every term is nonnegative; the last inequality is $D_m\le\|\varphi\|H_m$.
(c) Taylor's theorem with integral remainder for $\nabla f_\varphi$ along the segment from $\theta_0$ to $\theta(s)$, which lies in the ball of radius $\ell(s)$, then integrate as in (b).
(d) The bound on the drift of $g_\varphi$ is the same expansion as in (b) or (c), applied to $g_\varphi$ itself along the segment from $\theta_0$ to $\theta(s)$, with the $m = 1$ term removed.
For a unit $\varphi$, $\ip{\varphi}{(K(s) - K_0)\varphi} = |g_\varphi(s)|^2 - |g_\varphi(0)|^2 = 2g_\varphi(0)\cdot\delta + |\delta|^2$ with $\delta = g_\varphi(s) - g_\varphi(0)$ and $|g_\varphi(0)|\le\lambda_1(0)^{1/2}$, and the norm of a symmetric operator is the largest absolute value of its quadratic form on unit vectors.
For the second bound, for a unit $\psi$, $\ip{\psi}{(K(s) - K_0)v_k} = (g_\psi(s) - g_\psi(0))\cdot g_{v_k}(s) + g_\psi(0)\cdot(g_{v_k}(s) - g_{v_k}(0))$, with $|g_{v_k}(s)|\le\lambda_k^{1/2} + \sigma_k$ and $|g_\psi(0)|\le\lambda_1(0)^{1/2}$.
\end{proof}

Part (b) is the central estimate of the path calculus: the cost of a direction along a path of length $\ell$ is at most the Taylor series of the network's derivatives at the start, evaluated at $\ell$, with $\ell\|\varphi\|_{K_0}$ as its first-order term and every further order kept.
Part (c) is what the series becomes when its convergence is not assumed: any number of terms, with the whole tail resummed into one remainder by the largest derivative of the next order over the ball that the path stays in; the form with $\bar H_2$ is the whole tail, not a second-order approximation.
Part (d) says how the kernel itself moves.
The second bound in Eq.~\eqref{eq:Kdrift} is the one that matters for the slow directions: the kernel's action on $v_k$ drifts by a term proportional to the drift of all gradients, $\sigma$, times the size $\lambda_k^{1/2}$ of $v_k$'s own gradient, plus a term proportional to the drift of $v_k$'s own gradient, $\sigma_k$, times $\lambda_1^{1/2}$.
A slow direction is protected from the first term by its small $\lambda_k$ and from the second only to the extent that its own higher derivatives $D_m(v_k)$ are small.
That second term is where the kernel can grow by large factors on small eigenvalues with small absolute motion, which is what the runs of the edge-of-stability application show (Part~\ref{part:apps}).
By Corollary~\ref{cor:budget}, an eigendirection $v_k$ whose roughness $\Rgh(v_k)$ is large compared with $\sR$ is slow, $\lambda_k\le\sR^2/\Rgh(v_k)^2$; so the growth of the kernel on rough directions, which is what learning a rough dataset requires, is controlled by the higher derivatives of the network along those directions.

\subsection{The floor along a path, and from time 0}

The floor of Part~\ref{part:onetime} holds at time $t$ with the output and kernel of that time.
The next two theorems bound the error at time $t$ by quantities fixed at an earlier time and the path in between.

\begin{theorem}[Floor along a path]\label{thm:path}
Let $\Net$ be continuously differentiable in $\theta$.
For every time $t$,
\begin{equation}\label{eq:pathfloor}
\|E(t)\|\ \ge\ F_t := \sup_{\|\varphi\|\le1}\Big(\ip{\varphi}{E_0} - \cost_t(\varphi)\Big),
\end{equation}
and $F_t = \min\{\|E_0 + u\| : u\in\mathcal{C}_t\}$, where $\mathcal{C}_t$ is the compact convex symmetric set whose support function is $\cost_t$ and which contains $\Net(t) - \Net(0)$.
For every roughness function $\Rgh$, with $\sR(s)$ computed from $K(s)$,
\begin{equation}\label{eq:pathrough}
F_t\ \ge\ \frac{1}{\CR}\Big(\Rgh(Y) - \Rgh(\hat Y_0) - \int_0^t\sR(s)\,|\dot\theta(s)|\,ds\Big).
\end{equation}
\end{theorem}

\begin{proof}
For $\|\varphi\|\le1$, $\|E(t)\|\ge\ip{\varphi}{E(t)} = \ip{\varphi}{E_0} + \int_0^tg_\varphi\cdot\dot\theta\,ds\ge\ip{\varphi}{E_0} - \cost_t(\varphi)$.
$\cost_t$ is a seminorm on sampled functions, since each $\|\cdot\|_{K(s)}$ is, and it is bounded by $\ell(t)\max_s\lambda_1(s)^{1/2}\|\cdot\|$; it is therefore the support function of a compact convex symmetric set $\mathcal{C}_t$, and $\Net(t) - \Net(0)\in\mathcal{C}_t$ because $\ip{\varphi}{\Net(t) - \Net(0)}\le\cost_t(\varphi)$ for every $\varphi$.
The minimum form is the duality between the distance from a point to a compact convex set and the support function of that set, as in the proof of Theorem~\ref{thm:stepsup}(b).
For Eq.~\eqref{eq:pathrough}, let $\varphi_0$ be the element of the dual ball $D_\Rgh$ of Appendix~\ref{app:step} that attains $\Rgh(E_0)$, for which $\|\varphi_0\|\le\CR$ and $\|\varphi_0\|_{K(s)}\le\sR(s)$ at every $s$, take $\varphi = \varphi_0/\CR$, and use the reverse triangle inequality for $E_0 = \hat Y_0 - Y$.
\end{proof}

Theorem~\ref{thm:path} is Theorem~\ref{thm:dataset} for a stretch of training.
The step of length $s$ is replaced by the path, and the peak feature roughness over the set that one step can reach is replaced by its integral against the speed.
Expanded about $\theta_0$ (Theorem~\ref{thm:time0} below), $F_t$ is at least $F(\ell(t))$ of Eq.~\eqref{eq:Fs} with the features at $\theta_0$, minus a remainder of second order in $\ell(t)$: to first order, the whole stretch of training is charged as one step of length $\ell(t)$, and the Krylov series of Theorem~\ref{thm:stepkrylov} computes that term.
The dataset enters through $\Rgh(Y)$ alone, as in Part~\ref{part:onetime}: to bring the error below $(\Rgh(Y) - \Rgh(\hat Y_0))/\CR$, the path must supply roughness, at rate at most $\sR$ per unit length.

\begin{theorem}[The floor at time 0 bounds the error at time $t$]\label{thm:time0}
Let $F_0(\ell) = \max_{\|\varphi\|\le1}(\ip{\varphi}{E_0} - \ell\|\varphi\|_{K_0})$; by the duality of Theorem~\ref{thm:stepsup}, valid also when $F(\ell) = 0$, it equals $F(\ell)$ of Eq.~\eqref{eq:Fs} with $c_p = 1$ and the features at $\theta_0$, the smallest error of the first-order expansion about $\theta_0$ of a step of length at most $\ell$ (Theorem~\ref{thm:stepsup}, computed by the Krylov series of Theorem~\ref{thm:stepkrylov}).
Under the hypotheses of Lemma~\ref{lem:drift}, for every $t$ with path length $\ell = \ell(t)$:
\begin{equation}\label{eq:time0series}
\|E(t)\|\ \ge\ \sup_{\|\varphi\|\le1}\Big(\ip{\varphi}{E_0} - \sum_{m\ge1}\frac{D_m(\varphi)}{m!}\,\ell^m\Big)\ \ge\ F_0(\ell) - \sum_{m\ge2}\frac{H_m}{m!}\,\ell^m
\end{equation}
under the hypotheses of Proposition~\ref{prop:closure}, and in general
\begin{equation}\label{eq:time0}
\|E(t)\|\ \ge\ F_0(\ell) - \tfrac12\,\bar H_2(\ell)\,\ell^2,\qquad
\|E(t)\|\ \ge\ \frac{\Rgh(Y) - \Rgh(\hat Y_0) - \ell\,\sR - \tfrac12\,\CR\,\bar H_2(\ell)\,\ell^2}{\CR}
\end{equation}
for every roughness function $\Rgh$ with $\sR$ computed from $K_0$.
\end{theorem}

\begin{proof}
Insert Eq.~\eqref{eq:costseries} or Eq.~\eqref{eq:costresum} into Eq.~\eqref{eq:pathfloor}, using $\|\varphi\|\le1$ in the tail; the supremum of the first term is $F_0(\ell)$.
For the roughness form, along the functional $\varphi_0/\CR$ of Theorem~\ref{thm:path}, $\sR(s)\le\sR + \CR\bar H_2(\ell(s))\ell(s)$ by Lemma~\ref{lem:drift}(d).
\end{proof}

Theorem~\ref{thm:time0} is a prediction: its right side is made of the initial error, the initial kernel, the path length, and the derivatives of the network at the start.
For a local complexity function, $\Rgh(Y)$ is computed before training, so the absolute roughness of the dataset bounds the error after any stretch of training, minus the roughness of the initial output, minus the one-step term of Theorem~\ref{thm:dataset}, which is linear in $\ell$, and minus the kernel-motion series, which starts at second order in $\ell$.
The bound is informative while the series is below $F_0(\ell)$, up to a path length $\ell_\ast$; beyond it, only the floor at time $t$ with the kernel at time $t$ remains.
This is the integrated form of the main statement of Part~\ref{part:onetime}: the roughness that the labels demand must be supplied by parameter motion, the path length measures that motion, and as long as the kernel has not moved far the supply rate is the one at the start.

\subsection{A rough dataset needs a minimum path length and a minimum training time}\label{sec:mintime}

Theorems~\ref{thm:path} and~\ref{thm:time0} bound the error at time $t$ from below.
Read the other way, they say how much parameter motion, and under gradient flow how much training time, is needed before the error can reach a given value.
This form holds at every time, with no horizon: the kernel motion enters only through the largest second derivative over the ball that the path stays in.

For a roughness function $\Rgh$ and $\ell\ge0$, let
\begin{equation}\label{eq:budgetfn}
\mathcal{S}_\Rgh(\ell) = \ell\,\sR(\theta_0) + \tfrac12\,\CR\,\bar H_2(\ell)\,\ell^2 ,
\end{equation}
the most roughness that a path of length $\ell$ from $\theta_0$ can add to the output, with $\sR(\theta_0)$ the peak feature roughness at the start and $\bar H_2(\ell)$ as in Lemma~\ref{lem:drift}.
Under the hypotheses of Proposition~\ref{prop:closure} the second term can be replaced by $\CR\sum_{m\ge2}H_m\ell^m/m!$.
$\mathcal{S}_\Rgh$ is continuous and nondecreasing, with $\mathcal{S}_\Rgh(0) = 0$.
For $\varepsilon\ge0$, let
\begin{equation}\label{eq:elleps}
\ell_\varepsilon^\Rgh = \inf\big\{\ell\ge0 : \mathcal{S}_\Rgh(\ell)\ \ge\ \Rgh(Y) - \Rgh(\hat Y_0) - \CR\,\varepsilon\big\},
\end{equation}
which is $0$ when the right side is not positive, and let $\ell_\varepsilon^\star = \inf\{\ell\ge0 : F_0(\ell) - \tfrac12\bar H_2(\ell)\,\ell^2\le\varepsilon\}$, with $F_0$ from Theorem~\ref{thm:time0}.
An empty set has infimum $+\infty$, in which case part (b) below says that the error never reaches $\varepsilon$.
Under the hypotheses of Proposition~\ref{prop:closure}, $\ell_\varepsilon^\star$ may be defined with $\sum_{m\ge2}H_m\ell^m/m!$ in place of $\tfrac12\bar H_2(\ell)\ell^2$; Eq.~\eqref{eq:minpath} below holds when $\ell_\varepsilon^\star$ and $\mathcal{S}_\Rgh$ are taken in the same form.

\begin{corollary}[Minimum path length and training time]\label{cor:mintime}
Let $\Net$ be twice continuously differentiable in $\theta$, and let $\varepsilon\ge0$.
\begin{enumerate}[label=(\alph*)]
\item \emph{Roughness that a path can supply.} For every continuous, piecewise continuously differentiable parameter path, whatever update rule produced it, and every roughness function $\Rgh$,
\begin{equation}\label{eq:supplybound}
\|E(t)\|\ \ge\ \frac{\Rgh(Y) - \Rgh(\hat Y_0) - \mathcal{S}_\Rgh(\ell(t))}{\CR}.
\end{equation}
\item \emph{Minimum path length.} If $\|E(t)\|\le\varepsilon$, then
\begin{equation}\label{eq:minpath}
\ell(t)\ \ge\ \ell_\varepsilon^\star\ \ge\ \sup_\Rgh\,\ell_\varepsilon^\Rgh ,
\end{equation}
the supremum over all roughness functions.
\item \emph{Minimum training time.} Under gradient flow, if $\|E(t)\|\le\varepsilon<\|E_0\|$, then
\begin{equation}\label{eq:mintime}
t\ \ge\ \frac{2\,(\ell_\varepsilon^\star)^2}{\|E_0\|^2 - \varepsilon^2}\ \ge\ \frac{2\,(\ell_\varepsilon^\star)^2}{\|E_0\|^2} .
\end{equation}
\end{enumerate}
\end{corollary}

\begin{proof}
(a) Write $\ell = \ell(t)$.
The path up to time $t$ lies in the ball of radius $\ell$ about $\theta_0$, and $\theta(s)$ lies in the ball of radius $\ell(s)$.
The peak feature roughness is the largest kernel size over the dual ball of $\Rgh$, $\sR(\theta) = \max_{\varphi\in D_\Rgh}\|\varphi\|_{K(\theta)}$, since $\Rgh(\Psi g) = \max_{\varphi\in D_\Rgh}g\cdot g_\varphi$ (Appendix~\ref{app:step}).
Every $\varphi\in D_\Rgh$ has $\|\varphi\|\le\CR$, so by Lemma~\ref{lem:drift}(d) in its general form, $\|\varphi\|_{K(\theta(s))}\le\|\varphi\|_{K_0} + \CR\,\bar H_2(\ell(s))\,\ell(s)$, and hence $\sR(\theta(s))\le\sR(\theta_0) + \CR\,\bar H_2(\ell)\,\ell(s)$.
With $|\dot\theta(s)|\,ds = d\ell(s)$,
\[
\int_0^t\sR(\theta(s))\,|\dot\theta(s)|\,ds\ \le\ \ell\,\sR(\theta_0) + \CR\,\bar H_2(\ell)\int_0^t\ell(s)\,d\ell(s)\ =\ \mathcal{S}_\Rgh(\ell).
\]
Insert this into Eq.~\eqref{eq:pathrough} of Theorem~\ref{thm:path}, whose proof uses only $\frac{d}{ds}\ip{\varphi}{\Net} = g_\varphi\cdot\dot\theta$ and therefore holds for every such path.
Under the hypotheses of Proposition~\ref{prop:closure}, Lemma~\ref{lem:drift}(b) gives the series form in the same way.
(b) Theorem~\ref{thm:time0} gives $\|E(t)\|\ge F_0(\ell) - \tfrac12\bar H_2(\ell)\ell^2$; its proof uses only Theorem~\ref{thm:path} and Lemma~\ref{lem:drift}, so it holds for every such path.
$F_0$ is nonincreasing and $\bar H_2(\ell)\ell^2$ is nondecreasing, so the right side is nonincreasing in $\ell$; if $\|E(t)\|\le\varepsilon$ then $\ell(t)$ belongs to the set that defines $\ell_\varepsilon^\star$, and $\ell(t)\ge\ell_\varepsilon^\star$.
For the second inequality, Theorem~\ref{thm:stepsup}(a) and the reverse triangle inequality give, for every $\Rgh$ and $\ell$,
\[
F_0(\ell) - \tfrac12\bar H_2(\ell)\ell^2\ \ge\ \frac{\Rgh(E_0) - \ell\,\sR(\theta_0) - \tfrac12\CR\bar H_2(\ell)\ell^2}{\CR}\ \ge\ \frac{\Rgh(Y) - \Rgh(\hat Y_0) - \mathcal{S}_\Rgh(\ell)}{\CR}.
\]
If $\ell<\ell_\varepsilon^\Rgh$, then $\mathcal{S}_\Rgh(\ell) < \Rgh(Y) - \Rgh(\hat Y_0) - \CR\varepsilon$, so the right side, and with it the left side, exceeds $\varepsilon$, and $\ell$ is not in the set that defines $\ell_\varepsilon^\star$; that set is an interval $[\ell_\varepsilon^\star,\infty)$ or $(\ell_\varepsilon^\star,\infty)$ because the left side is nonincreasing, so $\ell_\varepsilon^\star\ge\ell_\varepsilon^\Rgh$.
(c) Since $\|E_0\|>\varepsilon\ge\|E(t)\|$ and the error is continuous, there is a first time $t'\le t$ with $\|E(t')\| = \varepsilon$.
By (b) at $t'$, $\ell(t')\ge\ell_\varepsilon^\star$, and by Proposition~\ref{prop:ell}, $\ell(t')^2\le\tfrac{t'}{2}\big(\|E_0\|^2 - \varepsilon^2\big)$.
Hence $t\ge t'\ge2(\ell_\varepsilon^\star)^2/(\|E_0\|^2 - \varepsilon^2)$.
\end{proof}

Corollary~\ref{cor:mintime} is the statement of the theory about training as a whole.
The labels demand roughness $\Rgh(Y)$; the network can supply roughness only at the rate $\sR$ per unit of parameter motion, a rate that can grow along the path only through the second derivative of the network; the path length measures the parameter motion; and under gradient flow the path length can grow no faster than $\sqrt{t/2}\,\|E_0\|$.
So a rougher dataset needs a longer path, and a longer training time, before the error can fall to a given value, and the bound holds at every time with no horizon.
For a local complexity function, every quantity in $\ell_\varepsilon^\Rgh$ except $\bar H_2$ is computed before training, and $\bar H_2$ is the largest second derivative of the network over the ball that the path stays in.
The Krylov form $\ell_\varepsilon^\star$ needs no choice of seminorm and is never worse than any roughness function.

The corollary is one-sided, like everything in Part~\ref{part:onetime}: it says that a rough dataset cannot be learned with a short path, not that a long path learns it.
It is informative when the peak feature roughness and the second-derivative term are small against the absolute roughness of the labels, so that $\ell_\varepsilon$ is large.
Part (a) and (b) hold for every update rule, including Adam with $\Pi = I$ in the definition of the path length; part (c) is for gradient flow, and for another update rule the path length has to be measured or bounded from the step sizes and gradient norms.
For two problems compared at the same tolerance, such as two learned time steppers in state units, the corollary bounds from below the training each one needs; whether the one with the larger bound in fact trains longer is not a consequence of it.

\subsection{Two-sided bounds for short paths}

The floor is a lower bound.
The identity $\dot E = -KE$, expanded about the kernel at $\theta_0$ with the whole motion of the kernel kept, also gives an upper bound, for as long as that motion is small.
Nothing in the theorem holds the kernel fixed: $\delta K(s)$ is the exact motion along the path, and the bound on it is the series of Lemma~\ref{lem:drift}.

\begin{theorem}[Expansion of the error about the kernel at the start]\label{thm:duhamel}
Along gradient flow, under the hypotheses of Lemma~\ref{lem:drift}, with $\delta K(s) = K(s) - K_0$,
\begin{equation}\label{eq:duhamel}
E(t) = e^{-K_0t}E_0 - \int_0^te^{-K_0(t-s)}\,\delta K(s)\,E(s)\,ds ,
\end{equation}
and therefore
\begin{equation}\label{eq:duhamelbound}
\begin{aligned}
\Big|\,\|E(t)\| - \|e^{-K_0t}E_0\|\,\Big|\ &\le\ \int_0^t\|\delta K(s)\|_{op}\,\|E(s)\|\,ds\\
&\le\ \|E_0\|\int_0^t\Big(2\lambda_1(0)^{1/2}\sigma(\ell(s)) + \sigma(\ell(s))^2\Big)ds ,
\end{aligned}
\end{equation}
and for every eigenpair $(\lambda_k, v_k)$ of $K_0$,
\begin{equation}\label{eq:duhamelk}
\Big|\,\hat e_k(t) - e^{-\lambda_kt}\hat e_k(0)\,\Big|\ \le\ \int_0^te^{-\lambda_k(t-s)}\,\|\delta K(s)v_k\|\,\|E(s)\|\,ds ,
\end{equation}
with $\|\delta K(s)v_k\|$ bounded by Eq.~\eqref{eq:Kdrift}.
\end{theorem}

\begin{proof}
$\dot E = -K_0E - \delta K(s)E$ is a linear equation with a forcing term, and Eq.~\eqref{eq:duhamel} is its variation-of-constants form.
The first bound uses $\|e^{-K_0(t-s)}\|_{op}\le1$ and the triangle inequality, the second Lemma~\ref{lem:drift}(d) and $\|E(s)\|\le\|E_0\|$ from Theorem~\ref{thm:identities}(a).
Eq.~\eqref{eq:duhamelk} is Eq.~\eqref{eq:duhamel} projected on $v_k$, using $\ip{v_k}{\delta K(s)E(s)} = \ip{\delta K(s)v_k}{E(s)}$ because $\delta K(s)$ is symmetric.
\end{proof}

Eq.~\eqref{eq:duhamelbound} is two-sided: the error at time $t$ lies within a band around the zeroth-order term $e^{-K_0t}E_0$ of the expansion, whose half-width is $\|E_0\|$ times the time integral of $2\lambda_1(0)^{1/2}\sigma(\ell(s)) + \sigma(\ell(s))^2$, a series in the path length starting at first order with coefficient $2\lambda_1(0)^{1/2}H_2$; since $\ell$ is nondecreasing, the half-width is at most $\|E_0\|\,t\,\big(2\lambda_1(0)^{1/2}\sigma(\ell(t)) + \sigma(\ell(t))^2\big)$.
It is the companion of Theorem~\ref{thm:time0}: that theorem's band is one-sided and its correction starts at second order in $\ell$, this one's is two-sided and starts at first order in $\ell$ with a factor of the time, and both become uninformative once the series exceed the quantities they correct.
Eq.~\eqref{eq:duhamelk} is the statement for one component.
On a slow direction the factor $e^{-\lambda_k(t-s)}$ forgets nothing, so the component is held near its initial value as long as the kernel's action on $v_k$ has not drifted, and by Eq.~\eqref{eq:Kdrift} that drift is governed by $\sigma$ and $\sigma_k$, the series of the network's own higher derivatives.
This is the floor of Part~\ref{part:onetime} seen over a stretch of training: the component of the error on a slow direction, and so by Corollary~\ref{cor:budget} on every eigendirection that is rough compared with the peak feature roughness, stays where it was until the kernel moves on that direction.
Nothing here says that the slow components do decay; the theorem says that whenever they decay beyond the factor $e^{-\lambda_kt}$, it is through $\delta K(s)v_k$, whose size is bounded but whose direction is not.

% ===============================================================
\section{The error to all orders}\label{sec:allorders}
% ===============================================================

The identities of Section~\ref{sec:identities} and the bounds of Section~\ref{sec:path} are the first terms of expansions that can be carried to any order, with every order computable and with equality in the limit.
This section writes the expansions out, proves where each converges, and states what each one needs in order to describe a stretch of training from the initial state.
There are three, with different expansion variables.
The one that matters is the one in which the linear part of the dynamics has already been summed, so that the expansion variable is the motion of the kernel and not the time.

\subsection{The series in time}

Assume that $V = -\nabla_\theta L$ extends to a bounded holomorphic map on the complex ball $|\theta - \theta_0|\le R$ in $\mathbb{C}^{d_\theta}$, and let $M_R = \sup$ of $|V|$ on that ball.
Then the path is analytic in $t$ and
\begin{equation}\label{eq:timeseries}
E(t) = \sum_{n\ge0}\frac{t^n}{n!}\,E^{(n)}(0)\qquad\text{for }|t| < \rho_t,\qquad \rho_t\ \ge\ \frac{R}{2M_R} .
\end{equation}
The radius is Cauchy's theorem on analytic differential equations in the form of a Picard iteration in complex time: on the disc $|t| < R/(2M_R)$, iterates started at $\theta_0$ stay in the ball of radius $R/2$, where $V$ is Lipschitz with constant at most $2M_R/R$ by Cauchy's estimate, so the iteration contracts and its uniform limit is analytic.
By Proposition~\ref{prop:timederiv}, every coefficient is explicit in the tower at $\theta_0$: $E^{(n)}(0) = (D^n\Net)(\theta_0)$, a sum of contractions of $\nabla\Net,\dots,\nabla^n\Net$ at $\theta_0$ with $E_0$, with $n = 1, 2, 3$ written out there.
This is the series that ``one order higher'' refers to.
It is exact inside its radius, but it is the least efficient expansion for training, because its natural variable is $\lambda_1t$.
Among the terms of $E^{(n)}(0)$ are $(-K_0)^nE_0$, whose sum over $n$ is the exponential series of $e^{-K_0t}E_0$; truncated at order $p$, that part alone leaves the error $\big(\sum_k\hat e_k(0)^2\,r_p(\lambda_kt)^2\big)^{1/2}$, with $|r_p(x)|\le x^{p+1}/(p+1)!$ the remainder of the exponential series, so the series in time is useful only for $\lambda_1t$ of order one unless $p$ grows with $\lambda_1t$: its truncation error is set by the fast directions, which have long decayed and which the floor says carry no error, while the slow directions, where the error lies, need $\lambda_1t$ of the order of $\lambda_1/\lambda_k$.
The fast part of the dynamics has to be summed before anything is expanded.

\subsection{The Dyson series in the motion of the kernel}

\begin{theorem}[Dyson series]\label{thm:dyson}
Let $\Net$ be twice continuously differentiable in $\theta$, let $\theta(t)$ be the gradient-flow path, let $\delta K(s) = K(s) - K_0$ along it, and let $\beta(t) = \int_0^t\|\delta K(s)\|_{op}\,ds$.
Then for every $t$
\begin{equation}\label{eq:dyson}
\begin{aligned}
E(t) &= \sum_{n\ge0}(-1)^n\mathcal{I}_n(t),\\
\mathcal{I}_n(t) &= \int_{t > s_1 > \dots > s_n > 0}e^{-K_0(t - s_1)}\,\delta K(s_1)\,e^{-K_0(s_1 - s_2)}\,\delta K(s_2)\cdots\delta K(s_n)\,e^{-K_0s_n}E_0\;ds_n\cdots ds_1 ,
\end{aligned}
\end{equation}
with $\mathcal{I}_0(t) = e^{-K_0t}E_0$, the series converging absolutely for every $t$ with
\begin{equation}\label{eq:dysonbound}
\|\mathcal{I}_n(t)\|\le\|E_0\|\,\frac{\beta(t)^n}{n!},\qquad
\Big\|E(t) - \sum_{n = 0}^{p}(-1)^n\mathcal{I}_n(t)\Big\|\le\|E_0\|\,\frac{\beta(t)^{p+1}}{(p+1)!} .
\end{equation}
Under the hypotheses of Lemma~\ref{lem:drift}, $\beta(t)\le\int_0^t\big(2\lambda_1(0)^{1/2}\sigma(\ell(s)) + \sigma(\ell(s))^2\big)ds$.
\end{theorem}

\begin{proof}
Eq.~\eqref{eq:duhamel} is $E = \mathcal{I}_0 - \mathcal{V}E$ with $(\mathcal{V}E)(t) = \int_0^te^{-K_0(t-s)}\delta K(s)E(s)\,ds$, a linear Volterra operator with a continuous kernel.
Substituting the equation into itself $p$ times gives $E = \sum_{n\le p}(-1)^n\mathcal{V}^n\mathcal{I}_0 + (-1)^{p+1}\mathcal{V}^{p+1}E$, and $\mathcal{V}^n\mathcal{I}_0 = \mathcal{I}_n$.
Since $\|e^{-K_0\tau}\|_{op}\le1$,
\[
\|(\mathcal{V}^nF)(t)\|\ \le\ \sup_{s\le t}\|F(s)\|\int_{t>s_1>\dots>s_n>0}\prod_i\|\delta K(s_i)\|_{op}\,ds\ =\ \sup_{s\le t}\|F(s)\|\,\frac{\beta(t)^n}{n!},
\]
the last step because the integral of a product of one function over the ordered simplex is $1/n!$ times the $n$-th power of its integral.
With $F = \mathcal{I}_0$ this is the term bound, and with $F = E$ and $\|E(s)\|\le\|E_0\|$ (Theorem~\ref{thm:identities}(a)) the remainder is at most $\|E_0\|\beta^{p+1}/(p+1)!$; the series converges absolutely by comparison with $e^{\beta}$.
The bound on $\beta$ is Lemma~\ref{lem:drift}(d).
\end{proof}

The expansion variable is $\beta$, the integrated motion of the kernel, and not the time: the exponentials carry the whole of the linear dynamics, at every order.
Term $n = 0$ is $e^{-K_0t}E_0$, and term $n = 1$ is the integral of Theorem~\ref{thm:duhamel}.
The series is exact for every $t$ and converges for every $t$, but it is written in terms of $\delta K(s)$ along the path that was taken.
It is a reorganization of the path, not a prediction from the start.
To predict, $\delta K(s)$ must itself be expanded in $\Delta(s)$ by the Taylor closure, and $\Delta(s)$ in turn, which is the scheme of the next subsection.

\subsection{Picard iteration from the initial state}

The closed system of Proposition~\ref{prop:closure} is $\dot\Delta = -\nabla_\theta L(\theta_0 + \Delta)$.
Split its right side into its first-order Taylor expansion about $\theta_0$ in the outputs and the rest:
\begin{equation}\label{eq:Nsplit}
\nabla_\theta L(\theta_0 + \Delta) = g_{E_0} + G_0\Delta + N(\Delta),\qquad N(\Delta) = \nabla_\theta L(\theta_0 + \Delta) - g_{E_0} - G_0\Delta ,
\end{equation}
with $G_0 = \frac1n\Psi_0^\top\Psi_0$.
$N(0) = 0$, $\nabla N(\Delta) = \Hess(\theta_0 + \Delta) - G_0$, and under the Taylor closure $N$ is the series of all the terms of $\frac1n\Psi(\theta_0 + \Delta)^\top\big(\Net(\theta_0 + \Delta) - Y\big)$ beyond the two kept, whose first term is $S(E_0;\theta_0)\Delta$, with $S(E_0;\theta_0)$ the matrix $S(E_0)$ at $\theta_0$.
Let
\begin{equation}\label{eq:LambdaN}
\Lambda_N(r) = \sup_{|\theta - \theta_0|\le r}\|\Hess(\theta) - G_0\|_{op},
\end{equation}
the Lipschitz constant of $N$ on the ball of radius $r$: the curvature of the loss beyond its kernel term at $\theta_0$.
By Proposition~\ref{prop:hessian} and Lemma~\ref{lem:drift}(d),
\begin{equation}\label{eq:LambdaNbound}
\Lambda_N(r)\ \le\ \bar H_2(r)\Big(\|E_0\| + \big(\lambda_1(0)^{1/2} + \sigma(r)\big)r\Big) + \big(2\lambda_1(0)^{1/2} + \sigma(r)\big)\sigma(r),
\end{equation}
and $\Lambda_N(r)\to\|S(E_0;\theta_0)\|_{op}\le H_2\|E_0\|$ as $r\to0$, a limit that cannot be used at a fixed time, since the ball must contain the path up to that time.

\begin{theorem}[Picard scheme from the initial state]\label{thm:picard}
Let $\Net$ be twice continuously differentiable in $\theta$.
Let $\Delta^{\mathrm{lin}}(t) = -\int_0^te^{-G_0(t-s)}g_{E_0}\,ds$, the solution of $\dot\Delta = -(g_{E_0} + G_0\Delta)$, the zeroth iterate, for which $E_0 + \Psi_0\Delta^{\mathrm{lin}}(t) = e^{-K_0t}E_0$, and define
\begin{equation}\label{eq:picard}
\begin{aligned}
\Delta_{[0]} &= \Delta^{\mathrm{lin}},\\
\Delta_{[k+1]}(t) &= \Delta^{\mathrm{lin}}(t) - \int_0^te^{-G_0(t-s)}\,N\big(\Delta_{[k]}(s)\big)\,ds,\\
E_{[k]}(t) &= \Net\big(\theta_0 + \Delta_{[k]}(t)\big) - Y ,
\end{aligned}
\end{equation}
the brackets distinguishing the iterates from the time derivatives $E^{(n)}$ of Proposition~\ref{prop:timederiv}.
\begin{enumerate}[label=(\alph*)]
\item \emph{Short windows, geometric rate.} Let $r > 0$ and $T > 0$ satisfy
\begin{equation}\label{eq:picardcond}
\mu := T\,\Lambda_N(r)\ \le\ \tfrac12\qquad\text{and}\qquad \sqrt{T/2}\,\|E_0\|\ \le\ \tfrac12\,r .
\end{equation}
Then the gradient-flow path stays in the ball of radius $r$ on $[0,T]$, and for all $t\le T$
\begin{equation}\label{eq:picardconv}
\begin{aligned}
|\Delta_{[k]}(t) - \Delta(t)|\ &\le\ \frac{\mu^k}{1 - \mu}\sup_{s\le T}|\Delta_{[1]}(s) - \Delta_{[0]}(s)|,\\
\|E_{[k]}(t) - E(t)\|\ &\le\ \big(\lambda_1(0)^{1/2} + \sigma(r)\big)\,|\Delta_{[k]}(t) - \Delta(t)| .
\end{aligned}
\end{equation}
\item \emph{Any finite window.} Let $T > 0$ and $r\ge\sqrt{2T}\,\|E_0\|$, and let $P_r$ be the nearest-point projection onto the closed ball of radius $r$.
Define the iterates of Eq.~\eqref{eq:picard} with $N\circ P_r$ in place of $N$, and $E_{[k]}(t) = \Net\big(\theta_0 + P_r\Delta_{[k]}(t)\big) - Y$.
Then, with no condition on $T\Lambda_N(r)$, for all $t\le T$
\begin{equation}\label{eq:picardany}
\begin{aligned}
|\Delta_{[k]}(t) - \Delta(t)|\ &\le\ \frac{(\Lambda_N(r)\,t)^k}{k!}\,\sqrt{2t}\,\|E_0\|,\\
\|E_{[k]}(t) - E(t)\|\ &\le\ \big(\lambda_1(0)^{1/2} + \sigma(r)\big)\,|\Delta_{[k]}(t) - \Delta(t)| .
\end{aligned}
\end{equation}
\end{enumerate}
In both cases the iterates evaluate the network only on the ball of radius $r$, and if each $\Net(x_i;\cdot)$ extends holomorphically to a complex ball of radius larger than $r$ about $\theta_0$, every iterate is computed from the tower at $\theta_0$ alone.
\end{theorem}

\begin{proof}
(a) By variation of constants, the closed system is equivalent on $[0,T]$ to the integral equation $\Delta = \mathcal{P}[\Delta]$ with $\mathcal{P}[\Delta](t) = \Delta^{\mathrm{lin}}(t) - \int_0^te^{-G_0(t-s)}N(\Delta(s))\,ds$, and the iteration is $\Delta_{[k+1]} = \mathcal{P}[\Delta_{[k]}]$.
Let $X$ be the continuous paths on $[0,T]$ with $\sup_t|\Delta(t)|\le r$, with the supremum norm.
Since $G_0\succeq0$, $\|e^{-G_0\tau}\|_{op}\le1$, so for $\Delta_1,\Delta_2\in X$, $\|\mathcal{P}[\Delta_1] - \mathcal{P}[\Delta_2]\|_\infty\le T\Lambda_N(r)\|\Delta_1 - \Delta_2\|_\infty = \mu\|\Delta_1 - \Delta_2\|_\infty$.
$\mathcal{P}$ maps $X$ into $X$: $\Delta^{\mathrm{lin}}$ is gradient flow on $\frac12\|E_0 + \Psi_0\Delta\|^2$, which obeys the energy identity with $K_0$, so $|\Delta^{\mathrm{lin}}(t)|\le\sqrt{t/2}\,\|E_0\|\le r/2$ by Proposition~\ref{prop:ell}, and $|N(\Delta)|\le\Lambda_N(r)|\Delta|\le\Lambda_N(r)r$ on the ball since $N(0) = 0$, so the integral term is at most $\mu r\le r/2$.
By the contraction principle $\mathcal{P}$ has a unique fixed point in $X$ and the iterates converge to it at the stated rate.
The true path lies in $X$, because $|\Delta(t)|\le\ell(t)\le\sqrt{t/2}\,\|E_0\|\le r/2$ by Proposition~\ref{prop:ell}, and it is a fixed point, so it is the limit.
Finally $\|\Net(\theta_0 + \Delta_1) - \Net(\theta_0 + \Delta_2)\|\le\sup_{B_r}\lambda_1(\theta)^{1/2}|\Delta_1 - \Delta_2|$, and $\lambda_1(\theta)^{1/2} = \sup_{\|\varphi\|\le1}\|\varphi\|_{K(\theta)}\le\lambda_1(0)^{1/2} + \sigma(r)$ on the ball by the proof of Lemma~\ref{lem:drift}(d) applied to the segment from $\theta_0$ to $\theta$.
(b) $N$ is Lipschitz with constant $\Lambda_N(r)$ on the convex ball of radius $r$, by the mean-value inequality, and $P_r$ is 1-Lipschitz with values in that ball, so $\tilde N = N\circ P_r$ is Lipschitz with constant $\Lambda_N(r)$ on all of parameter space and equals $N$ on the ball.
Let $\tilde{\mathcal{P}}$ be $\mathcal{P}$ with $\tilde N$ in place of $N$.
For any two continuous paths $a$ and $b$, $\|e^{-G_0\tau}\|_{op}\le1$ gives $|\tilde{\mathcal{P}}[a](t) - \tilde{\mathcal{P}}[b](t)|\le\Lambda_N(r)\int_0^t|a(s) - b(s)|\,ds$, and by induction on $k$, $|\tilde{\mathcal{P}}^k[a](t) - \tilde{\mathcal{P}}^k[b](t)|\le\frac{(\Lambda_N(r)t)^k}{k!}\sup_{s\le t}|a(s) - b(s)|$.
By Proposition~\ref{prop:ell}, $|\Delta(t)|\le\sqrt{t/2}\,\|E_0\|\le r$ for $t\le T$, so $\tilde N(\Delta) = N(\Delta)$ along the true path and $\Delta = \tilde{\mathcal{P}}[\Delta]$ on $[0,T]$.
With $a = \Delta^{\mathrm{lin}}$ and $b = \Delta$, and $|\Delta^{\mathrm{lin}}(s)|, |\Delta(s)|\le\sqrt{s/2}\,\|E_0\|$, this gives the first bound of Eq.~\eqref{eq:picardany}.
For the second, $|P_r\Delta_{[k]} - \Delta| = |P_r\Delta_{[k]} - P_r\Delta|\le|\Delta_{[k]} - \Delta|$, and the Lipschitz bound on the network over the ball is the one of (a).
For Eq.~\eqref{eq:LambdaNbound}: $\Hess - G_0 = S(E) + (G - G_0)$; $\|S(E;\theta)\|_{op}\le H_2(\theta)\|E(\theta)\|$ and $\|E(\theta)\|\le\|E_0\| + \sup_{B_r}\lambda_1^{1/2}\,|\theta - \theta_0|$; and for a unit parameter vector $u$, $u^\top(G - G_0)u = \|\Psi u\|^2 - \|\Psi_0u\|^2$ with $\big|\|\Psi u\| - \|\Psi_0u\|\big|\le\|(\Psi - \Psi_0)u\| = \sup_{\|\varphi\|\le1}(g_\varphi(\theta) - g_\varphi(\theta_0))\cdot u\le\sigma(r)$ and $\|\Psi u\| + \|\Psi_0u\|\le2\lambda_1(0)^{1/2} + \sigma(r)$.
\end{proof}

Theorem~\ref{thm:picard} is the framework for gradient flow on the squared loss; Section~\ref{sec:predict} states it for every update rule (Theorem~\ref{thm:predict}).
Every iterate is a function of the initial error, the initial features, and the network on a ball about $\theta_0$, which within its radius of convergence is the tower at $\theta_0$, and the limit is the true path, with equality.
The convergence is not one order of a series per iterate: $\Delta_{[1]}$ already contains every order of the Taylor series of $\Net$ inside $N$, and the $k$-th iterate is exact through the $k$-th power of a bookkeeping parameter multiplying $N$, the successive approximations of the integral equation in $N$.
The zeroth parameter iterate solves the gradient flow of the first-order expansion about $\theta_0$; the error $E_{[0]}$ of Eq.~\eqref{eq:picard} is the full network evaluated on that path, not the expansion, and every later iterate corrects the path with the whole of $N$.
The scheme is an expansion about $\theta_0$; it does not hold the kernel fixed, and its limit is the true path with the kernel moving as it does.
The quadratic model is the same scheme with the series of $\Net$ in Eq.~\eqref{eq:closureE} cut after its second-order term, so that $N$ is a polynomial of degree three in $\Delta$; it is a truncation, and this paper does not make it.
The rate is set by $\Lambda_N$, the curvature of the loss beyond its kernel term, and not by $\lambda_1$, because the linear dynamics have been summed into $e^{-G_0t}$ and $e^{-K_0t}$.
The condition $\mu\le\frac12$ of part (a) gives geometric convergence on a short window; it is not a horizon.
By part (b) the scheme converges from $\theta_0$ on every finite window, at the rate $(\Lambda_NT)^k/k!$, with the ball chosen to contain the path, $r\ge\sqrt{2T}\,\|E_0\|$, and $\Lambda_N$ evaluated at that radius and not at $r\to0$.
What grows with $T$ is the cost: by Stirling's formula, $(\Lambda_NT)^k/k!\le(e\Lambda_NT/k)^k$, so a fixed accuracy needs about $e\,\Lambda_N\big(\sqrt{2T}\,\|E_0\|\big)\,T$ iterates, and the iterates evaluate the network on a ball of radius $\sqrt{2T}\,\|E_0\|$.
By Eq.~\eqref{eq:LambdaNbound}, $\Lambda_N$ at that radius contains, besides $H_2\|E_0\|$, the drift of the kernel term over the ball, whose terms linear in $r$ have coefficients of the order of $\lambda_1(0)^{1/2}H_2$, and those terms are the larger ones when $H_2\|E_0\|\ll\lambda_1(0)$; no closed form for the number of iterates is claimed.
The scheme can also be restarted at an intermediate time from the state there, which lowers the number of iterates; the concatenation of the windows is exact.

\subsection{What each expansion gives over a training run}

Each of the three expansions is an equality in its limit.
None of them gives a bound that is uniform over all of training from the initial state, and the reasons are different for each.
\begin{enumerate}[label=(\roman*)]
\item The Taylor series in $\Delta$ has a finite radius of convergence for $\tanh$ networks: $\tanh$ has poles at $\pm i\pi/2$ in its argument, so the series of $\Net$ in $\theta$ converges only within a distance of $\theta_0$ set by the pre-activations, and nothing keeps the path of a rough target, whose length $\ell$ grows while its slow eigenvalues grow by large factors, inside that ball.
Where it leaves it, the resummed forms of Lemma~\ref{lem:drift}(c) remain exact, but their coefficients are maxima over the ball, which is no longer a prediction from $\theta_0$.
\item The series in time, Eq.~\eqref{eq:timeseries}, resolves the fast directions first, and already its terms $(-K_0)^nE_0$ leave a truncation error at order $p$ of $x^{p+1}/(p+1)!$ with $x = \lambda_1t$ in the fastest direction; it is the expansion to avoid.
\item The Dyson series converges for every $t$, but only with $\delta K(s)$ taken along the path; from the initial state alone it reduces to (i).
\item The Picard scheme converges from the initial state on every finite window (Theorem~\ref{thm:picard}(b)).
What it does not give is an estimate uniform in time: the number of iterates needed for a given accuracy grows like $\Lambda_N(\sqrt{2T}\,\|E_0\|)\,T$, and the iterates evaluate the network on a ball of radius $\sqrt{2T}\,\|E_0\|$, which beyond the radius of item (i) is no longer the tower at $\theta_0$.
No finite iterate is a closed-form prediction of late-time quantities.
\end{enumerate}
The questions that the one-time theory cannot answer, whether a slow component decays, the sign of $\kappa_S$ late in training, the growth of the small eigenvalues, are determined by the limit of the Picard scheme, but no finite iterate with a bound independent of $T$ settles them.
That is why going one order higher, or any fixed number of orders higher, does not settle them, and why Section~\ref{sec:closing} states them as assumptions on the path rather than as terms to compute.

% ===============================================================
\section{Predicting the error along training: the general framework}\label{sec:predict}
% ===============================================================

This section states the framework in the generality in which it is meant to be used.
Any design, meaning any network, labels, known part, output scale, loss or initialization, trained by any update rule that is a continuously differentiable flow or map on a finite-dimensional state, has an error $E(t)$ that is predicted from the initial state by an explicit convergent scheme, with an error bound that is computed along with the prediction, and the prediction is certified to describe training whenever a check computed from it holds.
Nothing is assumed about the architecture or the problem.
This section lifts the conventions $\Pi = I$ and $c_p = 1$ of Part~\ref{part:path}: the update rule is arbitrary.
The Picard scheme of Theorem~\ref{thm:picard} is the instance for gradient flow on the squared loss with the kernel term as reference; the comparison of designs in Part~\ref{part:apps} is an application of Corollary~\ref{cor:compare} below.

\subsection{Training as a dynamical system}

Let the state $z\in\R^{d_z}$ collect the parameters $\theta$ and the state $\varsigma$ of the optimizer, $z = (\theta,\varsigma)$, with the Euclidean norm $|\cdot|$.
Training is either a flow or a map:
\begin{equation}\label{eq:trainingsystem}
\dot z = \mathcal{F}(t, z),\quad z(0) = z_0,\qquad\text{or}\qquad z_{j+1} = M_j(z_j),\quad j = 0, 1, \dots
\end{equation}
The flow $\mathcal{F}$ is assumed continuous in $(t, z)$ and continuously differentiable in $z$, with $D_z\mathcal{F}$ continuous in $(t, z)$; each map $M_j$ is assumed continuously differentiable on the ball about $z_0$ considered below.
The error is a function of the state, $E(z) = c_p\Net(\theta) - Y$ at the training inputs, and so is every other quantity of the theory: the loss, the kernel, the slow part, the Krylov values.
The following are instances.
\begin{itemize}
\item \emph{Gradient flow} on any twice continuously differentiable loss: $z = \theta$, $\mathcal{F} = -\nabla_\theta L$.
\item \emph{Gradient descent} with step sizes $\eta_j$, which may follow a schedule: $M_j(\theta) = \theta - \eta_j\nabla_\theta L(\theta)$; with a fixed preconditioner, $\eta_j$ is replaced by $\eta_j\Pi$.
\item \emph{Minibatch training} with a given sequence of batches $\mathcal{B}_j$ and step sizes, as for one run with a fixed seed: $M_j(\theta) = \theta - \eta_j\nabla_\theta L_{\mathcal{B}_j}(\theta)$, with $L_{\mathcal{B}_j}$ the loss on the batch $\mathcal{B}_j$; the maps differ from step to step, which the framework allows.
\item \emph{Momentum}: $z = (\theta, v)$, $v_{j+1} = \beta_{\mathrm m}v_j - \eta_j\nabla_\theta L(\theta_j)$, $\theta_{j+1} = \theta_j + v_{j+1}$, or its continuous-time form $\dot\theta = v$, $\dot v = -\beta_{\mathrm m}v - \nabla_\theta L(\theta)$.
\item \emph{Adam}: $z = (\theta, m, v)$ with the first and second moment estimates.
The second moments start at $v = 0$, where the square root in the update is not differentiable, so standard Adam is not covered as it stands.
The framework applies to the variant in which the square root is replaced by a function that is continuously differentiable for all real arguments, such as $(v^2 + \epsilon_{\mathrm A}^4)^{1/4}$, which differs from $\sqrt{|v|}$ by at most $\epsilon_{\mathrm A}$; the prediction is then for that variant.
\end{itemize}

\paragraph{Reference and remainder.}
Choose a reference linear part: a family $\mathcal{J}(t)$ of matrices, continuous in $t$, for the flow, or matrices $\mathcal{J}_j$ for the map.
The default is the linearization at the initial state, $\mathcal{J}(t) = D_z\mathcal{F}(t, z_0)$ or $\mathcal{J}_j = DM_j(z_0)$; for gradient flow on the squared loss it is $-\Hess(\theta_0) = -\big(G_0 + S(E_0;\theta_0)\big)$.
Theorem~\ref{thm:picard} uses instead the kernel term, $\mathcal{J} = -G_0$, which keeps the propagator a contraction, at the price of a remainder whose Lipschitz constant does not vanish as $r\to0$ but tends to $\|S(E_0;\theta_0)\|_{op}$.
Let $\mathcal{U}(t,s)$ be the propagator of $\dot x = \mathcal{J}(t)x$, and $\mathcal{U}(j,i) = \mathcal{J}_{j-1}\cdots\mathcal{J}_i$ with $\mathcal{U}(i,i) = I$, and let $\Gamma\ge1$ and real $\xi$, or $\zeta > 0$, satisfy
\begin{equation}\label{eq:propbound}
\|\mathcal{U}(t,s)\|_{op}\le\Gamma\,e^{\xi(t-s)}\quad(s\le t),\qquad\text{or}\qquad \|\mathcal{U}(j,i)\|_{op}\le\Gamma\,\zeta^{\,j-i}\quad(i\le j).
\end{equation}
Write the displacement $\delta = z - z_0$; for gradient flow it is $\Delta$ of Section~\ref{sec:allorders}.
The flow is $\dot\delta = \mathcal{F}(t,z_0) + \mathcal{J}(t)\delta + N_t(\delta)$ with remainder $N_t(\delta) = \mathcal{F}(t, z_0 + \delta) - \mathcal{F}(t, z_0) - \mathcal{J}(t)\delta$, and the map is $\delta_{j+1} = c_j + \mathcal{J}_j\delta_j + N_j(\delta_j)$ with $c_j = M_j(z_0) - z_0$ and $N_j(\delta) = M_j(z_0 + \delta) - z_0 - c_j - \mathcal{J}_j\delta$; in both cases $N(0) = 0$.
For gradient flow with $\mathcal{J} = -G_0$, $N_t = -N$ with $N$ of Eq.~\eqref{eq:Nsplit}.
For $r > 0$, the Lipschitz constant of the remainder on the ball $|\delta|\le r$ over the window considered is
\begin{equation}\label{eq:LambdaF}
\Lambda_F(r) = \sup_{t\le T,\ |\delta|\le r}\|D_z\mathcal{F}(t, z_0 + \delta) - \mathcal{J}(t)\|_{op}\qquad\text{or}\qquad \Lambda_F(r) = \sup_{j<J,\ |\delta|\le r}\|DM_j(z_0 + \delta) - \mathcal{J}_j\|_{op},
\end{equation}
by the mean-value inequality on the convex ball.
With the default reference, $\Lambda_F(r)\to0$ as $r\to0$.
Let $P_r$ be the nearest-point projection onto the ball of radius $r$.

\paragraph{The scheme.}
The zeroth iterate is the reference path, $\delta_{[0]}(t) = \int_0^t\mathcal{U}(t,s)\mathcal{F}(s,z_0)\,ds$, or $\delta_{[0],j} = \sum_{i<j}\mathcal{U}(j,i+1)c_i$, and
\begin{equation}\label{eq:genpicard}
\delta_{[k+1]}(t) = \delta_{[0]}(t) + \int_0^t\mathcal{U}(t,s)\,N_s\big(P_r\delta_{[k]}(s)\big)\,ds,\qquad
\delta_{[k+1],j} = \delta_{[0],j} + \sum_{i<j}\mathcal{U}(j,i+1)\,N_i\big(P_r\delta_{[k],i}\big).
\end{equation}
The predicted state is $\hat z_{[k]} = z_0 + P_r\delta_{[k]}$, and the predicted error is $E(\hat z_{[k]})$.
Every iterate is computed from the initial state, the update rule and the network on the ball of radius $r$, without running training.

\subsection{The prediction theorem}

\begin{theorem}[Prediction from the initial state]\label{thm:predict}
Let $T > 0$, or a number of steps $J$, and $r > 0$, and write $\xi^+ = \max(\xi, 0)$, $\zeta^+ = \max(\zeta, 1)$ and $\Lambda = \Lambda_F(r)$.
\begin{enumerate}[label=(\alph*)]
\item \emph{Convergence.} The iterates converge uniformly on $[0,T]$, or on $j\le J$, to the unique solution $\delta^\ast$ of the clipped equation, Eq.~\eqref{eq:genpicard} with $\delta_{[k]} = \delta_{[k+1]} = \delta^\ast$, and
\begin{equation}\label{eq:genbound}
\begin{aligned}
|\delta_{[k]}(t) - \delta^\ast(t)|\ &\le\ \varepsilon_k(t) := e^{\xi^+t}\,\frac{(\Gamma\Lambda t)^k}{k!}\,e^{\Gamma\Lambda t}\,\sup_{s\le t}|\delta_{[1]}(s) - \delta_{[0]}(s)|,\\
|\delta_{[k],j} - \delta^\ast_j|\ &\le\ \varepsilon_{k,j} := (\zeta^+)^j\binom jk\Big(\frac{\Gamma\Lambda}{\zeta^+}\Big)^k\Big(1 + \frac{\Gamma\Lambda}{\zeta^+}\Big)^{j-k}\sup_{i\le j}|\delta_{[1],i} - \delta_{[0],i}|,
\end{aligned}
\end{equation}
with $\binom jk = 0$ for $k > j$; for the map, moreover, $\delta_{[k],j} = \delta^\ast_j$ exactly for $j\le k + 1$.
\item \emph{Certification.} If $|\delta_{[k]}(t)| + \varepsilon_k(t)\le r$ for every $t\le T$, or $|\delta_{[k],j}| + \varepsilon_{k,j}\le r$ for every $j\le J$, then the training trajectory stays in the ball of radius $r$ about $z_0$, equals $z_0 + \delta^\ast$, and
\begin{equation}\label{eq:gencert}
|z(t) - z_0 - \delta_{[k]}(t)|\le\varepsilon_k(t),\qquad |z_j - z_0 - \delta_{[k],j}|\le\varepsilon_{k,j}.
\end{equation}
If the trajectory is known by another argument to stay in the ball, Eq.~\eqref{eq:gencert} holds without the check.
\item \emph{Readout.} Under (b), for every quantity $Q$ of the state that is Lipschitz with constant $\mathrm{Lip}_r(Q)$ on the ball, $|Q(\hat z_{[k]}(t)) - Q(z(t))|\le\mathrm{Lip}_r(Q)\,\varepsilon_k(t)$, and the same for the map.
For the error, $\mathrm{Lip}_r(E)\le|c_p|\sup_{|\theta-\theta_0|\le r}\lambda_1(\theta)^{1/2}$, with the kernel of the features with $\Pi = I$.
\item \emph{Restarting.} Two solutions of the clipped flow that start at displacements $a$ and $b$ at time $t_1$ differ at time $t\ge t_1$ by at most $\Gamma\,e^{(\xi^+ + \Gamma\Lambda)(t - t_1)}|a - b|$, and two solutions of the clipped map that start at step $j_1$ differ at step $j$ by at most $\Gamma\,(\zeta^+ + \Gamma\Lambda)^{j - j_1}|a - b|$.
So the scheme can be restarted from a predicted state, with the error of that state carried forward by this factor.
\end{enumerate}
\end{theorem}

\begin{proof}
Let $\tilde N = N\circ P_r$.
$N$ is Lipschitz with constant $\Lambda$ on the convex ball, by the mean-value inequality, and $P_r$ is 1-Lipschitz with values in the ball, so $\tilde N$ is Lipschitz with constant $\Lambda$ everywhere.

(a) For the flow, let $\tilde{\mathcal P}$ be the map on continuous paths given by the right side of Eq.~\eqref{eq:genpicard}.
By Eq.~\eqref{eq:propbound}, $|\tilde{\mathcal P}[a](t) - \tilde{\mathcal P}[b](t)|\le\Gamma\Lambda\int_0^te^{\xi(t-s)}|a(s) - b(s)|\,ds$.
Let $d_m = \delta_{[m+1]} - \delta_{[m]}$ and $w_m(t) = e^{-\xi^+t}|d_m(t)|$.
Since $e^{\xi(t-s)}\le e^{\xi^+(t-s)}$, $w_m(t)\le\Gamma\Lambda\int_0^tw_{m-1}(s)\,ds$, and by induction $w_m(t)\le\frac{(\Gamma\Lambda t)^m}{m!}\sup_{s\le t}w_0(s)\le\frac{(\Gamma\Lambda t)^m}{m!}\sup_{s\le t}|d_0(s)|$.
Summing over $m\ge k$ with $\sum_{m\ge k}x^m/m!\le\frac{x^k}{k!}e^x$ shows that the iterates are uniformly Cauchy on $[0,T]$; their uniform limit $\delta^\ast$ is continuous, is a fixed point of $\tilde{\mathcal P}$ because $\tilde{\mathcal P}$ is continuous in the supremum norm, and satisfies the first bound of Eq.~\eqref{eq:genbound}.
Two fixed points $a$, $b$ have $D(t) = e^{-\xi^+t}|a(t) - b(t)|\le\Gamma\Lambda\int_0^tD$, so $D = 0$ by Gr\"onwall's inequality.
For the map, the same computation with sums gives $w_{m,j}\le\frac{\Gamma\Lambda}{\zeta^+}\sum_{i<j}w_{m-1,i}$ for $w_{m,j} = (\zeta^+)^{-j}|d_{m,j}|$, hence $w_{m,j}\le\binom jm(\Gamma\Lambda/\zeta^+)^m\sup_{i\le j}w_{0,i}$, and the tail bound $\sum_{m\ge k}\binom jmx^m\le\binom jkx^k(1+x)^{j-k}$, from $\binom jm\le\binom jk\binom{j-k}{m-k}$, gives the second bound; the fixed point of the clipped recursion is unique because the recursion determines it step by step.
The map is strictly causal: $\delta_{[k+1],j}$ depends on $\delta_{[k],i}$ only for $i < j$.
Since $\delta_{[0],0} = 0 = \delta^\ast_0$ and $\delta_{[0],1} = c_0 = \delta^\ast_1$, because $N_0(0) = 0$, induction on $k$ gives $\delta_{[k],j} = \delta^\ast_j$ for $j\le k + 1$.

(b) Under the hypothesis, $|\delta^\ast(t)|\le|\delta_{[k]}(t)| + \varepsilon_k(t)\le r$, so $P_r\delta^\ast = \delta^\ast$ and $\delta^\ast$ solves the integral equation with $N$ in place of $\tilde N$.
For the flow, that integral equation is the variation-of-constants form of $\dot\delta = \mathcal{F}(t, z_0 + \delta)$, so $z_0 + \delta^\ast$ solves the flow with $z(0) = z_0$, and by uniqueness of solutions for a right side that is continuous in $(t, z)$ and locally Lipschitz in $z$ it is the training trajectory.
For the map, $\delta^\ast_{j+1} = c_j + \mathcal{J}_j\delta^\ast_j + N_j(\delta^\ast_j)$ for every $j < J$, which is the training recursion, and $\delta^\ast_0 = 0$.
The bound is then (a).
If the trajectory is known to stay in the ball, it solves the clipped equation, so it is $\delta^\ast$ by the uniqueness in (a).

(c) Both $\hat z_{[k]}$ and $z$ lie in the ball, and $|P_r\delta_{[k]} - \delta^\ast| = |P_r\delta_{[k]} - P_r\delta^\ast|\le|\delta_{[k]} - \delta^\ast|$.
The parameters of two states in the ball differ by at most the distance between the states.
For the error, $c_p\Net$ is Lipschitz in $\theta$ on the ball with constant $|c_p|$ times the largest operator norm of its Jacobian there, which is $|c_p|\sup\lambda_1^{1/2}$ for the RMS norm, with the kernel of the features with $\Pi = I$.

(d) For the flow, $D(t) = |a(t) - b(t)|$ satisfies $D(t)\le\Gamma e^{\xi^+(t-t_1)}|a - b| + \Gamma\Lambda\int_{t_1}^te^{\xi^+(t-s)}D(s)\,ds$, and Gr\"onwall's inequality applied to $e^{-\xi^+(t-t_1)}D$ gives the bound.
For the map, $D_j\le\Gamma(\zeta^+)^{j-j_1}|a - b| + \Gamma\Lambda\sum_{j_1\le i<j}(\zeta^+)^{j-1-i}D_i$, and induction gives $D_j\le\Gamma(\zeta^+ + \Gamma\Lambda)^{j-j_1}|a - b|$.
\end{proof}

Theorem~\ref{thm:predict} is the general form of the framework.
Its hypotheses are smoothness of the update rule on a ball about the initial state and a bound on the reference propagator; the architecture, the loss and the labels enter only through $\mathcal{F}$ or $M_j$, and through the constants $\Gamma$, $\xi$ or $\zeta$, and $\Lambda_F$.
Part (a) says that the iterates converge, to the solution of the clipped equation, for every window, with no horizon.
Part (b) says when that solution is the training trajectory, with a check that is computed from the iterate itself.
For gradient flow on a loss $L\ge0$ the check is automatic once $r\ge\sqrt{T\,L(\theta_0)}$, since the Cauchy--Schwarz inequality and $\dot L = -|\dot\theta|^2$ bound the path length up to $T$ by $\sqrt{T\,(L(\theta_0) - L(\theta(T)))}$; on the squared loss with $\Pi = I$ and $c_p = 1$ this is $\sqrt{T/2}\,\|E_0\|$ (Proposition~\ref{prop:ell}).
Part (c) turns the prediction of the state into a prediction of the error, the loss, or any other Lipschitz quantity, with a proved error.
Part (d) allows the window to be split.

\paragraph{The instances.}
\begin{itemize}
\item \emph{Gradient flow on the squared loss.} With the kernel term as reference, $\mathcal{J} = -G_0$, $\Gamma = 1$, $\xi = 0$ and $\Lambda_F = \Lambda_N$ of Eq.~\eqref{eq:LambdaN}; the iterates are those of Theorem~\ref{thm:picard}(b), with $\delta = \Delta$ and $N_t = -N$, and the bound there replaces $\sup_s|\delta_{[1]}(s) - \delta_{[0]}(s)|\,e^{\Lambda t}$ by the a priori $\sqrt{2t}\,\|E_0\|$.
The zeroth iterate is $\delta_{[0]}(t) = -G_0^{+}\big(I - e^{-G_0t}\big)g_{E_0}$, with $G_0^+$ the pseudo-inverse, since $g_{E_0}$ lies in the range of $G_0$.
With the default reference $-\Hess(\theta_0)$, which is symmetric, $\Gamma = 1$ and $\xi = \max\big(0, -\lambda_{\min}(\Hess(\theta_0))\big)\le\|S(E_0;\theta_0)\|_{op}$, and $\Lambda_F(r)\to0$ as $r\to0$.
\item \emph{Gradient descent.} With the reference $\mathcal{J}_j = I - \eta_jG_0$, $\Lambda_F = \max_j\eta_j\,\Lambda_N$.
The $\mathcal{J}_j$ commute, and $|1 - \eta_j\lambda|$ over $\lambda\in[0,\lambda_1(0)]$ is largest at an endpoint, so $\Gamma = 1$ and $\zeta = \max\big(1, \max_j|1 - \eta_j\lambda_1(0)|\big)$; if $\eta_j\lambda_1(0)\le2$ for every $j$, then $\zeta = 1$.
Only the kernel at the start enters the reference: sharpening along the path, including the edge of stability, is inside $\Lambda_F$, at the price of the factor $(\zeta^+)^j$ when $\eta_j\lambda_1(0) > 2$.
\item \emph{Minibatches.} With the reference $\mathcal{J}_j = I - \eta_jG_0^{(\mathcal{B}_j)}$, with $G_0^{(\mathcal{B}_j)}$ the kernel term of the batch loss at $\theta_0$, submultiplicativity gives $\Gamma = 1$ and $\zeta = \max_j\max\big(1, |1 - \eta_j\lambda_1^{(\mathcal{B}_j)}(0)|\big)$, with $\lambda_1^{(\mathcal{B}_j)}(0)$ the top eigenvalue of $G_0^{(\mathcal{B}_j)}$, and $\Lambda_F$ is the largest Lipschitz constant of the batch remainders.
The prediction is for the given sequence of batches.
\item \emph{Heavy-ball momentum} with a constant step size $\eta$, and the reference built on $G_0$: in the eigenbasis of $G_0$ the reference splits into the $2\times2$ blocks $\begin{psmallmatrix}1 - \eta\lambda & \beta_{\mathrm m}\\ -\eta\lambda & \beta_{\mathrm m}\end{psmallmatrix}$, one for each eigenvalue $\lambda$, and $\Gamma$ and $\zeta$ must bound the powers of these blocks uniformly over $\lambda\in[0,\lambda_1(0)]$.
The blocks are not normal, so $\Gamma > 1$ in general, and $\Gamma$ grows as $\zeta$ approaches the largest spectral radius of the blocks; with varying step sizes, $\zeta$ must bound the products of the blocks, not only their spectral radii.
\item \emph{Momentum and Adam in general.} Any reference whose propagator is bounded as in Eq.~\eqref{eq:propbound} can be used, for instance the Jacobian of the update map at $z_0$, with $\Gamma$ and $\zeta$ read off its products.
The energy identity is not available, and the containment is certified by Theorem~\ref{thm:predict}(b).
\end{itemize}

\subsection{What is predicted, and at what cost}

The prediction is two-sided.
Under the check of Theorem~\ref{thm:predict}(b), at every time in the window the error of the trained network lies within $\mathrm{Lip}_r(E)\,\varepsilon_k(t)$ of the error of the predicted state, and every Lipschitz function of the state, such as the loss, lies within the corresponding band.
The quantities of Part~\ref{part:onetime} at the predicted state, the slow part, the floor and the Krylov values, are functions of the error and of the features, and inherit bands from Lipschitz bounds on those functions over the ball.

The accuracy is set by $X = \Gamma\Lambda_F T$.
By Stirling's formula, $X^k/k!\le(eX/k)^k$, and with the factor $e^X$ in $\varepsilon_k$, a fixed accuracy needs a number of iterates proportional to $X$, about $3.6\,X$ plus a term logarithmic in the accuracy; for the map, the governing quantity is $J\,\Gamma\Lambda_F/\zeta^+$, together with the factor $(\zeta^+)^J$.
Each iterate costs one evaluation of the update rule along a path over the window, which is about the cost of training over that window.
For the map, the scheme is exact at step $j$ after $k\ge j - 1$ iterates: carried to the end, at a cost of order $J^2$ steps, it reproduces the clipped recursion, which is training whenever the trajectory stays in the ball.
The scheme is therefore a convergent, certified expansion of training from the initial state, and not, in general, a cheaper way to compute it.
It is cheap when $X$ is small, which with the default reference means that the update rule stays close to its linearization over the ball that the window covers: then the zeroth and first iterates, which are explicit or need one pass along an explicit path, already carry a small certified error.
How small $X$ is along real training, and whether a reference that moves with the iterate reduces $\Lambda_F$ to the part of the curvature that the error drives, are the open questions of the framework (Section~\ref{sec:measure}).

\subsection{Comparing designs}

A comparison of two designs needs an upper bound on the error of one and a lower bound on the error of the other, at the same time.
The framework supplies both, and so do the lower bounds of Parts~\ref{part:onetime} and~\ref{part:path}.

\begin{corollary}[Comparing two designs]\label{cor:compare}
Let designs A and B be trained by update rules that satisfy the hypotheses of Theorem~\ref{thm:predict}, and let $T$ be a time in both windows.
Let $\overline{E}_A(T)$ be any proved upper bound on $\|E_A(T)\|$ and $\underline{E}_B(T)$ any proved lower bound on $\|E_B(T)\|$, for instance
\begin{itemize}
\item $\overline{E}_A(T) = \|E_A(\hat z^A_{[k]}(T))\| + \mathrm{Lip}_r(E_A)\,\varepsilon^A_k(T)$ and $\underline{E}_B(T) = \|E_B(\hat z^B_{[k]}(T))\| - \mathrm{Lip}_r(E_B)\,\varepsilon^B_k(T)$, from Theorem~\ref{thm:predict}, when the check of Theorem~\ref{thm:predict}(b) holds for each design;
\item $\overline{E}_A(T) = \|E_A(0)\|$ under preconditioned gradient flow on the squared loss with a fixed step-size matrix (Lemma~\ref{lem:monoexp});
\item for B trained with $\Pi = I$ and $c_p = 1$, and $\ell_B$ an upper bound on its path length up to $T$, such as $\sqrt{T/2}\,\|E_B(0)\|$ under gradient flow on the squared loss: $\underline{E}_B(T)$ from the floor at time 0 or from the roughness that a path can supply (Theorem~\ref{thm:time0}, Corollary~\ref{cor:mintime}(a)), as in Theorem~\ref{thm:designcompare}(ii) and (iii); or $\|E_B(0)\|\exp(-\int_0^T\lambda_1^B)$ under preconditioned gradient flow on the squared loss (Lemma~\ref{lem:monoexp}).
\end{itemize}
If $\overline{E}_A(T) < \underline{E}_B(T)$, then $\|E_A(T)\| < \|E_B(T)\|$.
\end{corollary}

\begin{proof}
$\|E_A(T)\|\le\overline{E}_A(T) < \underline{E}_B(T)\le\|E_B(T)\|$.
\end{proof}

The corollary is the general principle; the bounds it combines are what make it informative.
The predicted bands of Theorem~\ref{thm:predict} give a sharp comparison at any time, at the cost of computing iterates; the cheapest bounds, an initial error that gradient flow cannot increase and a lower bound that holds at every time, give a comparison in closed form, valid for every time up to $T$ at once.
Section~\ref{sec:designcompare} develops the closed-form case and applies it to learned time steppers.

% ===============================================================
\section{The loss curve through the theory}\label{sec:loss}
% ===============================================================

Theorem~\ref{thm:identities}(c) writes the loss as a functional of $\kappa = \kappa_G + \kappa_S$.
This section reads the two parts through the objects of Part~\ref{part:onetime}.

\subsection{The kernel term is set by where the error lies}

\begin{proposition}[$\kappa_G$ and the slow part]\label{prop:kappaG}
At any time, for every $\gamma > 0$,
\begin{equation}\label{eq:KEsplit}
\|KE\|\ \le\ \gamma\,\|E\| + \lambda_1\,\|E - E_\gamma\|,\qquad \|E - E_\gamma\|\le\lambda_1^{1/2}\,\Upsilon(\gamma)^{1/2},
\end{equation}
and, if $\|E\|_K > 0$,
\begin{equation}\label{eq:kappaGsplit}
\kappa_G\ \le\ \gamma + \lambda_1\,\frac{\sum_{\lambda_k>\gamma}\lambda_k\hat e_k^2}{\|E\|_K^2};
\end{equation}
hence, the second line under gradient flow,
\begin{equation}\label{eq:rateslow}
\begin{aligned}
\kappa_G\,\|E\|_K^2 = \|KE\|^2\ &\le\ \big(\gamma\,\|E\| + \lambda_1^{3/2}\Upsilon(\gamma)^{1/2}\big)^2,\\
-\frac{d}{dt}\,\tfrac12\|E\|^2 = \|E\|_K^2\ &\le\ \|E\|\big(\gamma\,\|E\| + \lambda_1^{3/2}\Upsilon(\gamma)^{1/2}\big).
\end{aligned}
\end{equation}
\end{proposition}

\begin{proof}
$KE = KE_\gamma + K(E - E_\gamma)$; in the eigenbasis $KE_\gamma$ has components $\lambda_k\gamma\hat e_k/(\lambda_k + \gamma)\le\gamma\hat e_k$ in absolute value, and $\|K(E - E_\gamma)\|\le\lambda_1\|E - E_\gamma\|$; $\|E - E_\gamma\|^2 = \sum_k\lambda_k^2\hat e_k^2/(\lambda_k + \gamma)^2\le\lambda_1\Upsilon(\gamma)$.
For Eq.~\eqref{eq:kappaGsplit}, split $\kappa_G\|E\|_K^2 = \sum_k\lambda_k\cdot\lambda_k\hat e_k^2$ at $\lambda_k = \gamma$, and bound $\lambda_k$ by $\gamma$ below the split and by $\lambda_1$ above it.
For the rate, $-\frac{d}{dt}\frac12\|E\|^2 = \|E\|_K^2$ under gradient flow, and $\|E\|_K^2 = \ip{E}{KE}\le\|E\|\,\|KE\|$.
\end{proof}

In words: the loss decreases at a relative rate at most $2\gamma$, plus a term that is the part of the error the kernel removes quickly, measured by the length $\Upsilon(\gamma)^{1/2}$ of the step that removes it.
Once the fast part has been removed, which the floor of Part~\ref{part:onetime} says happens while the slow part stays, the loss creeps down at the rate of the slow directions, and by Eq.~\eqref{eq:speed} the speed decays at the rate $\kappa = \kappa_G + \kappa_S$, with $|\kappa_S|\le H_2\|E\|$ and, by Eq.~\eqref{eq:kappaGsplit}, $\kappa_G$ of the order of the slow eigenvalues once the kernel energy $\|E\|_K^2$, and not only the error, lies on eigenvalues below $\gamma$.
This is the floor in dynamic form: not a bound on the error at time $t$, but a bound on how fast it can be falling, at every time, in terms of the kernel of that time.
For a rough dataset Theorem~\ref{thm:floor} bounds $\|E_\gamma\|$ from below by the absolute roughness of the labels, so the error that remains is the error that falls slowly.

\subsection{The feature-learning part is one scalar-valued function of time}

At each instant, $\kappa_S$ is the whole departure of $\ddot L$ from the value that the kernel of that instant would give; the motion of the features up to that instant is already inside $K(t)$ and $\kappa_G$.
Three exact facts about it:
\begin{enumerate}[label=(\roman*)]
\item $\kappa_S\|E\|_K^2 = \ip{E}{\nabla^2_\theta\Net[g_E,g_E]}$: the component along the error of the curvature of the output along the gradient.
It is the same sampled function that appears in the acceleration of the error, Eq.~\eqref{eq:accel}, and, for one step, $\ip{E_j}{r_j} = \frac{\eta^2}{2}\kappa_S\|E_j\|_K^2 + \varrho_j$ with, exactly, $\varrho_j = \eta^2\int_0^1(1 - \tau)\ip{E_j}{\big(\nabla^2_\theta\Net(\theta_j - \tau\eta g_j) - \nabla^2_\theta\Net(\theta_j)\big)[g_j,g_j]}\,d\tau$, which under the convergence hypothesis of Eq.~\eqref{eq:rseries} is the component along $E_j$ of the terms of order $\eta^3$ and higher of that series; so $\kappa_S > 0$ means that the second-order part of a gradient step points along the error, and $\kappa_S < 0$ that it points against it.
\item $|\kappa_S|\le H_2(\theta)\|E\|$, so by Eq.~\eqref{eq:speed},
\begin{equation}\label{eq:sandwich}
\exp\Big(-2\int_0^t\big(\kappa_G + H_2\|E\|\big)ds\Big)\ \le\ \frac{\|E(t)\|_K^2}{\|E_0\|_{K_0}^2}\ \le\ \exp\Big(-2\int_0^t\big(\kappa_G - H_2\|E\|\big)ds\Big),
\end{equation}
with $\kappa_G$ and $H_2$ evaluated along the path.
\item $\kappa_S$ is the only part of $\kappa$ that can be negative; the loss curve is convex in $t$ wherever $\kappa_S\ge-\kappa_G$, and convex on a logarithmic axis wherever $\kappa_S\ge-\mathrm{Var}_w(\lambda)/\mathrm{mean}_w(\lambda)$, Proposition~\ref{prop:whykappa}(b).
\end{enumerate}
The sign of $\kappa_S$ is not given by the theory.
It is an open sign of the same kind as that of $\ip{q}{E}$ for the top eigenvalue in Proposition~\ref{prop:lambda1}, but not the same one: both are contractions of the second derivative of the network, with the error in the output slot for $\kappa_S$, $\kappa_S\|E\|_K^2 = \frac1n\sum_iE_i\cdot\nabla^2_\theta\Net(x_i)[g_E,g_E]$, and in a parameter slot, through $g_E$, for $\ip{q}{E} = \frac1n\sum_iv_{1,i}\cdot\nabla^2_\theta\Net(x_i)[g_{v_1},g_E]$.

% ===============================================================
\section{Closing assumptions}\label{sec:closing}
% ===============================================================

The tower cannot be closed without a hypothesis, and this section states the hypotheses that close it, each as an assumption with its consequence proved from it.
None is proved here, and each is a quantity that can be measured along a run.

\begin{assumption}[Curvature along the gradient]\label{ass:curv}
There are constants $\kappa_-\le\kappa_+$ such that $\kappa_-\le\kappa(s)\le\kappa_+$ for all $s\in[t_0, t]$.
\end{assumption}

\begin{proposition}[Two-sided loss curve]\label{prop:curv}
Along gradient flow with $g_E(t_0)\neq0$, under Assumption~\ref{ass:curv}, with $\nu_0 = \|E(t_0)\|_{K(t_0)}^2$,
\begin{equation}\label{eq:twosided}
\nu_0\,e^{-2\kappa_+(t - t_0)}\ \le\ \|E(t)\|_{K(t)}^2\ \le\ \nu_0\,e^{-2\kappa_-(t - t_0)},
\end{equation}
\begin{equation}\label{eq:twosidedloss}
\begin{aligned}
\|E(t)\|^2\ &\le\ \|E(t_0)\|^2 - \nu_0\,\frac{1 - e^{-2\kappa_+(t - t_0)}}{\kappa_+},\\
\|E(t)\|^2\ &\ge\ \|E(t_0)\|^2 - \nu_0\,\frac{1 - e^{-2\kappa_-(t - t_0)}}{\kappa_-},
\end{aligned}
\end{equation}
with $(1 - e^{-2\kappa\tau})/\kappa$ read as $2\tau$ at $\kappa = 0$; and in differential form $-2\kappa_-L'\le L''\le-2\kappa_+L'$ for $L = \frac12\|E\|^2$, since $-L' = \|E\|_K^2\ge0$.
\end{proposition}

\begin{proof}
Eq.~\eqref{eq:speed} from $t_0$ with the bounds on $\kappa$, then Eq.~\eqref{eq:energy}; the differential form is Eq.~\eqref{eq:losscurve}.
The upper line of Eq.~\eqref{eq:twosidedloss} uses only $\kappa\le\kappa_+$, and the lower line only $\kappa\ge\kappa_-$.
\end{proof}

The upper bound in Eq.~\eqref{eq:twosidedloss} is the first statement of this theory that bounds the error from above beyond a short path: under a cap on the curvature along the gradient, $\|E\|^2$ must fall by at least $\nu_0\big(1 - e^{-2\kappa_+(t - t_0)}\big)/\kappa_+$, which tends to the kernel energy of the current error over the cap, from every time.
The cap is implied by a cap on the sharpness, since $\kappa\le\lambda_{\max}(\nabla^2_\theta L)$ and $\kappa_G\le\lambda_1$, so Assumption~\ref{ass:curv} with $\kappa_+ = \Lambda$ holds along any gradient-flow path whose sharpness stays below $\Lambda$; for gradient descent the discrete identity Eq.~\eqref{eq:speedstep} gives $|g_{j+1}|^2\ge(1 - 2\eta\bar\kappa_j)|g_j|^2$, so a cap on $\bar\kappa_j$ still bounds the speed from below, but the term $\eta^2|\bar\Hess_jg_j|^2$ is not controlled by the cap, and the one-step energy identity, Proposition~\ref{prop:onestep}(c), is not $-2\eta|g_j|^2$; so neither line of Eq.~\eqref{eq:twosidedloss} transfers to gradient descent as stated, which at $\eta\lambda_1 = 2$ is the known fact that the descent lemma says nothing.
The lower bound in Eq.~\eqref{eq:twosidedloss} needs $\kappa_-$, which for the kernel term is a lower bound on $\kappa_G$, that is, on the eigenvalues that carry the error; for a rough dataset the floor says those are small, and then, if $\kappa_S$ is also small, the lower bound is close to $\|E(t_0)\|^2 - 2\nu_0(t - t_0)$, the loss falling at most linearly at its current rate.

\begin{assumption}[Kernel-term dominance along the gradient]\label{ass:GN}
$\kappa_S(s)\ge0$ for all $s\in[t_0,t]$: the loss is at least as curved along the gradient as its kernel term, equivalently the second-order part of a gradient step has a nonnegative component along the error.
\end{assumption}

\begin{proposition}[Loss bounded below by kernel quantities]\label{prop:GN}
Along gradient flow with $g_E(t_0)\neq0$, under Assumption~\ref{ass:GN},
\begin{equation}\label{eq:GN}
\begin{aligned}
\|E(t)\|_{K(t)}^2\ &\le\ \|E(t_0)\|_{K(t_0)}^2\exp\Big(-2\int_{t_0}^t\kappa_G\,ds\Big),\\
\|E(t)\|^2\ &\ge\ \|E(t_0)\|^2 - 2\|E(t_0)\|_{K(t_0)}^2\int_{t_0}^t\exp\Big(-2\int_{t_0}^s\kappa_G\Big)ds ,
\end{aligned}
\end{equation}
with $\kappa_G$ the error-weighted mean eigenvalue of the kernel along the path.
\end{proposition}

\begin{proof}
Eq.~\eqref{eq:speed} with $\kappa\ge\kappa_G$, then Eq.~\eqref{eq:energy}.
\end{proof}

Under Assumption~\ref{ass:GN} feature learning can only slow the loss curve relative to the kernel's own prediction at each time; the loss is then bounded below by a quantity built from the kernels along the path and nothing else.
Whether $\kappa_S\ge0$ holds, and where it fails, is the measurement that decides whether this is the right closing hypothesis; in one step, $\ip{E_j}{r_j} - \varrho_j = \frac{\eta^2}{2}\kappa_S\|E_j\|_K^2$ exactly (Section~\ref{sec:loss}), so $\kappa_S$ itself is the quantity to measure, from one Hessian-vector product at each saved state of a run.

\begin{assumption}[Sharpening on average]\label{ass:sharp}
For gradient descent, let the centre $\theta_{c,j}$ of step $j$ be the midpoint of the iterates $\theta_j$ and $\theta_{j+1}$, and let $\alpha_j = -2\ip{q}{E}$ evaluated at $\theta_{c,j}$, with $q$ from Proposition~\ref{prop:lambda1} and $\lambda_1$ simple at each $\theta_{c,j}$.
There is a window length $m$ such that, along the stretch of training considered, the mean of $\alpha_j$ over any $m$ consecutive steps is positive.
\end{assumption}

By Proposition~\ref{prop:lambda1}, $-2\ip{q}{E}$ is the rate of growth of the top eigenvalue under gradient flow, so the assumption says that the top of the spectrum is pushed up on average.
Its consequence, the balance of self-stabilization at the edge of stability and the size of the oscillation, is a heuristic from the literature, not a result, and is stated with the edge-of-stability application (Part~\ref{part:apps}).
It closes the top of the spectrum; Assumptions~\ref{ass:curv} and~\ref{ass:GN} close the loss curve.
What none of them closes is the growth of the kernel on the slow directions, which is the one closure that would turn the floor into a prediction of when a rough target is learned; it would be a lower bound on $\ip{v_k}{\dot Kv_k} = -2g_{v_k}\cdot S(v_k)g_E$ for the slow $v_k$, and the theory has no sign for it.

% ===============================================================
\section{What to measure}\label{sec:measure}
% ===============================================================

Every quantity named in Part~\ref{part:path} is computable from a saved state or along a run, except the maxima over balls ($\bar H_m$, $\bar D_m$, $\Lambda_N$), which sampling can only estimate from below; none has been measured yet.
\begin{enumerate}[label=(\arabic*)]
\item $\kappa(t)$ from one Hessian-vector product $\Hess g_E$ per state, with its split into $\kappa_G = \|KE\|^2/\|E\|_K^2$, which needs one kernel product, and $\kappa_S$ as the difference; their signs and sizes along each run, the check of Eq.~\eqref{eq:speed} against the logged $\|E\|_K^2$, and $\rho = 2\|E\|_K^2/\|E\|^2$ against $2\kappa$, to find the stretches where Eq.~\eqref{eq:logconcave} holds and the loss falls faster than exponentially.
\item The path length $\ell_N = \eta\sum_j|g_j|$ along each run, against Eq.~\eqref{eq:ellbound} and against its exact form in Proposition~\ref{prop:ell}; both are proved for gradient flow, so this tests the flow prediction against the discrete run, not a proved bound on the run.
\item The derivative sizes $H_m(\theta)$ for $m = 2, 3, 4$ at saved states, by nested automatic differentiation, and $\bar H_5$ over the ball reached, so that the resummed bound Eq.~\eqref{eq:costresum} with $M = 4$ can be evaluated at the path lengths reached in the following 100, 1000 and 10000 steps, against $F_0(\ell)$ and against the measured error, to find $\ell_\ast$; the sum of the four measured terms alone is a truncation and is not used.
\item The direction-specific drifts $\sigma_k$ for the slow eigenvectors of a saved state, against $\sigma$, to see whether the second term of Eq.~\eqref{eq:Kdrift} is the one through which the slow eigenvalues grow.
\item The smallest $\kappa_-$ and largest $\kappa_+$ along each run, for Assumption~\ref{ass:curv}, and the fraction of states with $\kappa_S\ge0$, for Assumption~\ref{ass:GN}.
\item $\Lambda_N$ at saved states: the largest eigenvalue in absolute value of $\Hess - G$, by power iteration on Hessian-vector and kernel products, gives only $\Lambda_N(0^+)$, a lower bound on $\Lambda_N(r)$, and $\Lambda_N(r)$ needs the supremum over the ball, which sampling estimates from below; with it, the number of Picard iterates $e\,\Lambda_N(\sqrt{2T}\|E\|)\,T$ of Theorem~\ref{thm:picard}(b) for a window of $T$, in units of steps, and the window of part (a); and the first two Picard iterates from a saved state, computed for gradient flow, against the error of the discrete run 100 and 1000 steps later.
\item For the update rule actually used (gradient descent with its schedule, minibatches, Adam), the constants $\Gamma$, $\zeta$ and $\Lambda_F$ of Theorem~\ref{thm:predict} with the default reference and with the kernel-term reference, the quantity $\Gamma\Lambda_FT$ for windows of 100, 1000 and 10000 steps, the first iterates of Eq.~\eqref{eq:genpicard} against the run, and the check of Theorem~\ref{thm:predict}(b); and whether a reference that moves with the iterate lowers $\Lambda_F$.
\end{enumerate}

% ===============================================================
\section{Scope and status}\label{sec:status}
% ===============================================================

\paragraph{Part~\ref{part:onetime}.}
Theorems~\ref{thm:fast}, \ref{thm:floor}, \ref{thm:slow}, \ref{thm:stepsup} and~\ref{thm:stepkrylov} and Proposition~\ref{prop:rates} hold at an arbitrary time during training, for any labels and any update rule.
They need no model of the network.
Every statement uses the gradient features and the kernel at one time.
Theorem~\ref{thm:dataset} and Theorem~\ref{thm:slow}(b) use in addition the peak feature roughness over the set that one step can reach, and nothing else about how the kernel moves.
These results describe where the error lies in the eigenbasis of the kernel, the rate of change of the error, and what one step of a given length can do.

The one-step results have three readings.
(i) With $\Psi = \bar\Psi$ from Eq.~\eqref{eq:avgstep}, the change of the error lies in $\mathcal{D}_s$ with $s = |g|$, and every lower bound holds exactly for the step actually taken.
(ii) Proposition~\ref{prop:stepexact}, Theorem~\ref{thm:dataset} and Theorem~\ref{thm:slow}(b) hold exactly for every step of length at most $s$; they take the peak feature roughness over the set $B_s$ that the step can reach, and make no assumption on how the gradient features change.
(iii) With $\Psi = \Psi(t)$, Theorems~\ref{thm:stepsup} and~\ref{thm:stepkrylov}, Proposition~\ref{prop:krylovrate} and Proposition~\ref{prop:floorstep} are exact statements about the first-order Taylor expansion of the step about $\theta(t)$, and Proposition~\ref{prop:taylorstep} bounds the difference from the error that steps actually reach by $\tfrac12|c_p|\bar H_2(s)s^2$, with no assumption on how the gradient features change.
No statement in Part~\ref{part:onetime} assumes that the kernel stays fixed.
Lemma~\ref{lem:mech} holds under the stated assumptions.
The lower bounds hold on the training inputs, or on validation inputs with their own kernel.
The floor of Theorem~\ref{thm:floor} for a fixed seminorm is informative only when the leak term is small, that is, when $\sR^2 \ll \gamma$ or when the part of the error that the kernel removes quickly has already been removed.
The supremum over all seminorms, $F(s)$, has no leak term, but its optimal functional depends on the current error and kernel, not on the training inputs alone.
The absolute roughness of the dataset under a local complexity function depends on the data alone, and it enters the bounds only through the last term of Eqs.~\eqref{eq:main} and~\eqref{eq:mainstep}, which is exact when the output is a multiple $\alpha Y$ of the labels with $0 \le \alpha \le 1$ and loose when the roughness of the output is not aligned with that of the labels, for instance when the output is rougher than the labels.

\paragraph{Part~\ref{part:path}: exact, no hypothesis beyond smoothness.}
Proposition~\ref{prop:hessian} (the loss Hessian); Proposition~\ref{prop:tower} (the tower); Proposition~\ref{prop:lambda1} (rate of the top eigenvalue, for a simple eigenvalue); Proposition~\ref{prop:lossderiv} (every derivative of the loss as an error part plus an error-free part); Proposition~\ref{prop:timederiv} (every time derivative of the error and of the gradient by one derivation, and the hierarchy of the loss along the gradient); Theorem~\ref{thm:identities} (energy, speed, loss curve, split, acceleration; with $g_E(0)\neq0$); Proposition~\ref{prop:whykappa} (what $\kappa$ determines: the active eigenvalue, the shape of the loss curve, the bound by the sharpness, the exact rate of $\kappa$, the last with three derivatives); Proposition~\ref{prop:onestep} (the same for one step, with exact integral remainders); Proposition~\ref{prop:taylorstep} (the first-order step against the steps actually taken, from Part~\ref{part:onetime}); Proposition~\ref{prop:ell} (path length); Lemma~\ref{lem:drift} (kernel drift, integral and resummed forms); Theorem~\ref{thm:path} (floor along a path); Corollary~\ref{cor:mintime} (minimum path length for every update rule, minimum training time under gradient flow); Theorem~\ref{thm:duhamel} (exact expansion of the error about the kernel at the start, two-sided bound); Theorem~\ref{thm:dyson} (the Dyson series, convergent for every time); Theorem~\ref{thm:picard} (the Picard scheme from the initial state, convergent on every finite window, geometrically on windows set by the curvature beyond the kernel term); Proposition~\ref{prop:kappaG} ($\kappa_G$ and the slow part); Theorem~\ref{thm:predict} (prediction from the initial state for every update rule that is a continuously differentiable flow or map, with the certification of part (b)); Corollary~\ref{cor:compare} (comparison of designs).
The identities and bounds of Part~\ref{part:path} are proved for gradient flow, with the one-step forms of Proposition~\ref{prop:onestep} for gradient descent; Theorem~\ref{thm:path} holds for any update rule, since its proof uses only $\frac{d}{dt}\ip{\varphi}{\Net} = g_\varphi\cdot\dot\theta$.

\paragraph{Part~\ref{part:path}: exact under analyticity, with the series kept whole.}
Proposition~\ref{prop:closure}; the series in time, Eq.~\eqref{eq:timeseries}; the series forms in Lemma~\ref{lem:drift}(b), Theorem~\ref{thm:time0} and Proposition~\ref{prop:onestep}.

\paragraph{Part~\ref{part:path}: proved from an explicit assumption.}
Propositions~\ref{prop:curv} and~\ref{prop:GN}.

\paragraph{Open.}
(i) The sign of $\kappa_S$, equivalently of the second-order part of a gradient step along the error; the sign of $\ip{q}{E}$; the sign of $\ip{v_k}{\dot Kv_k}$ on slow directions. The last two are one question, since $\ip{q}{E} = -\frac12\ip{v_1}{\dot Kv_1}$; the first is a different contraction of the same second derivative (Section~\ref{sec:loss}).
(ii) Whether $\ell_\ast$, the path length up to which the time-0 bounds are informative, covers a useful stretch of training.
(iii) Whether Assumptions~\ref{ass:curv} and~\ref{ass:GN} hold on the runs of Part~\ref{part:apps}, and with what constants.
(iv) A closure of the growth of the kernel on the slow directions, which is what a prediction of when a rough target is learned would need.
(v) An upper bound on the error at a given roughness of the dataset beyond the short-path band of Theorem~\ref{thm:duhamel} and the assumptions of Section~\ref{sec:closing}; it needs a model of how well a trained network fits a target of a given roughness, and such a model has not been established.

% ===============================================================
\part{Applications}\label{part:apps}
% ===============================================================

% TODO(author): the applications are to be added here, after the theory.
% Planned content, from the earlier drafts:
%  1. Numerical checks of Part I on sine targets: error against the roughness
%     of the target, the slow parts of target and output, the Krylov series
%     against the cluster, high-pass and Lipschitz seminorms, and the replayed
%     Adam steps (main.tex, Section "Numerical checks of the framework").
%  2. Learned time steppers (the design comparison along training is now in
%     Section sec:designcompare; the rest is still to be added): the step map Q + c_p N, the target U = Phi_h - Q,
%     Assumption "Regularity", the local complexity R(Phi_h - Q) (needs a new
%     symbol; \mathcal{H} is now the loss Hessian), Proposition "Absolute
%     roughness is a property of the state", Lemma "Local complexity of the
%     two steps", Theorem "The direct network must also build the identity",
%     and the Lorenz-63 results (main.tex, Sections 7 and 8).
%  3. Gradient descent at the edge of stability (eos_section.tex), which can
%     now cite Propositions prop:hessian and prop:lambda1 of Part II in place
%     of its own Proposition "Sharpness and top eigenvalue" and part (d) of
%     "The kernel after one step", and states the consequence of
%     Assumption ass:sharp (the self-stabilization balance, a heuristic).
%  4. Measurements of the quantities of Section sec:measure.
%  The bibliography entries of eos_section (eos_bibitems.tex) return with 3.

This part applies the framework to design choices.
Section~\ref{sec:designcompare} gives the first application: the closed-form case of the comparison principle of Corollary~\ref{cor:compare}, with learned time steppers as the example.
Still to be added are the numerical checks of Part~\ref{part:onetime} on targets of increasing roughness, the rest of the application to learned time steppers with the comparison on Lorenz-63, the application to gradient descent at the edge of stability, which uses Propositions~\ref{prop:hessian} and~\ref{prop:lambda1} and Assumption~\ref{ass:sharp}, and the measurements listed in Section~\ref{sec:measure}.

% ===============================================================
\section{Comparing two designs at a given training time}\label{sec:designcompare}
% ===============================================================

A design choice changes the problem that the network is given: the labels it must fit, the known part of a step map, the scale at which its output enters, its initialization.
This section proves that one design has the smaller error than another at every time up to a given training time, from the data, the initial parameters, and bounds on the network over a ball about them.
It is Corollary~\ref{cor:compare} with the cheapest bounds that the framework supplies, which need no iterates of Theorem~\ref{thm:predict}: one upper bound, on the error of the first design, and one lower bound, on the error of the second, both valid at every time up to $T$.
A sharper comparison at a given time uses the predicted bands of Theorem~\ref{thm:predict} in the same corollary.
Under gradient flow the upper bound is free: the error never increases.
The lower bound comes from the theory: the error cannot fall faster than the top eigenvalue of the kernel allows, and it cannot fall below what the parameter path can supply (Theorems~\ref{thm:path} and~\ref{thm:time0}).
Nothing is assumed about the architecture.

\paragraph{Setting.}
Two designs, A and B, are trained on the same training inputs.
Each has its own network, parameters, output scale and labels, and its error $E_A$, $E_B$ is measured in the same units, for a learned time stepper in state units.
Both networks are twice continuously differentiable in their parameters, and each is trained by preconditioned gradient flow on its loss $\frac12\|\Net - Y/c_p\|^2$, as in Section~\ref{sec:kernel}, with the flow defined on $[0,T]$; then $\dot E = -KE$ with the kernel of that design, whatever its output scale.
Times $t$ are the gradient-flow times of each training, and the comparison is at equal training time, which is a comparison at equal step size of the underlying optimizer on these losses; the output scale enters only through the kernel of the network each design uses (Proposition~\ref{prop:scale}).

\begin{lemma}[The error falls monotonically, and at most exponentially fast]\label{lem:monoexp}
Under preconditioned gradient flow with a fixed step-size matrix, for every $t\ge0$,
\begin{equation}\label{eq:monoexp}
\|E(0)\|\,\exp\Big(-\int_0^t\lambda_1(s)\,ds\Big)\ \le\ \|E(t)\|\ \le\ \|E(0)\| ,
\end{equation}
where $\lambda_1(s)$ is the top eigenvalue of the kernel at time $s$.
\end{lemma}

\begin{proof}
By Eq.~\eqref{eq:Edotgf}, $\frac{d}{dt}\|E\|^2 = -2\ip{E}{KE}$, and $0\le\ip{E}{KE}\le\lambda_1\|E\|^2$ because $K$ is positive semidefinite.
So $\|E\|^2$ is nonincreasing, and $\frac{d}{dt}\log\|E\|^2\ge-2\lambda_1$ wherever $E\neq0$; integrating gives the lower bound, and if $E(0)\neq0$ the lower bound shows that $E$ never vanishes.
\end{proof}

The lower bound uses only the top eigenvalue along the path.
It can be bounded by the initial kernel and the second derivatives of the network over a ball about the initial parameters: with $\Pi = I$ and $c_p = 1$, the path up to time $t$ stays in the ball of radius $\sqrt{t/2}\,\|E(0)\|$ about $\theta_0$ (Proposition~\ref{prop:ell}), and on that ball $\lambda_1(\theta)^{1/2}\le\lambda_1(0)^{1/2} + \sigma(r)$ by Lemma~\ref{lem:drift}(d), as in the proof of Theorem~\ref{thm:picard}.

\begin{theorem}[Design comparison along training]\label{thm:designcompare}
Let designs A and B be trained by preconditioned gradient flow with fixed step-size matrices, and let $T > 0$.
Let $L_B(T)$ be any of the following lower bounds on $\|E_B(T)\|$:
\begin{enumerate}[label=(\roman*)]
\item $\|E_B(0)\|\exp\big(-\int_0^T\lambda_1^B(s)\,ds\big)$, or $\|E_B(0)\|\,e^{-\bar\lambda_BT}$ for any $\bar\lambda_B$ at least the top eigenvalue of the kernel of B along its path up to $T$ (Lemma~\ref{lem:monoexp});
\item with $\Pi = I$ and $c_p = 1$ for B, $F_0^B(\ell_B) - \tfrac12\bar H_2^B(\ell_B)\,\ell_B^2$, where $F_0^B$ is $F_0$ of Theorem~\ref{thm:time0} for the initial error and kernel of B and $\ell_B\ge\ell_B(T)$, for instance $\ell_B = \sqrt{T/2}\,\|E_B(0)\|$ (Theorem~\ref{thm:time0} and Proposition~\ref{prop:ell});
\item with $\Pi = I$ and $c_p = 1$ for B, and any roughness function $\Rgh$, $\big(\Rgh(Y_B) - \Rgh(\hat Y_{B,0}) - \mathcal{S}_\Rgh(\ell_B)\big)/\CR$ with $\ell_B$ as in (ii) (Corollary~\ref{cor:mintime}(a)).
\end{enumerate}
If $\|E_A(0)\| < L_B(T)$, then $\|E_A(t)\| < \|E_B(t)\|$ for every $t\in[0,T]$.
\end{theorem}

\begin{proof}
By Lemma~\ref{lem:monoexp}, $\|E_A(t)\|\le\|E_A(0)\|$ for every $t$.
Each of (i)--(iii) is a lower bound on $\|E_B(t)\|$ for every $t\le T$: (i) by Lemma~\ref{lem:monoexp}, since $\int_0^t\lambda_1^B\le\int_0^T\lambda_1^B$; (ii) by Theorem~\ref{thm:time0}, since $F_0^B(\ell) - \tfrac12\bar H_2^B(\ell)\ell^2$ is nonincreasing in $\ell$ and $\ell_B(t)\le\ell_B(T)\le\ell_B$; (iii) by Corollary~\ref{cor:mintime}(a), since $\mathcal{S}_\Rgh$ is nondecreasing.
So $\|E_A(t)\|\le\|E_A(0)\| < L_B(T)\le\|E_B(t)\|$.
\end{proof}

Theorem~\ref{thm:designcompare} says that within a training budget $T$, a design that starts closer to its target wins against a design that cannot close its own gap in time.
The bound (i) charges the second design the fastest possible exponential decay; the bound (ii) charges it the part of its initial error that its features at the start cannot remove with a path of the length it can travel, computed by the Krylov series, minus the motion of the kernel; the bound (iii) is the same in terms of the absolute roughness of its labels.
Every quantity is defined for any twice continuously differentiable network and is computed from the data and the initial parameters, except $\lambda_1^B$ along the path, or its bound $\bar\lambda_B$, and $\bar H_2^B$, which are maxima over a ball.
The theorem covers $[0,T]$, and every lower bound becomes vacuous as $T\to\infty$ (for (ii) and (iii) when $\bar H_2^B > 0$), so it says nothing about the limit of long training.
A sharp comparison at a given $T$, beyond the trivial upper bound $\|E_A(0)\|$, needs two-sided enclosures, such as the Picard iterates of Theorem~\ref{thm:picard}.

\paragraph{Known parts.}
For learned time steppers (Section~\ref{sec:setting}), the step map is $x_{j+1} = Q(x_j) + c_p\Net(x_j)$, the labels are $U = \Phi_h - Q$, and the error is $E = c_p\Net - U$ in state units.
If the output of the network is zero at initialization, for instance because its output layer is initialized to zero, then $E(0) = -U$ and $\|E(0)\| = \|\Phi_h - Q\|$, the size in state units of what the known part leaves to the network.

\begin{corollary}[A known part that leaves less wins within a budget]\label{cor:knownpart}
Let A and B be steppers with known parts $Q_A$ and $Q_B$, trained as in Theorem~\ref{thm:designcompare}, where the network of A has output zero at initialization.
If $\|\Phi_h - Q_A\| < L_B(T)$ for one of the lower bounds of Theorem~\ref{thm:designcompare}, then stepper A has the smaller error in state units at every time in $[0,T]$.
If the network of B also has output zero at initialization, then $\|E_B(0)\| = \|\Phi_h - Q_B\|$ in (i).
\end{corollary}

\begin{proof}
Theorem~\ref{thm:designcompare} with $E_A(0) = -(\Phi_h - Q_A)$, and $E_B(0) = -(\Phi_h - Q_B)$ when the output of B is zero at initialization.
\end{proof}

\paragraph{The Euler step and the direct step.}
The Euler step has $Q_{\mathrm E} = \id$ and the direct step $Q_{\mathrm D} = 0$.
Let $\Phi_t$ be the flow of $\dot x = f(x)$, fix $h_0 > 0$, and let
\[
M_1 = \max_i\ \sup_{0\le t\le h_0}\big|f(\Phi_t(x_i))\big| ,
\]
the largest speed along the training trajectories over a time $h_0$; it does not depend on the time step $h\le h_0$, and we assume $M_1 > 0$.
Then, for every $h\le h_0$,
\[
|\Phi_h(x_i) - x_i| = \Big|\int_0^hf(\Phi_t(x_i))\,dt\Big|\le M_1h ,
\]
so $\|\Phi_h - \id\|\le M_1h$ and $\|\Phi_h\|\ge\|\id\| - M_1h$, where $\id$ is the sampled function with values $x_i$.
For the bound (ii), let $F_{0,\id}^{\mathrm D}(\ell) = \max_{\|\varphi\|\le1}\big(\ip{\varphi}{\id} - \ell\,\|\varphi\|_{K_0^{\mathrm D}}\big)$, the part of the identity that the features of the direct network at initialization cannot build with a path of length $\ell$, computed by the Krylov series of Theorem~\ref{thm:stepkrylov}.

\begin{corollary}[The direct step is worse than the Euler step for small time steps]\label{cor:eulerdirect}
Let the Euler and the direct networks be trained as in Theorem~\ref{thm:designcompare}, with outputs zero at initialization and with initial parameters that do not depend on $h$, and let $T > 0$ and $h\le h_0$.
\begin{enumerate}[label=(\alph*)]
\item Let $\bar\lambda_{\mathrm D}$ bound the top eigenvalue of the kernel of the direct network along its path up to $T$, for every time step in $(0, h_0]$; with $\Pi = I$ and $c_p = 1$, one such bound is the largest top eigenvalue over the ball of radius $\sqrt{T/2}\,(\|\id\| + M_1h_0)$ about the initial parameters (Proposition~\ref{prop:ell}).
If
\begin{equation}\label{eq:EDexp}
M_1h\ <\ \big(\|\id\| - M_1h\big)\,e^{-\bar\lambda_{\mathrm D}T},
\end{equation}
then the direct step has the larger error at every time in $[0,T]$.
In particular this holds for every $h < h^\ast(T)$ with
\begin{equation}\label{eq:hstar}
h^\ast(T) = \min\Big(h_0,\ \frac{\|\id\|}{M_1}\,\frac{e^{-\bar\lambda_{\mathrm D}T}}{1 + e^{-\bar\lambda_{\mathrm D}T}}\Big) > 0 .
\end{equation}
\item With $\Pi = I$ and $c_p = 1$ for the direct network, let $\ell_{\mathrm D} = \sqrt{T/2}\,\big(\|\id\| + M_1h\big)$. If
\begin{equation}\label{eq:EDpath}
2M_1h\ <\ F_{0,\id}^{\mathrm D}(\ell_{\mathrm D}) - \tfrac12\bar H_2^{\mathrm D}(\ell_{\mathrm D})\,\ell_{\mathrm D}^2 ,
\end{equation}
then the direct step has the larger error at every time in $[0,T]$.
\end{enumerate}
So for every training budget $T$ there is a time step below which the direct step is worse than the Euler step throughout $[0,T]$.
\end{corollary}

\begin{proof}
For the Euler network, $\|E_{\mathrm E}(t)\|\le\|E_{\mathrm E}(0)\| = \|\Phi_h - \id\|\le M_1h$ by Lemma~\ref{lem:monoexp}.
(a) For the direct network, $\|E_{\mathrm D}(0)\| = \|\Phi_h\|\ge\|\id\| - M_1h$, and Lemma~\ref{lem:monoexp} gives $\|E_{\mathrm D}(t)\|\ge(\|\id\| - M_1h)e^{-\bar\lambda_{\mathrm D}t}\ge(\|\id\| - M_1h)e^{-\bar\lambda_{\mathrm D}T}$.
Under Eq.~\eqref{eq:EDexp} this exceeds $M_1h$.
For $h < h^\ast(T)$, $h\le h_0$, so $M_1$ and $\bar\lambda_{\mathrm D}$ apply, and $M_1h(1 + e^{-\bar\lambda_{\mathrm D}T}) < \|\id\|e^{-\bar\lambda_{\mathrm D}T}$, which is Eq.~\eqref{eq:EDexp}.
The ball in the bound on $\bar\lambda_{\mathrm D}$ contains the direct path up to $T$ for every $h\le h_0$, since $\ell_{\mathrm D}(t)\le\sqrt{t/2}\,\|\Phi_h\|\le\sqrt{T/2}\,(\|\id\| + M_1h_0)$.
(b) The function $E\mapsto\max_{\|\varphi\|\le1}(\ip{\varphi}{E} - \ell\|\varphi\|_{K_0})$ is 1-Lipschitz, being a maximum of 1-Lipschitz functions, and it is unchanged when $E$ is replaced by $-E$ and $\varphi$ by $-\varphi$.
With $E_{\mathrm D}(0) = -\Phi_h$ and $\|\Phi_h - \id\|\le M_1h$, $F_0^{\mathrm D}(\ell)\ge F_{0,\id}^{\mathrm D}(\ell) - M_1h$.
By Proposition~\ref{prop:ell}, $\ell_{\mathrm D}(t)\le\sqrt{t/2}\,\|\Phi_h\|\le\ell_{\mathrm D}$, and, as in the proof of Theorem~\ref{thm:designcompare}(ii), the right side of Theorem~\ref{thm:time0} is nonincreasing in $\ell$, so Theorem~\ref{thm:time0} gives $\|E_{\mathrm D}(t)\|\ge F_{0,\id}^{\mathrm D}(\ell_{\mathrm D}) - \tfrac12\bar H_2^{\mathrm D}(\ell_{\mathrm D})\ell_{\mathrm D}^2 - M_1h$, which under Eq.~\eqref{eq:EDpath} exceeds $M_1h$.
The last statement follows from (a), since $h^\ast(T) > 0$.
\end{proof}

The mechanism is the size of the target in state units, not its roughness.
The Euler network starts within $\|\Phi_h - \id\| = O(h)$ of its target, and gradient flow never moves it further away; the direct network starts at a distance $\|\Phi_h\|$ that does not shrink with $h$, and Lemma~\ref{lem:monoexp} and Theorem~\ref{thm:time0} bound how fast it can close it.
The bound (a) needs only the top eigenvalue of the direct kernel and gives a threshold $h^\ast(T)$ that decreases exponentially in $T$; the bound (b) uses the part of the identity that the direct features at initialization cannot build, is not exponential in $T$, and can be the larger for moderate $T$, as long as the path length stays below the length at which its kernel-motion term, which grows linearly in $T$, catches up with $F_{0,\id}^{\mathrm D}$.
If the network has output biases and the targets are centred, the identity is replaced by the centred states throughout.

What the corollary does not cover is the limit of long training at a fixed time step.
As $T\to\infty$ both thresholds tend to zero, while whether the Euler step remains the more accurate one after long training at small $h$ is to be checked in the Lorenz-63 runs.
That limit is set by how well each network can fit its target at all, and a proof of it needs either the two-sided Picard enclosures of Theorem~\ref{thm:picard} out to the training time or an assumption on how the kernel grows on the slow directions (Section~\ref{sec:closing}).
The corollary is also stated for gradient flow, under which the error is monotone; under Adam the error can rise, and the upper bound on the Euler error does not hold as stated.
It concerns the error of one step, not the error of a rollout.

% ===============================================================
\section*{Related work}\label{sec:related}
\addcontentsline{toc}{section}{Related work}
% ===============================================================

The training kernel of Section~\ref{sec:kernel} is the neural tangent kernel of \cite{jacot2018} at the current parameters, with the preconditioner of the optimizer absorbed into the features.
Nothing here assumes that it stays fixed.
That gradient descent fits the low-frequency components of a target before the high-frequency ones is the spectral bias, or frequency principle, of \cite{rahaman2019,basri2019,xu2020}; \cite{basri2020} extends it to inputs of non-uniform density, where the slow directions are the high eigenfunctions of a Laplacian on the data.
Those results describe the dynamics of a fixed kernel and give rates for smooth targets.
Part~\ref{part:onetime} gives lower bounds only, holds at every time with the kernel of that time, makes no assumption on the update rule, and replaces frequency by an arbitrary seminorm, fixed by the training inputs or built from the error and the kernel.
Part~\ref{part:path} does not hold the kernel fixed either, and does not truncate the tower of derivatives at the linearized or the quadratic model; it expands about a point of the path to any order, with every remainder kept.
Residual and Euler-type steppers and their continuous-time limits are studied in \cite{chen2018}, and the role of the time step and of the known part of a learned stepper in \cite{krishnapriyan2023}; these enter with the applications of Part~\ref{part:apps}.

% ===============================================================
\appendix
\section{Symbols}\label{app:symbols}

Table~\ref{tab:symbols} lists the symbols used in Parts~\ref{part:onetime} and~\ref{part:path}.

\begin{table}[H]
\centering
\caption{Symbols.}\label{tab:symbols}
\footnotesize
\renewcommand{\arraystretch}{0.9}
\begin{tabular}{@{}p{0.25\textwidth}p{0.70\textwidth}@{}}
\toprule
Symbol & Meaning \\
\midrule
$\Aset$ & data set, a compact set containing the training inputs \\
$x_1,\dots,x_n$, $Y$ & training inputs and labels; $(X,Y)$ is the dataset \\
$\ip{\cdot}{\cdot}$, $\|\cdot\|$ & inner product and root-mean-square norm of sampled functions \\
$\supA{u}$, $\Lip(u)$ & $\sup_{x\in\Aset}|u(x)|$ and the Lipschitz constant of $u$ on $\Aset$ \\
$\Rgh$ & roughness function, a seminorm on sampled functions; a local complexity function when fixed by the inputs \\
$\CR$ & norm constant of $\Rgh$ \\
$\Rgh_{\mathrm{cl}}$, $\Rgh_m$, $\Lipn$ & cluster seminorm, manifold high-pass above the $m$th graph-Laplacian eigenvector, Lipschitz seminorm on the training inputs \\
$\Net$, $\theta$, $c_p$ & network, its parameters, and the scale at which its output enters the prediction \\
$\hat Y$, $E$ & prediction $c_p\Net$ and error $\hat Y - Y$ at the training inputs, in the units of the labels \\
$\Rgh(Y)$, $\Rgh(Y/\|Y\|)$ & absolute and relative roughness of the dataset under $\Rgh$ \\
$\Psi$, $\Pi$, $K$ & gradient features, step-size matrix, training kernel \\
$\lambda_i$, $v_i$, $\hat e_i$ & eigenvalues and eigenvectors of $K$, components $\ip{v_i}{E}$ of the error \\
$\gamma$, $E_\gamma$, $\Upsilon(\gamma)$ & rate, slow part $\gamma(K+\gamma I)^{-1}E$ of the error at rate $\gamma$, squared length of the shortest parameter step that produces the removed part $E - E_\gamma$ \\
$\sR$ & peak feature roughness under $\Rgh$ \\
$\bar\sigma_\Rgh(s)$ & largest $\sR$ at any parameters that a step of rescaled length at most $s$ can reach \\
$\hat g$ & rescaled parameter velocity \\
$s$, $s_\gamma$ & rescaled length of a parameter step; the length $\Upsilon(\gamma)^{1/2}/|c_p|$ whose best step leaves $E_\gamma$ \\
$g_\varphi$, $\|\varphi\|_K$ & gradient $\frac1n\Psi^\top\varphi$ and kernel size $|g_\varphi|$ of a sampled function $\varphi$ \\
$b_\Rgh(s)$ & step bound of $\Rgh$, $(\Rgh(E) - |c_p|\,s\,\sR)/\CR$ \\
$\mathcal{D}_s$, $F(s)$, $F_M(s)$ & first-order output changes of steps of length at most $s$; smallest error among them; its Krylov series \\
$V_M$, $q_E$ & Krylov subspace spanned by $E, KE, \dots, K^{M-1}E$; number of distinct eigenvalues of $K$ on which $E$ has a component \\
$\varphi^\star$ & best roughness functional, the unit vector along the smallest reachable error $E^\star$ \\
\midrule
$L$, $g_E$ & loss $\frac12\|E\|^2$ and its gradient $\frac1n\Psi^\top E$ (Part~\ref{part:path}: $\Pi = I$, $c_p = 1$) \\
$\Hess$, $G$, $S(\varphi)$ & loss Hessian, its kernel term $\frac1n\Psi^\top\Psi$, the Hessian of $\ip{\varphi}{\Net}$ in $\theta$ \\
$T_m(\varphi)$, $D_m(\varphi)$, $H_m$ & $m$-th derivative of $\ip{\varphi}{\Net}$ in $\theta$, its operator norm, the root-mean-square size of $\nabla^m_\theta\Net$ \\
$\bar D_m(\varphi;r)$, $\bar H_m(r)$ & largest $D_m(\varphi;\theta)$, $H_m(\theta)$ over the ball of radius $r$ about $\theta_0$ \\
$S_k(E)$, $G_k$ & error part and error-free part of $\nabla^k_\theta L$ \\
$q$, $a_1$ & output change $\Psi a_1$ and parameter direction $a_1 = S(v_1)g_{v_1}$ of the rate of the top eigenvalue \\
$\Delta$ & displacement $\theta(t) - \theta_0$ \\
$\kappa$, $\kappa_G$, $\kappa_S$ & curvature of the loss along the gradient, its kernel term and its remaining part \\
$\rho$, $w_k$ & instantaneous rate $|g_E|^2/L$ of the loss, distribution $\hat e_k^2/\|E\|^2$ of the error over the eigenvectors \\
$m_k$ & $\nabla^k_\theta L$ contracted with $k$ copies of $g_E$ \\
$D$ & derivative along the vector field $V = -\nabla_\theta L$ of gradient flow \\
$\ell$, $\cost_t(\varphi)$ & length of the parameter path; cost of the direction $\varphi$ along it \\
$F_t$, $F_0(\ell)$ & floor along the path; $F(\ell)$ with the features at $\theta_0$ \\
$\sigma(r)$, $\sigma_k(r)$ & drift of all gradients and of the gradient of $v_k$ over a path of length $r$ \\
$\delta K$, $\beta$, $\mathcal{I}_n$ & kernel motion $K(s) - K_0$, its integrated size, the terms of the Dyson series \\
$N$, $\Lambda_N(r)$ & part of the gradient beyond its first-order expansion about $\theta_0$, its Lipschitz constant on the ball of radius $r$ \\
$z = (\theta,\varsigma)$, $F$, $M_j$ & state of training (parameters and optimizer state), the flow or the maps of the update rule \\
$A(t)$, $B_j$, $U$, $\Gamma$, $\xi$, $\zeta$ & reference linear part, its propagator, and the constants of the bound on the propagator \\
$\Lambda_F(r)$, $\varepsilon_k$, $P_r$ & Lipschitz constant of the remainder on the ball, certified error of the $k$-th iterate, projection onto the ball \\
$\Delta_{[k]}$, $E_{[k]}$, $\mu$ & Picard iterates and the contraction factor \\
\bottomrule
\end{tabular}
\end{table}

% ===============================================================
\section{Proofs for Part~\ref{part:onetime}}\label{app:proofs}
% ===============================================================

\subsection{Proofs of the lower bounds}\label{app:floor}

\paragraph{Lemma~\ref{lem:mech}.}
By Eq.~\eqref{eq:kernel}, $(Kw)_j = \tfrac1n\sum_l\Psi_j\Psi_l^\top w_l = \Psi_j g_w$, so $Kw = \Psi g_w$.
By Eq.~\eqref{eq:sigma}, $\Rgh(\Psi g_w) \le \sR|g_w|$.
Finally $|g_w|^2 = \tfrac1n\sum_i w_i\cdot(\Psi_i g_w) = \ip{w}{Kw}$. \qed

\paragraph{Corollary~\ref{cor:budget}.}
$\Rgh(Kv_i) = \lambda_i\Rgh(v_i)$ by homogeneity, and Lemma~\ref{lem:mech} with $w = v_i$ gives $\lambda_i\Rgh(v_i) \le \sR\ip{v_i}{Kv_i}^{1/2} = \sR\lambda_i^{1/2}$; squaring gives $\lambda_i^2\Rgh(v_i)^2\le\sR^2\lambda_i$, and Eq.~\eqref{eq:budget} follows by dividing by $\lambda_i$ when $\lambda_i > 0$, and trivially when $\lambda_i = 0$. \qed

\paragraph{Theorem~\ref{thm:fast}.}
(a) Let $w = (K+\gamma I)^{-1}u = \sum_i \hat u_i v_i/(\lambda_i+\gamma)$.
Then $Kw = u - u_\gamma$ and
\begin{equation*}
\ip{w}{Kw} = \sum_i\frac{\lambda_i\,\hat u_i^2}{(\lambda_i+\gamma)^2} = \Upsilon_u(\gamma).
\end{equation*}
Lemma~\ref{lem:mech} with this $w$ gives the first inequality of Eq.~\eqref{eq:fastsmooth}.
Since $\lambda/(\lambda+\gamma)^2 \le 1/(4\gamma)$ for $\lambda \ge 0$, $\Upsilon_u(\gamma) \le \|u\|^2/(4\gamma)$, which gives the second.
(b) By the triangle inequality, Eq.~\eqref{eq:CR}, and part (a),
\begin{equation*}
\Rgh(u) \le \Rgh(u_\gamma) + \Rgh(u - u_\gamma) \le \CR\|u_\gamma\| + \sR\,\Upsilon_u(\gamma)^{1/2}.
\end{equation*}
Rearranging gives Eq.~\eqref{eq:slowlower}. \qed

\paragraph{Theorem~\ref{thm:floor}.}
Since $E = \hat Y - Y$, the reverse triangle inequality gives $\Rgh(E) \ge \Rgh(Y) - \Rgh(\hat Y)$.
Insert this into Theorem~\ref{thm:fast}(b) with $u = E$. \qed

\paragraph{Theorem~\ref{thm:dataset}.}
Eq.~\eqref{eq:dataset} is the second inequality of Proposition~\ref{prop:stepexact} in Appendix~\ref{app:aux}, with $E = \hat Y(t) - Y$ and $\Rgh(\hat Y(t) - Y) \ge \Rgh(Y) - \Rgh(\hat Y(t))$. \qed

\paragraph{Theorem~\ref{thm:slow}.}
(a) The first inequality holds because the weights $\gamma/(\lambda_i+\gamma)$ lie in $(0,1]$.
The second is Theorem~\ref{thm:fast}(b) with $u = E(t)$, and the third is $\Rgh(E(t)) \ge \Rgh(Y) - \Rgh(\hat Y(t))$, the reverse triangle inequality.
(b) The first inequality is Proposition~\ref{prop:stepexact}, and the second is the reverse triangle inequality.
(c) With $s_\gamma = \Upsilon(\gamma)^{1/2}/|c_p|$, the middle term of Eq.~\eqref{eq:main} is the step bound $b_\Rgh(s_\gamma)$ of Eq.~\eqref{eq:bR}, and $F(s_\gamma) = \|E_\gamma(t)\| > 0$ by Proposition~\ref{prop:floorstep}.
By Theorem~\ref{thm:stepsup}(b), $\sup_\Rgh b_\Rgh(s_\gamma) = F(s_\gamma)$, attained by the rank-one seminorm along the minimizer $E^\star$, which is $E_\gamma(t)$ by Theorem~\ref{thm:stepsup}(c) and the proof of Proposition~\ref{prop:floorstep}.
With $\Psi = \Psi(t)$, Theorem~\ref{thm:stepsup}(b) gives the supremum $F(s)$ of $b_\Rgh(s)$ and the functional $\varphi^\star$, Theorem~\ref{thm:stepkrylov}(b) computes it, and Proposition~\ref{prop:taylorstep} gives the two bounds on the error that steps actually reach.
For Eq.~\eqref{eq:mainexact}, the proof of Theorem~\ref{thm:stepsup}(b) gives $F(s) = g_s(\varphi^\star) = \ip{\varphi^\star}{E(t)} - |c_p|s\|\varphi^\star\|_K$, and the first identity is this with $s = s_\gamma$ and $\varphi^\star = \varphi_\gamma$; in both, $\ip{\varphi}{E(t)} = \ip{\varphi}{\hat Y(t)} - \ip{\varphi}{Y}$. \qed

\subsection{Auxiliary one-step results}\label{app:aux}

The four propositions below are used in Sections~\ref{sec:step} and~\ref{sec:lower}; their proofs are in Appendix~\ref{app:step}.
The first bounds every step of length at most $s$ with no assumption on how the gradient features change along it.

\begin{proposition}[Every step of length at most $s$]\label{prop:stepexact}
Let $B_s = \{\theta + \Pi^{1/2}g' : |g'| \le s\}$.
For a sampled function $\varphi$ let $\|\varphi\|_{K,s} = \max_{\theta'\in B_s}\|\varphi\|_{K(\theta')}$, and for a roughness function $\Rgh$ let $\bar\sigma_\Rgh(s) = \max_{\theta'\in B_s}\sR(\theta')$, where $K(\theta')$ and $\sR(\theta')$ are computed from the gradient features at $\theta'$.
Every step of rescaled length at most $s$ leaves an error $E_{\mathrm{new}}$ with
\begin{equation}\label{eq:stepexact}
\|E_{\mathrm{new}}\| \ \ge\ \sup_{\|\varphi\|\le1}\Big(\ip{\varphi}{E} - |c_p|\,s\,\|\varphi\|_{K,s}\Big) \ \ge\ \frac{1}{\CR}\Big(\Rgh(E) - |c_p|\,s\,\bar\sigma_\Rgh(s)\Big)
\end{equation}
for every roughness function $\Rgh$.
\end{proposition}

The second identifies the floor of Theorem~\ref{thm:floor} with the step bound of one roughness function.

\begin{proposition}[The floor is a one-step bound]\label{prop:floorstep}
Let $\Psi = \Psi(t)$.
For $\gamma > 0$, let $s_\gamma = \Upsilon(\gamma)^{1/2}/|c_p|$. Then $F(s_\gamma) = \|E_\gamma\|$, and for every roughness function $\Rgh$,
\begin{equation}\label{eq:floorstep}
\frac{1}{\CR}\Big(\Rgh(Y) - \Rgh(\hat Y) - \sR\,\Upsilon(\gamma)^{1/2}\Big) \ \le\ b_\Rgh(s_\gamma) \ \le\ F(s_\gamma) = \|E_\gamma\|.
\end{equation}
\end{proposition}

The third states that no bound depends on the scale at which the network's output enters the prediction.

\begin{proposition}[Output-scale invariance]\label{prop:scale}
Let $c \neq 0$, and replace $c_p$ by $c\,c_p$ and $\Net$ by $\Net/c$ at the same parameters.
Then the prediction $\hat Y$, the error $E$, and $\Rgh(\hat Y)$ are unchanged; the gradient features, the kernel and the peak feature roughness become $\Psi/c$, $K/c^2$ and $\sR/|c|$.
With the rate $\gamma$ replaced by $\gamma/c^2$, the slow part $E_\gamma$, the product $\sR\,\Upsilon(\gamma)^{1/2}$, and the floor of Theorem~\ref{thm:floor} are unchanged.
For every $s$, the set $\mathcal{D}_s$, the step bounds $b_\Rgh(s)$, the smallest reachable error $F(s)$, the Krylov subspaces $V_M$, and the Krylov values $F_M(s)$ are unchanged, and so are the bounds of Proposition~\ref{prop:rates} for a given $\hat g$.
The quantities $\bar\sigma_\Rgh(s)$, $\|\varphi\|_{K,s}$ and $\bar H_2(s)$ become $\bar\sigma_\Rgh(s)/|c|$, $\|\varphi\|_{K,s}/|c|$ and $\bar H_2(s)/|c|$, so the bounds of Theorem~\ref{thm:dataset} and of Propositions~\ref{prop:stepexact} and~\ref{prop:taylorstep} are unchanged.
\end{proposition}

The fourth gives the rate at which the Krylov series of Theorem~\ref{thm:stepkrylov} converges.

\begin{proposition}[Convergence of the Krylov series]\label{prop:krylovrate}
Let $F(s) > 0$, let $\varphi^\star$ be the functional of Theorem~\ref{thm:stepsup}(b), and let $d_M = \min_{\varphi\in V_M}\|\varphi - \varphi^\star\|$.
\begin{enumerate}[label=(\alph*)]
\item $\tfrac12 F(s)\,d_M^2 \le F(s) - F_M(s) \le \big(\|E\| + |c_p|\,s\,\lambda_1^{1/2}\big)\,d_M$. Hence $F_M(s) \to F(s)$ if and only if $d_M \to 0$, and $F_M(s) = F(s)$ if and only if $\varphi^\star \in V_M$.
\item If $0 < (c_ps)^2 < \Upsilon(0^+)$, let $\chi = (\lambda_1+\gamma)/\gamma$ with $\gamma$ from Theorem~\ref{thm:stepsup}(c). Then
\begin{equation}\label{eq:krylovratestep}
d_M \le \frac{\chi-1}{2}\Big(\frac{\sqrt\chi-1}{\sqrt\chi+1}\Big)^{M-1}.
\end{equation}
\end{enumerate}
\end{proposition}

Part (a) holds for every nested family of subspaces, including subspaces fixed by the training inputs; such a family converges as fast as it captures $\varphi^\star$.
The rate in part (b) is the Chebyshev bound of the conjugate gradient method \cite{saad2003}, with $\chi$ at least the condition number of $K + \gamma I$, and equal to it when $K$ is singular.
A long step gives a small $\gamma$ and a large $\chi$, so the series converges slowly for steps that nearly fit the error.

\subsection{Proofs of the one-step results}\label{app:step}

Write $g_s(\varphi) = \ip{\varphi}{E} - |c_p|\,s\,\|\varphi\|_K$.

\paragraph{Proposition~\ref{prop:rates}.}
(a) By Eq.~\eqref{eq:Edot}, $\tfrac{d}{dt}\ip{\varphi}{E} = c_p\,\hat g\cdot g_\varphi$, and the Cauchy--Schwarz inequality gives the bound.
(b) $E(t+\tau) - E(t) = \tau\dot E + o(\tau)$.
The reverse triangle inequality and Eq.~\eqref{eq:CR} give $|\Rgh(E(t+\tau)) - \Rgh(E(t))| \le |\tau|\,\Rgh(\dot E) + \CR\,o(\tau)$, and $\Rgh(\dot E) = |c_p|\,\Rgh(\Psi\hat g) \le |c_p|\,\sR|\hat g|$ by Eq.~\eqref{eq:sigma}.
(c) $\tfrac{d}{dt}\|E\| = \ip{\hat E}{\dot E}$, and (a) with $\varphi = \hat E$ gives the bound.
Under gradient flow, $\ip{\hat E}{\dot E} = -\ip{\hat E}{KE} = -\|E\|_K^2/\|E\|$, and $|c_p||\hat g| = |g_E| = \|E\|_K = \|E\|\,\|\hat E\|_K$. \qed

\paragraph{Proposition~\ref{prop:stepexact}.}
The step is $\Delta\theta = \Pi^{1/2}g$ with $|g| \le s$, so every point $\theta + \tau\Delta\theta$ with $\tau\in[0,1]$ lies in $B_s$.
By Eq.~\eqref{eq:avgstep}, $E_{\mathrm{new}} = E + c_p\bar\Psi g$, and for $\|\varphi\|\le1$,
$|\ip{\varphi}{\bar\Psi g}| = \big|\int_0^1 g\cdot g_\varphi(\theta+\tau\Delta\theta)\,d\tau\big| \le |g|\,\|\varphi\|_{K,s}$, so $\ip{\varphi}{c_p\bar\Psi g}\ge-|c_p|\,s\,\|\varphi\|_{K,s}$, where $g_\varphi(\theta') = \tfrac1n\Psi(\theta')^\top\varphi$ has length $\|\varphi\|_{K(\theta')}$.
Hence $\|E_{\mathrm{new}}\| \ge \ip{\varphi}{E_{\mathrm{new}}} \ge \ip{\varphi}{E} - |c_p|s\|\varphi\|_{K,s}$.
For the second inequality, let $\varphi_0\in D_\Rgh$ attain $\Rgh(E)$ as in the proof of Theorem~\ref{thm:stepsup}(a); then $\|\varphi_0\| \le \CR$ and $\|\varphi_0\|_{K(\theta')} \le \sR(\theta')$ for every $\theta'$, so $\varphi = \varphi_0/\CR$ gives the bound.
The maxima exist because the network is continuously differentiable and $B_s$ is compact. \qed

\paragraph{Theorem~\ref{thm:stepsup}.}
Weak duality: for $|g| \le s$ and $\|\varphi\| \le 1$, $\ip{\varphi}{c_p\Psi g} = c_p\,g\cdot g_\varphi \ge -|c_p|\,s\,\|\varphi\|_K$, so $g_s(\varphi) \le \ip{\varphi}{E + c_p\Psi g} \le \|E + c_p\Psi g\|$.
Hence $g_s(\varphi) \le F(s)$.

(a) Let $D_\Rgh = \{\varphi : \ip{\varphi}{u} \le \Rgh(u) \text{ for all } u\}$, the dual ball of $\Rgh$.
By the separation theorem, $\Rgh(u) = \max_{\varphi\in D_\Rgh}\ip{\varphi}{u}$.
Every $\varphi\in D_\Rgh$ has $\|\varphi\| \le \CR$ and $\|\varphi\|_K = \max_{|g|=1}\ip{\varphi}{\Psi g} \le \sR$.
Let $\varphi_0\in D_\Rgh$ attain $\Rgh(E)$.
Then $b_\Rgh(s) \le \big(\ip{\varphi_0}{E} - |c_p|s\|\varphi_0\|_K\big)/\CR = g_s(\varphi_0/\CR) \le F(s)$.

(b) The set $\mathcal{D}_s$ is compact, convex, and symmetric, and its support function is $\max_{\Delta\in\mathcal{D}_s}\ip{\varphi}{\Delta} = |c_p|\,s\,\|\varphi\|_K$.
Let $-\Delta^\star$ be the point of $\mathcal{D}_s$ closest to $E$; it exists and is unique, and $E^\star = E + \Delta^\star$ is the unique minimizer in Eq.~\eqref{eq:Fs}.
Since $F(s) > 0$, $E^\star \neq 0$.
The projection inequality gives $\ip{E^\star}{\Delta'} \le \ip{E^\star}{-\Delta^\star}$ for all $\Delta'\in\mathcal{D}_s$, so $|c_p|s\|\varphi^\star\|_K = \ip{\varphi^\star}{-\Delta^\star}$ and $g_s(\varphi^\star) = \ip{\varphi^\star}{E^\star} = F(s)$.
For $\Rgh^\star$, $C_{\Rgh^\star} = 1$ and $\sigma_{\Rgh^\star} = \max_{|g|=1}|\ip{\varphi^\star}{\Psi g}| = \|\varphi^\star\|_K$.
Since $\ip{\varphi^\star}{E} = F(s) + |c_p|s\|\varphi^\star\|_K > 0$, $b_{\Rgh^\star}(s) = g_s(\varphi^\star) = F(s)$.
With (a), this gives the supremum.
If some $\varphi$ with $\|\varphi\|\le1$ has $g_s(\varphi) = F(s)$, then weak duality with $g = g^\star$ gives $F(s) \le \ip{\varphi}{E^\star} \le \|E^\star\| = F(s)$, and equality in the Cauchy--Schwarz inequality gives $\varphi = \varphi^\star$.

(c) For $\Delta$ in the range of $K$, the smallest $|g|$ with $\Psi g = \Delta$ satisfies $|g|^2 = \ip{\Delta}{K^{+}\Delta}$, so $\mathcal{D}_s = \{\Delta \in \operatorname{ran}K : \ip{\Delta}{K^{+}\Delta} \le (c_ps)^2\}$.
The component of $E$ outside the range of $K$ is not changed by any step.
Without the constraint, the component of $E$ in the range can be removed, which requires $\ip{\Delta}{K^{+}\Delta} = \Upsilon(0^+) = \sum_{\lambda_i>0}\hat e_i^2/\lambda_i$.
If $(c_ps)^2 \ge \Upsilon(0^+)$ this is feasible, and only the component outside the range remains.
Otherwise the constraint is active; the problem is convex, so the Lagrange condition $E - \Delta = \gamma K^{+}\Delta$ on the range is sufficient.
It gives $\Delta = K(K+\gamma I)^{-1}E$, leaving the slow part $\gamma(K+\gamma I)^{-1}E$ with components $\gamma\hat e_i/(\lambda_i+\gamma)$, and the constraint $\Upsilon(\gamma) = (c_ps)^2$, whose solution is unique because $\Upsilon$ decreases strictly.

(d) The bound is $g_s(\hat E) \le F(s)$.
If $\|\hat E\|_K = 0$, then $E$ is orthogonal to the range of $K$ and $F(s) = \|E\|$.
Otherwise take $g = -\operatorname{sign}(c_p)\,s\,g_{\hat E}/\|\hat E\|_K$, so that $c_p\Psi g = -|c_p|s\,K\hat E/\|\hat E\|_K$.
Since $\ip{E}{K\hat E} = \|E\|\,\|\hat E\|_K^2$, $\|E + c_p\Psi g\|^2 = \|E\|^2 - 2|c_p|s\|E\|\|\hat E\|_K + O(s^2)$, whose square root is $\|E\| - |c_p|s\|\hat E\|_K + O(s^2)$. \qed

\paragraph{Proposition~\ref{prop:floorstep}.}
By definition, $|c_p|\,s_\gamma = \Upsilon(\gamma)^{1/2}$, so $b_\Rgh(s_\gamma) = \big(\Rgh(E) - \sR\,\Upsilon(\gamma)^{1/2}\big)/\CR$.
The reverse triangle inequality gives $\Rgh(E) \ge \Rgh(Y) - \Rgh(\hat Y)$, which is the first inequality of Eq.~\eqref{eq:floorstep}.
The second is Theorem~\ref{thm:stepsup}(a).
For the equality, $\Upsilon$ is strictly decreasing on $(0,\infty)$ if $E$ has a component in the range of $K$, so $(c_ps_\gamma)^2 = \Upsilon(\gamma) < \Upsilon(0^+)$ and $\gamma$ is the unique solution of Theorem~\ref{thm:stepsup}(c), which gives $E^\star = E_\gamma$ and $F(s_\gamma) = \|\gamma(K+\gamma I)^{-1}E\| = \|E_\gamma\|$.
Otherwise $\Upsilon \equiv 0$, $s_\gamma = 0$, and $F(0) = \|E\| = \|E_\gamma\|$ because every component of $E$ lies on eigenvalue zero. \qed

\paragraph{Proposition~\ref{prop:scale}.}
The prediction is $(c\,c_p)(\Net/c) = c_p\Net$, so it, the error $E$ and $\Rgh(\hat Y)$ are unchanged.
The gradient features of $\Net/c$ are $\Psi/c$, so the kernel is $K/c^2$, with eigenvalues $\lambda_i/c^2$ and the same eigenvectors, and Eq.~\eqref{eq:sigma} gives the peak feature roughness $\sR/|c|$.
With $\gamma/c^2$ in place of $\gamma$, the weights $(\gamma/c^2)/(\lambda_i/c^2 + \gamma/c^2) = \gamma/(\lambda_i + \gamma)$ are unchanged, so $E_\gamma$ is unchanged, and $\Upsilon$ becomes $\sum_i (\lambda_i/c^2)\hat e_i^2/((\lambda_i + \gamma)/c^2)^2 = c^2\,\Upsilon(\gamma)$, so $\sR\,\Upsilon(\gamma)^{1/2}$ and the floor are unchanged.
The set $\mathcal{D}_s = \{(c\,c_p)(\Psi/c)\,g : |g| \le s\}$ is unchanged, hence so is $F(s)$, and $|c\,c_p|\,s\,\sR/|c| = |c_p|\,s\,\sR$ leaves $b_\Rgh(s)$ unchanged.
The Krylov subspaces of $K/c^2$ are those of $K$, and $|c\,c_p|\,s\,\|\varphi\|_{K/c^2} = |c_p|\,s\,\|\varphi\|_K$, so $F_M(s)$ is unchanged.
In Proposition~\ref{prop:rates}, $|c\,c_p|\,\|\varphi\|_{K/c^2} = |c_p|\,\|\varphi\|_K$ and $|c\,c_p|\,\sR/|c| = |c_p|\,\sR$. \qed

\paragraph{Theorem~\ref{thm:stepkrylov}.}
(a) For $\varphi\in V_M$ with $\|\varphi\| = 1$, $b_\Rgh(s)$ for $\Rgh(u) = |\ip{\varphi}{u}|$ is $|\ip{\varphi}{E}| - |c_p|s\|\varphi\|_K$.
Since $V_M$ is a subspace, the sign of $\ip{\varphi}{E}$ can be chosen, and since $g_s$ is positively homogeneous, the maximum over the unit ball of $V_M$ is the larger of $0$ and the maximum over its unit sphere.
The subspaces are nested, and every $F_M(s) \le F(s)$ by weak duality.
For $V_1 = \operatorname{span}\{E\}$, $g_s(\alpha\hat E) = \alpha\|E\| - |\alpha||c_p|s\|\hat E\|_K$ for $|\alpha|\le1$.
(b) Let $\mu_1,\dots,\mu_{q_E}$ be the distinct eigenvalues of $K$ on which $E$ has a component.
If $s = 0$, then $E^\star = E$ and $\varphi^\star = E/\|E\|\in V_1$.
If $0 < (c_ps)^2 < \Upsilon(0^+)$, let $\varpi$ be the polynomial of degree at most $q_E-1$ with $\varpi(\mu_j) = 1/(\mu_j+\gamma)$; then $\varpi(K)E = (K+\gamma I)^{-1}E \in V_{q_E}$, and $\varphi^\star$ is a multiple of it.
Otherwise $E^\star$ is the component of $E$ in the null space of $K$, which is nonzero because $F(s) > 0$, and it equals $\varpi(K)E$ for the polynomial with $\varpi(0) = 1$ and $\varpi(\mu_j) = 0$ for $\mu_j > 0$. \qed

\paragraph{Proposition~\ref{prop:krylovrate}.}
(a) For the lower half, every $\varphi$ with $\|\varphi\|\le1$ satisfies $g_s(\varphi) \le \ip{\varphi}{E^\star} = F(s)\ip{\varphi}{\varphi^\star}$ by weak duality with $g = g^\star$.
Since $\|\varphi - \varphi^\star\|^2 \le 2 - 2\ip{\varphi}{\varphi^\star}$, $g_s(\varphi) \le F(s)\big(1 - \tfrac12\|\varphi - \varphi^\star\|^2\big)$; apply this to the maximizer in $V_M$.
For the upper half, let $\psi$ be the orthogonal projection of $\varphi^\star$ onto $V_M$, so $\|\psi\|\le1$ and $\|\psi - \varphi^\star\| = d_M$.
Then $|\ip{\psi - \varphi^\star}{E}| \le \|E\|d_M$ and $\big|\|\psi\|_K - \|\varphi^\star\|_K\big| \le \|\psi - \varphi^\star\|_K \le \lambda_1^{1/2}d_M$, so $F_M(s) \ge g_s(\psi) \ge F(s) - (\|E\| + |c_p|s\lambda_1^{1/2})d_M$.
The last statement follows from the uniqueness in Theorem~\ref{thm:stepsup}(b).
(b) Let $w = (K+\gamma I)^{-1}E$, so $\varphi^\star = w/\|w\|$.
For every polynomial $\varpi$ of degree at most $M-1$, $\varpi(K)E/\|w\|\in V_M$, so $d_M \le \|w - \varpi(K)E\|/\|w\|$.
Expanding in the eigenvectors of $K$, $\|w - \varpi(K)E\| \le \max_{\lambda\in[0,\lambda_1]}|(\lambda+\gamma)^{-1} - \varpi(\lambda)|\,\|E\|$, and $\|w\| \ge \|E\|/(\lambda_1+\gamma)$.
The substitution $\lambda = \tfrac{\lambda_1}{2}(\tau+1)$ maps $(\lambda+\gamma)^{-1}$ on $[0,\lambda_1]$ to $\tfrac{2}{\lambda_1}(\tau - a)^{-1}$ on $[-1,1]$ with $a = -(1 + 2\gamma/\lambda_1)$.
The best uniform approximation of $(\tau - a)^{-1}$ on $[-1,1]$ by polynomials of degree $k$ has error $\big(|a| - \sqrt{a^2-1}\big)^k/(a^2-1)$ \cite{achieser1956}.
With $|a| - \sqrt{a^2-1} = (\sqrt\chi-1)/(\sqrt\chi+1)$ and $a^2-1 = 4\gamma(\lambda_1+\gamma)/\lambda_1^2$, the best error for degree $M-1$ is $\tfrac{\lambda_1}{2\gamma(\lambda_1+\gamma)}\big(\tfrac{\sqrt\chi-1}{\sqrt\chi+1}\big)^{M-1}$.
Multiplying by $\lambda_1+\gamma$ gives Eq.~\eqref{eq:krylovratestep}, since $\lambda_1/\gamma = \chi-1$. \qed

% ===============================================================
\section{Other update rules}\label{app:momentum}
% ===============================================================

\paragraph{Minibatches.}
A minibatch $B$ of $b$ training inputs uses the loss on $B$ alone.
Let $W$ be the random sampling matrix that multiplies the point of a sampled function at $x_i$ by $d_i = n/b$ if $i \in B$ and by $d_i = 0$ otherwise.
The minibatch loss is $\tfrac1{2n}\sum_i d_i\,|\Net(x_i) - Y_i/c_p|^2$, and for a fixed batch $B$, preconditioned gradient flow on its loss gives $\hat g = -g_{WE}/c_p$ and
\begin{equation}\label{eq:Edotbatch}
\dot E = -K W E .
\end{equation}
The error $E$ and the kernel $K$ are those of the whole training set, so a minibatch step acts in the same space of sampled functions on all $n$ training inputs, the stacked residuals, as a full-batch step.
The error changes at every training input, not only at those in the batch.
Only the sampling matrix is random.
If $B$ is drawn uniformly, $\mathbb{E}\,W = I$, so at fixed parameters the rate of Eq.~\eqref{eq:Edotbatch}, averaged over $B$, is that of Eq.~\eqref{eq:Edotgf}.
The statements about $K$ itself, its eigenvalues and eigenvectors, do not involve the batch.
Every statement about $Kw$ for a general sampled function $w$ holds for each minibatch step with $w = WE$.
In particular, by Lemma~\ref{lem:mech}, $\Rgh(\dot E) \le \sR\,\|WE\|_K$, and the rescaled parameter velocity has length $|\hat g| = \|WE\|_K/|c_p|$, which enters Proposition~\ref{prop:rates}.
The identities of Part~\ref{part:path} are stated for full-batch gradient flow; Theorem~\ref{thm:path}, which uses only the parameter velocity, holds for minibatch training as well.

\paragraph{Adam with momentum.}
For a continuous-time model of Adam with momentum, the update is $\dot\theta = -\Pi(t)\,q_{\mathrm m}$, where $\Pi(t)$ is the current preconditioner and $q_{\mathrm m}$ is a moving average of past minibatch gradients; the features and kernel are those of $\Pi(t)$ at each time.
The parameter velocity is then not a function of the current error, and Eq.~\eqref{eq:Edotgf} no longer gives the rate of change of the error.
The results of Section~\ref{sec:lower} use the parameter velocity only through the rescaled parameter velocity $\hat g = \Pi^{-1/2}\dot\theta$ of Eq.~\eqref{eq:Edot}, whatever update rule produced it.
Momentum therefore changes $\hat g$ and leaves the kernel, the peak feature roughness, and every bound of Section~\ref{sec:lower} unchanged.
An error with a small removal rate $\ip{E}{KE}/\|E\|^2$ is removed faster under momentum only through a larger $|\hat g|$, and the output gains roughness no faster than $|c_p|\,\sR\,|\hat g|$ (Proposition~\ref{prop:rates}).
The same computation on $n'$ validation inputs, with their own gradient features $\Psi'$, kernel $K'$ and norm $\|\cdot\|'$, bounds the rate of decrease of the validation error $E'$ by $|c_p|\,\|E'/\|E'\|'\|_{K'}\,|\hat g|$.

% ===============================================================
\begin{thebibliography}{99}
\bibitem{achieser1956}
N.~I.~Achieser.
\emph{Theory of Approximation}.
Frederick Ungar, New York, 1956.
% VERIFY: section and page for the best approximation of 1/(x-a) on [-1,1].

\bibitem{hiriart2001}
J.-B.~Hiriart-Urruty and C.~Lemar\'echal.
\emph{Fundamentals of Convex Analysis}.
Springer, Berlin, 2001.
% VERIFY: support functions and the distance to a convex set (Chapter C).

\bibitem{saad2003}
Y.~Saad.
\emph{Iterative Methods for Sparse Linear Systems}.
2nd edition, SIAM, Philadelphia, 2003.
% VERIFY: Chebyshev bound for the conjugate gradient method (Section 6.11).

\bibitem{jacot2018}
A.~Jacot, F.~Gabriel, and C.~Hongler.
Neural tangent kernel: convergence and generalization in neural networks.
In \emph{Advances in Neural Information Processing Systems} 31, 2018.

\bibitem{rahaman2019}
N.~Rahaman, A.~Baratin, D.~Arpit, F.~Draxler, M.~Lin, F.~Hamprecht, Y.~Bengio, and A.~Courville.
On the spectral bias of neural networks.
In \emph{Proceedings of the 36th International Conference on Machine Learning}, 2019.

\bibitem{basri2019}
R.~Basri, D.~Jacobs, Y.~Kasten, and S.~Kritchman.
The convergence rate of neural networks for learned functions of different frequencies.
In \emph{Advances in Neural Information Processing Systems} 32, 2019.

\bibitem{basri2020}
R.~Basri, M.~Galun, A.~Geifman, D.~Jacobs, Y.~Kasten, and S.~Kritchman.
Frequency bias in neural networks for input of non-uniform density.
In \emph{Proceedings of the 37th International Conference on Machine Learning}, 2020.

\bibitem{xu2020}
Z.-Q.~J.~Xu, Y.~Zhang, T.~Luo, Y.~Xiao, and Z.~Ma.
Frequency principle: Fourier analysis sheds light on deep neural networks.
\emph{Communications in Computational Physics}, 28(5):1746--1767, 2020.
% VERIFY: volume, issue and pages.

\bibitem{chen2018}
R.~T.~Q.~Chen, Y.~Rubanova, J.~Bettencourt, and D.~Duvenaud.
Neural ordinary differential equations.
In \emph{Advances in Neural Information Processing Systems} 31, 2018.

\bibitem{krishnapriyan2023}
A.~S.~Krishnapriyan, A.~F.~Queiruga, N.~B.~Erichson, and M.~W.~Mahoney.
Learning continuous models for continuous physics.
\emph{Communications Physics}, 6:319, 2023.
% VERIFY: volume and article number.
% TODO(author): \input{eos_bibitems} returns with the edge-of-stability application.
\end{thebibliography}

\end{document}
